% =====================================================================
%  Harald Høffding, EXPERIENCE AND INTERPRETATION.  Studies in the Philosophy of Religion.
%  English translation of: Oplevelse og Tydning. Religionsfilosofiske Studier.  København: Nyt Nordisk Forlag
%  (Otto Halling & Chr. Nielsen), 1918.
%
%  GENERATED; DO NOT EDIT BY HAND.  Built by ebook-method/oplevelse-og-tydning/tr_parts.py build --write from
%  tr_template.tex and the parts in ebook-method/oplevelse-og-tydning/tr/ (translated 2026-10-14 from
%  texts/hoeffding/oplevelse-og-tydning/transcription.tex, following TRANSLATOR-BRIEF.md and GLOSSARY.md there).
%  Page numbers in the margin [N] are those of the Danish print.  Footnote marks *) **) ***) †) as printed.
%  Terms: Oplevelse = experience; Erfaring = experience, marked [Erfaring] where the two are contrasted;
%  Tydning = interpretation (Hans, 2026-10-14).
%  ODD SPOTS: translating was also a second reading of the Danish.  The translators listed 163 places that looked
%  wrong or unclear (tr/odd-*.tsv); the 128 that touch the transcribed text were read at the image (the
%  transcription's header, (f)).  7 were transcription faults, mended in both texts (e.g. Modstykke p. 33, a stray ~
%  p. 36); the rest are the print's own: kept in the Danish, translated here as the sense requires (names as printed:
%  Eckchardt, Edwad Cooke).  The other 35 were questions of grammar or reference, settled in translating.
% =====================================================================
\documentclass[12pt,a4paper,openany]{book}
\usepackage[utf8]{inputenc}
\usepackage[T1]{fontenc}
\usepackage{libertinus}
\usepackage{textalpha}
\usepackage[english]{babel}
\usepackage{setspace}
\usepackage[margin=1.2in]{geometry}
\usepackage{fancyhdr}
\setlength{\headheight}{28pt}
\addtolength{\topmargin}{-13.5pt}
\usepackage{hyperref}
\hypersetup{pdftitle={Experience and Interpretation (1918)}, pdfauthor={Harald Høffding}, hidelinks}
\pagestyle{fancy}
\fancyhf{}
\fancyhead[LE,RO]{\thepage}
\fancyhead[RE]{\textit{Harald Høffding: Experience and Interpretation}}
\fancyhead[LO]{\textit{\leftmark}}
\fancypagestyle{plain}{\fancyhf{}\fancyhead[LE,RO]{\thepage}}
\setstretch{1.3}
\usepackage{marginnote}
\usepackage{graphicx}
\DeclareUnicodeCharacter{0254}{\reflectbox{c}}  % ɔ: the old Danish id-est mark (a reversed c)
\newcommand{\starmark}[1]{\ifnum#1>10 \ifcase\numexpr#1-10\relax\or\dag\or\dag\dag\or\dag\dag\dag\fi\else\ifcase#1\or*\or**\or***\or****\or*****\fi\fi)}   % \footnote[11], [12]: the print's †) ††)
\renewcommand{\thefootnote}{\starmark{\value{footnote}}}   % the print marks its notes *) **) ***) per page
\newcommand{\opage}[1]{\marginnote{\footnotesize\textit{[#1]}}}
\newcommand{\apage}[1]{\marginnote{\footnotesize\textit{[#1]}}}
\newcommand{\whiteline}{\par\vspace{\baselineskip}\par}
% the printed Indhold: an essay or chapter line (title, page), a section line
\newcommand{\ientry}[2]{\par\smallskip\noindent #1\dotfill\,#2\par}
\newcommand{\isub}[2]{\par\noindent\hspace*{2em}#1\dotfill\,#2\par}
% the headings (no \addcontentsline: the print has its own Indhold, and a footnote in a heading breaks a toc entry)
\pdfstringdefDisableCommands{\renewcommand{\footnote}[2][]{}}
\newcommand{\prefacehead}[1]{\chapter*{#1}\markboth{Fortale}{Fortale}}
\newcommand{\booktitle}[2]{\clearpage\begin{center}{\Large #1}\\[6pt]{\large #2}\end{center}\markboth{Experience and Interpretation}{Experience and Interpretation}\medskip}
\newcommand{\introhead}[1]{\begin{center}{\large #1}\end{center}\medskip}
\newcommand{\chapterhead}[2]{\bigskip\begin{center}#1\\[3pt]{\large #2}\end{center}\medskip}
\newcommand{\conclhead}[1]{\bigskip\begin{center}{\large #1}\end{center}\markboth{#1}{#1}\medskip}
\newcommand{\sectionhead}[1]{\par\medskip\begin{center}#1\end{center}\smallskip}
\newcommand{\essayhead}[1]{\clearpage\begin{center}{\Large #1}\end{center}\markboth{#1}{#1}\medskip}
\newcommand{\essaysource}[1]{\begin{center}{\small #1}\end{center}\medskip}
\newcommand{\subhead}[1]{\par\medskip\begin{center}#1\end{center}\smallskip}
\newcommand{\romanhead}[1]{\par\medskip\begin{center}#1.\end{center}\smallskip}

\begin{document}

\frontmatter
\begin{titlepage}
\centering
\vspace*{2cm}
{\large HARALD HØFFDING}\par
\vspace{0.6cm}
\rule{2cm}{0.4pt}\par
\vspace{0.6cm}
{\Huge EXPERIENCE\\[4pt] AND INTERPRETATION}\par
\vspace{1cm}
{\large STUDIES IN THE PHILOSOPHY\\ OF RELIGION}\par
\vfill
{\large NYT NORDISK FORLAG}\par
{\small (OTTO HALLING \& CHR. NIELSEN)}\par
{\large COPENHAGEN}\par
\vspace{1cm}
{\small 1918}\par
\vspace{1cm}
{\small\textit{Translated from the Danish}}\par
\end{titlepage}

% --- p. 5 ---
\setcounter{footnote}{0}%

\opage{5}\prefacehead{PREFACE}

The essay that has given this book its name was written on the occasion of my excellent
friend and colleague Théodore Flournoy of Geneva having published, in his ``Archives de Psychologie'' in
the year 1915, notes by a lady in which she described her ecstatic
states and the influence they had on her religious life. This description seemed to me not only
to have great psychological interest in and of itself, but also
to invite comparison with similar phenomena,
particularly as regards the relation between what is itself experienced
and the involuntary or fully conscious, perhaps willed
interpretations. I therefore resumed earlier studies
concerning the remarkable Spanish nun Santa Teresa.
As an introduction I have given shorter descriptions of
kindred phenomena from various times.

The title ``Experience and Interpretation'' also fits the other parts of the book reasonably well. With good will
it can also be made to apply to the piece on
Madvig as a Philosopher, since it deals in part with his
position on the religious problem, grounded in his own experience. But I admit that the essay on Aristotle's
conception of friendship can be subsumed under
the common title only if one sees it as describing
how the great Greek interpreted his experiences concerning
a significant human relation of life.

Carlsberg, August 19, 1918.
HARALD HØFFDING

\chapter*{CONTENTS}
\markboth{Contents}{Contents}
\opage{7}\noindent\hfill{\small Page}\par
\ientry{EXPERIENCE AND INTERPRETATION}{9}
\isub{Introduction}{9}
\isub{I. \emph{The Primitive Stage.}}{}
\isub{\hspace*{1em}1. Totemism}{17}
\isub{\hspace*{1em}2. Bacchantism}{22}
\isub{\hspace*{1em}3. Witch Trials and Demonic Possession}{23}
\isub{\hspace*{1em}4. Spiritualism}{25}
\isub{II. \emph{Prophetism and Mysticism.}}{}
\isub{\hspace*{1em}1. The Prophets of Israel}{34}
\isub{\hspace*{1em}2. Primitive Christianity}{37}
\isub{\hspace*{1em}3. Neoplatonism}{42}
\isub{\hspace*{1em}4. Augustine}{46}
\isub{\hspace*{1em}5. Medieval Mysticism}{48}
\isub{III. \emph{Saint Teresa.}}{}
\isub{\hspace*{1em}1. The First Stages}{58}
\isub{\hspace*{1em}2. The Highest Stage}{63}
\isub{\hspace*{1em}3. Teresa and the Doctrine of the Church}{67}
\isub{\hspace*{1em}4. Autobiography and System}{72}
\isub{IV. \emph{A Modern Mystic.}}{}
\isub{\hspace*{1em}1. Psychology and Mysticism}{80}
\isub{\hspace*{1em}2. A Modern Mystic's Ecstasies}{86}
\isub{\hspace*{1em}3. Ecstasy and Religion}{96}
\isub{\hspace*{1em}4. The Principle of Personality and Love of Humankind}{102}
\isub{\hspace*{1em}5. Cécile’s ``Seventh Mansion''}{107}
\isub{Conclusion}{113}
\ientry{SOCIOLOGY AND THE PHILOSOPHY OF RELIGION}{117}
\ientry{ON CONSOLATION}{142}
\ientry{FEATURES OF THE HISTORY OF PRAYER}{163}
\ientry{TWO PASTORS}{189}
\ientry{MADVIG AS A PHILOSOPHER}{210}
\ientry{ARISTOTLE ON FRIENDSHIP}{225}

\mainmatter

% --- p. 9 ---
\setcounter{footnote}{0}%
\opage{9}\booktitle{EXPERIENCE AND INTERPRETATION}{A COMPARATIVE STUDY IN THE PSYCHOLOGY OF RELIGION}

\introhead{INTRODUCTION}

T\textsc{here} is, in the strictest sense of the word, no entirely
simple and immediate experience. We are never at any time
wholly passive. Even the simplest sense-perceptions
have their character determined by their
relation to other sensations that occur simultaneously
or have occurred beforehand, as well as by our whole state of mind.
Strong emotion can thus be an
obstacle to the occurrence of sense-perceptions. Apart from this,
processes of recognition at once interweave themselves with
what is sensed or felt. Both in immediate and
in mediate recognition what is experienced is at once brought
into connection with something we knew beforehand or believed
we knew. And this connection can set in so quickly
that we do not notice the process that is taking place here,
but believe that we are behaving in a wholly receptive way. This
can be the case both when the recognition is valid,
and when it (as in an illusion of memory) is invalid.
That we have in fact put something into it, we discover
only when the recognition proves to be false. Then
we can distinguish between experience and interpretation. But in the first moment
they always appear in involuntary connection.
What appears in our consciousness is
always experience and interpretation at once. If by
experience [Erfaring] is understood experiences that have been brought into a definite
context, we may say that experience is a
% --- p. 10 ---
\setcounter{footnote}{0}%
\opage{10}limiting case of experience [Erfaring]. The element of interpretation is here at its
minimum.

In more specific and composite cases it is
easy to distinguish between the two elements. When an object
appears to one person as small, to another as
large, it is obvious that the difference rests on the
standard that each of them brings with him or applies,
without being aware of it. When a color in one case
appears to me as saturated, in another case not,
it can be discovered that the difference rests on which color
I saw together with it or before it. When one
and the same historical person or historical event
is judged differently, despite equally great knowledge of it,
the difference must spring from the judges' different
personal presuppositions. One is then at once ready with
the proposition that what is seen depends on the eyes that see.

In a little sketch of the earthquake in San Francisco (1906),
which he experienced at close quarters, William James
has given an analysis of what took place in him
during this sudden event. The earthquake lasted only
48 seconds, and by thinking back on
that moment --- which he did as soon as he came to his
senses --- he could trace no actual reflection (little
thought, and no reflection); sensation and emotion
dominated the moment. The emotion was delight
and admiration, and they presuppose a certain recognition
(gleeful recognition). The fact was that anyone who goes
to San Francisco must be aware of the possibility of earthquakes.
James thus at once had clarity about what he
was experiencing, and was intent on getting to experience
it. There was no fear present. One of the first thoughts
was how well the little wooden house he lived in
had held up. When there was a recognition, it did not concern
the earthquake as a process in which layers of earth
% --- p. 11 ---
\setcounter{footnote}{0}%
\opage{11}collapse; it stood there as a Something with a certain personification
and came over him as a unified whole. And
this conception of the event as an individual presented
itself quite involuntarily and, so to speak, in an unavoidable
and irresistible manner. No human action could
point more definitely to a living being as
source and origin. Others too received from the event
such an impression of individual action. But
not everyone stood so intellectually interested in
the experience and was so glad to be part of it,
as the philosopher. For some it stood as a demonic
being. For one lady it stood as the end of the world and
the beginning of the Judgment, and this theological interpretation
occupied her mind so completely that it was only after she
had come down into the street that it dawned on her that it
was an earthquake.\footnote[1]{\emph{William James}: \textit{On some mental effects of the earth-quake}. (Memories and Studies. 1911).}

Not least in our dreams interpretation and experiences
weave together. At the transition from waking to
sleep one can sometimes observe how sense-perceptions
and elementary hallucinations are involuntarily
combined into objects and events. Such hypnagogic
states make it possible to spy on the work of our
mental life in connecting its elements and
bringing intelligible connection among them. And in
this work memories of earlier
experiences naturally co-operate. Also within the dream itself
this work can sometimes be traced, especially in the so-called dramatic
dreams, in which we are in distress or embarrassment
and then find help and consolation. A deceased colleague told
me that he dreamed he was walking naked in the street and was
very unhappy about it; but then he had met
% --- p. 12 ---
\setcounter{footnote}{0}%
\opage{12}me, and I had consoled him with the remark that he was dreaming.
He must himself have devised this consolation, though it seemed
to come to him through someone else. During awakening
(in the hypnopompic state) work often goes on
that brings more connection into the dreams
one has had than there was. It has been established thereby
that when a dream is suddenly interrupted by awakening,
and one then at once notes its content, this presents
a series of scenes without connection.
The longer one postpones the noting, the more connection
the dream presents, and this connection
generally grows still greater through the later accounts.
The more one can, so to speak, take the dream by surprise,
the more chaotic a character it has.\footnote[1]{Cf. \emph{Foncault}: L'évolution du rêve pendant le réveil. (Revue philosophique. 1904).}

Emotions and changes of mood can arise suddenly
and apparently without motive, and then often become
very painful to the one who experiences them, because they do not
seem to connect with earlier experiences. It
then often turns out that it is some painful experience,
perhaps far back in time, and long involuntarily
repressed and shut out from memory, --- a
``trauma'', --- that, so to speak, in a subterranean way
exerts its after-effects through low spirits and depression.
By searching questioning and discussion, ---
so-called psychoanalysis, --- modern psychiatrists
seek to draw the hidden cause forth and to remove its
influence by placing it in the clear daylight
of consciousness.

Hallucinations, whether they appear as visions
or as auditions or move in the domains of the process of touch,
often make a great effect by their
% --- p. 13 ---
\setcounter{footnote}{0}%
\opage{13}sudden appearance, and yet it can in many cases be traced
that they have developed successively, do not stand at a stroke
in their completed form, but grow from the indefinite
and particular to more definite, clearer and more richly equipped
forms. And this growth then proves to be
conditioned by the individuals' world of memory and world of thought.
This explains why visions on the whole
agree with the whole spiritual horizon
within which the visionary individuals live. Besides
the examples to which I have referred in my Psykologi
(6th ed. p. 192), namely Benevenuto's vision in prison
and Bernadette's vision in the rock crevice at Lourdes, I shall
here refer to a vision in the Revelation of John
(I, 10 ff), where the audition precedes the vision, in that
the seer hears a voice behind him, and the sight comes only
when he turns around. In later mystics, whose
description is more exact, the indefinite sensation
of a presence often precedes the audition, and
this again the vision. The Spanish nun
Sta Teresa noticed in an ecstatic state that someone was
present; the unaccustomed state, in which nothing was heard
or seen, instilled fear in her, and she began to
weep. But then she heard calming and consoling words
and thereupon became convinced that it was Jesus who was present.
Although she was plagued by her confessor with the
question how she could know that it was Jesus,
when she saw nothing, she nevertheless held fast to her interpretation.
The sensation of presence persisted for several days, and
thereupon the hands of the figure showed themselves to her, some
days later the whole figure. That Jesus thus little
by little revealed himself to her, and not at once wholly,
she derives from her weakness and her imperfect faith.\footnote[1]{\textit{Vie de Ste Terèse écrite par elle-même}. Trad. Bruix. Chap. 27-28.}
% --- p. 14 ---
\setcounter{footnote}{0}%
\opage{14}--- Examples of what appeared in the Spanish nun's exalted
visions are also to be found in more vulgar
hallucinations.\footnote[1]{See my \textit{Religionsfilosofi} § 57.} And what must especially be noted, the
ideas that are woven into the experiences need not
be objects of recognition, even if it can be demonstrated
that they are due to the individual's earlier experiences [Erfaringer]. A lady
among Morton Prince's patients saw her deceased husband
before her and heard him speak consoling words to her.
It later turned out that these words she had some months
before heard from a friend, to whom the deceased had at one
time spoken them. When they sounded for her in the hallucination,
she did not recognize them, but they stood before
her as new. When she had been hypnotized, she remembered
both the conversation with the friend, and that it was
he who had heard the consoling words from the deceased.\footnote[2]{\emph{Morton Prince}: \textit{The Unconscious}. The fundamentals of human personality normal and abnormal. New York 1914. p. 40.}

From the experiences now discussed the purely ecstatic
states differ in being so concentrated
that nothing can be remembered from them. Every emotion
has a tendency to inhibit or exclude
all sensations and ideas; in ecstasy
this tendency appears at its culmination. In the ecstatic
state hallucinations may arise that perhaps
can be remembered later, just as dreams can be remembered.
But pure ecstasy is declared by those who have experienced
it to be indescribable, although it is at the same time declared to
be the highest bliss. As we shall see in the course of this
investigation, the mystics of a more recent age are able
better than those of earlier ages to distinguish between ecstasy itself
and the visions or images it can bring with
it, and it is precisely a chief aim of the present
% --- p. 15 ---
\setcounter{footnote}{0}%
\opage{15}essay to demonstrate this. --- If one asks about the
physiological side of ecstasy, hardly a
better answer can be given than that contained in the following statements
by Flournoy: ``As regards what takes place in the
nervous system in the few moments in which the crisis is
at its height, one must admit that we know nothing
about it. To judge from the mental effects that follow
it, --- certainty, joy, peace, courage, feeling and exercise
of greater energy and self-surrender --- it must be
a reorganization of the brain centers, the establishment of new
coordinations of a sensory and motor kind, which gives the
individual a new attitude toward the surroundings\dots{}.
One might present the matter by one of the anatomophysiological
schemas of the kind that have been proposed for
similar cases. But all that is still so vague that I
will not dwell on it.''\footnote[1]{Th. \emph{Flournoy}: \textit{Une Mystique moderne}. (Archives de Psychologie. Genève 1915). p. 181.}

By its character of involuntary concentration\footnote[2]{The ecstatic concentration can sometimes be prepared by a concentration of attention and thought prepared with effort and maintained. But this is not always possible, and at the height it has a character of being seized and gripped, which precisely conditions the enigmatic quality of the phenomenon.}
ecstasy stands in sharp contrast to the ordinary
life of consciousness, which is precisely conditioned by contrast
and alternation. Therefore ecstatic states are especially suited
to illuminate the relation between experience and interpretation.
Interpretation has in such states very little
to build on. The way in which ecstatic states
are brought into relation to people's other experiences
and to their life of ideas will vary with the
stages and forms of spiritual development.

Three historical types can be distinguished, as regards
% --- p. 16 ---
\setcounter{footnote}{0}%
\opage{16}this relation, --- an animistic type, a dogmatic
type and an empirical type.

In animism the explanation of ecstatic states is found
in the influence of, or rather the indwelling of, a foreign
mental life. This interpretation is carried out quite involuntarily,
in that the transmitted faith in the spirits of ancestors
(or similar beings) immediately dominates the mind.
Experience and tradition fuse from the very first,
and it is their common product that makes itself known
to consciousness, without its being able to analyze it.

At the dogmatic stage there is a developed worship
and a doctrine of faith carried through with reflection, which are
recognized by the individual as authorities that ought to be taken as
basis in the interpretation. The difference between experience and
interpretation begins to assert itself, but is abolished
again through reverence for the transmitted norm. In
Israelite prophetism and Christian mysticism there are
significant examples of this, and we shall dwell especially
on Saint Teresa, whose mysticism has great psychological
interest.

When reflection has developed so far that the relation
between experience and interpretation can become an object
of critical examination, experience comes to stand with
greater independence than in the two preceding forms,
and it can even assert its superior power against every
attempt at interpretation. Neither fusion
nor dogmatic adaptation takes place, and the immediately
given comes more into its own. In the most recent time
Flournoy in his essay ``Une Mystique Moderne'' (which
has the subtitle ``Documents pour la psychologie religieuse'')
has given an extraordinarily interesting example of
this type.

% --- p. 17 ---
\setcounter{footnote}{0}%
\opage{17}\chapterhead{I.}{THE PRIMITIVE STAGE}

\sectionhead{1. \textit{Totemism.}\footnote[1]{Cf. I. G. \emph{Frazer}: \textit{The Golden Bough}\textsuperscript{3}. (1913). Especially Part VII, 2. p. 225---278. --- Émile \emph{Durkheim}: \textit{Les formes élémentaires de la vie religieuse}. Le système totémique en Australie. (1912). p. 141---199; 223—233.}}

All over the earth, among the most diverse peoples,
there takes place at the transition to adult age
an initiation of the young men, which is to make them
participants in the tribe's tradition and faith. Among the most
primitive peoples of whom we have closer knowledge, the Central Australians,
this initiation stands in close connection with
totemism. A totem is an object, for the most part
a plant or an animal, that stands as expression of the tribe's
essence and as condition of its welfare. Mythologically
this is explained by the tribe's primal fathers in the morning
of the ages having transformed themselves into such an object.
The point now is to lead the young to feel themselves One, or rather
to become One, with the tribe's essence expressed in the totem. This
is achieved through an initiation, after which there is handed over
to the young man a sacred implement, a churinga, a sacred
stick or stone, which is provided with signs and is held to be
the common body of the man and of his totem primal father, so
that the young man is thereby connected at once with
the ancestor and with the tribe's totem, into which the 
% --- p. 18 ---
\setcounter{footnote}{0}%
\opage{18}ancestor had at one time transformed himself.\footnote[1]{See \emph{Strehlow}: \textit{Die Arunda- und Loritza-Stämme in Zentralaustralien}. Frankfurt 1911. II p. 75 ff. Churinga is said to mean ``the own secret''. --- The excellent missionary, to whom some of the best information about the Central Australians is due, testifies to the seriousness with which the initiations were performed. The ceremonies are regarded as a kind of divine service, just as the Christians regard their own religious exercises. ``This,'' he adds (III p. G.), ``was assured to me by both Christian and heathen natives.''} The feeling of unity
with the tribe's totem becomes after the initiation so strong
and intense that every single individual declares himself to
be the animal or the plant that is the totem. A Central Australian
can therefore say: ``I am a kangaroo!''
``I am an emu!''

The difference between essence and symbol does not exist at
this stage. When this difference is emphasized,
doubt has already begun to stir. The great psychological
struggle that extends through the whole history of religion
concerns precisely this difference between essence and
symbol (logically speaking: between identity and analogy). When
Luther at the colloquy on religion at Marburg maintained that
in the proposition: ``This is my body and blood'' the word
``is'' must be taken literally (as copula), whereas Zwingli
interprets ``is'' as ``means''\footnote[2]{See on the discussion between Luther and Zwingli: \emph{Eduard Zeller}: \textit{Das theologische System Zwingli's}. Tübingen 1853. p. 129.} --- then the spirit of the original,
primitive religion speaks undoubtedly through
Luther. He speaks like that totemist who ``was'' a kangaroo.
---

A totem is an expression of the tribe's history. The primal fathers,
in the old days, felt their fate tied to
a definite object, which thereby became the tribe's totem;
these fathers' souls are born, according to the people's faith, anew
in the following generations. But in order truly to become a partaker in
the ancestors' essence and become One with them, the natural
% --- p. 19 ---
\setcounter{footnote}{0}%
\opage{19}birth is not enough. A new birth is required, and the young
attain it through the initiation they undergo
at the transition to adult age. By strict exercises
and castigations, partly mutilation (such as circumcision and
cutting off of teeth), partly privation and self-torments, their
minds are prepared to receive the sacred tradition. They
are brought into a state of passivity, which is akin to
dream and ecstasy, and thereupon they are brought, under very
dramatic forms, into connection with the Holy,
whereupon the handing over of the sacred implements takes
place. Sometimes a real birth is imitated at the initiation, as well
as possible. From the paroxysms that
the initiation and the preparations for it have called forth,
the individual wakes up to new life, and some
time may pass before the memory of what precedes the
whole ordeal awakens again.

Not always is the individual's totem the same as
the tribe's. In the visions that can arise in a young
man under the exaltation aroused by strict exercises and self-torments,
he sees perhaps an animal or a plant
that comes to stand for him as tied to his
essence and fate. But such totemism is derived from
the social totemism, which shows itself in an older
relative guiding the young man during the preparation
that leads him to choose his totem.

In a similar way, only still more violently,
it goes when peculiar medicine men, shamans
or priests are initiated. Here too the point is to bring about
a state in which the mind is opened to the sacred
tradition. The paroxysms are here stronger, assume
what one now calls a hysterical character, and it
is presumably also individuals with a predisposition to hysteria who most readily
become medicine men, and through whom the spirits
of the departed can most easily come to speak.

% --- p. 20 ---
\setcounter{footnote}{0}%
\opage{20}An analogy to the Central Australians' initiation magic, thus
a similar relation between ecstasy and tradition,
occurs among the Eskimos. The angakoks are trained
and initiated, under the guidance of paid teachers, to see
the spirits, get answers from them and get them to help.
The training lasts about 10 years. After a kind of
self-hypnotization (by rubbing for hours a small stone
against a larger) a time is spent in fasting and tension.
Then the young man is devoured by a sea monster, but lives
again, --- is thus reborn. The ecstasy, which sets in
at the culmination of this whole process, is conceived as
the departure of the soul from the body.\footnote[1]{W. \emph{Thalbitzer}: \textit{The Heathen Priests of East Greenland} (10. Amerikanistkongres. Wien 1909).}

At a stage like that appearing in these examples
the difference between experience and interpretation does not
assert itself. After his initiation the young Central Australian
feels himself immediately One with his totem. He does not distinguish
between his state and that which in this state
appears to him. Such a distinction would
presuppose that he could experience similar states without
interpreting them, and that he on the other hand could stand
over against the tribe's tradition as a particular fact.
But neither of the two is the case. The more inclined
to ecstasy he is, the more quickly the transmitted
ideas present themselves, and the more he is convinced
that his own soul has been away, and that his
ancestor's soul has taken its place. Ecstatic states
are regarded not merely as particularly susceptible to
the intervention of spirits, but such states stand as
inexplicable if one does not assume such an intervention.
Ecstasy itself is a sign of the working of a spirit.\footnote[2]{I. G. \emph{Frazer}: \textit{The Belief in Immortality and the Worship of the Dead}. 1913. p. 15. 444.} In
% --- p. 21 ---
\setcounter{footnote}{0}%
\opage{21}Greenland older angakoks select as their pupils
especially boys who have shown themselves to be clairvoyant.\footnote[1]{W. \emph{Thalbitzer}: \textit{The Heathen Priests of East Greenland}. p. 454.}

There is indeed something mysterious both about the ecstatic
states and about the transmitted ideas
that are used to interpret them. Neither ecstasy nor
tradition seems able to be traced back to definite
causes. One does not know, and one does not ask, whence
they come. They both work through their mysteriousness,
both point beyond the individual's limited experience
and thereby exert a mutual attraction on each other,
imperceptible to the individual persons. Through the
solemn initiation the young undergo, the
two elements are united once for all for them.

It is fully consistent when the primitive
human being regards his ecstatic experiences partly as
effect of spirits, partly as expression of a communion with
spirits. The world of spirits is for him a reality, a
part of his world. Doubt perhaps stirs in individuals;
even among the Papuans there are said to be skeptics
who do not share the tribe's faith. But on the whole
tradition stands as sacred, i.e. inviolable, indisputable.
It is then fully consistent when what is related about
the world of spirits is involuntarily applied to the explanation of
otherwise inexplicable phenomena. ---

What is thus characteristic of the stage that we
call the primitive (without thereby presupposing that there
should not have been other stages preceding)\footnote[2]{In my criticism of the book by Durkheim mentioned above (see below the section ``Sociology and the Philosophy of Religion'') I have laid weight on the fact that the religion that is regarded as the most primitive we know, the Central Australian, we know only as fully finished, so to speak systematized. It must have reached this fixed form through a long series of individual initiatives and psychical processes (combinations, postponements and displacements).} asserts itself
% --- p. 22 ---
\setcounter{footnote}{0}%
\opage{22}later on under other forms, right up to our
own time. In what follows immediately I give some examples
of this.

\sectionhead{2. \textit{Bacchantism.}}

From the Orient strong currents of ecstatic
enthusiasm have at various times passed into Europe.
Thus the Dionysus cult came from Asia Minor through
Thrace into the Greek world at a very early
stage of its culture. True, the oriental
Dionysus underwent a thorough Hellenization, and the ecstasies
were brought, especially by the Delphic oracle's wise conduct,
into harmonious forms. But the original
ecstatic element shows itself constantly, even after Apollo
had entered into a pact with Dionysus. At Delphi the
ecstatic priestess's body is animated by the god, or her
soul is released from her body, in order to receive divine
revelations; one seems not sharply to
have distinguished between these two explanations, since the human being
by both ways can be said to become One with the god,
raised above his customary consciousness. The originally
Greek word enthusiasm, which has passed into more recent languages
to denote every kind of enthusiasm, properly means the
state of being in God. In the period from the eighth to the
sixth century B.C. there arose in Greece a whole
series of seers and prophets, in legend under the name of Sibyls
and Bacids. In them a need and longing found vent
that was not satisfied by the clearly formed Greek
worship. They were individual personalities who by their
enthusiasm and ecstatic nature won authority in larger or
smaller circles and in the following generations had to
give name to various written books of prophecy.\footnote[1]{Erwin \emph{Rhode}: \textit{Psyche}\textsuperscript{2}. II. p 59—65.}

% --- p. 23 ---
\setcounter{footnote}{0}%
\opage{23}The oriental spiritual tendency that expressed itself in the Dionysus cult
spreads at the close of antiquity in a
series of forms.\footnote[1]{\emph{Sam Wide}: \textit{De græsk-orientalske Mysteriereligioner} (1915).}

\sectionhead{3. \textit{Witch Trials and Demonic Possessions}\footnote[2]{\emph{Bror Gadelius}: \textit{Tro och öfvertro i gångna tider}. Stockholm 1912—13.}}

The medieval faith in demonic possessions and the witch trials
connected with them rests, as the
Swedish psychiatrist Bror Gadelius has shown in an interesting
work, on the same psychological presuppositions
as the primitive human being's religious beliefs. In
earlier times, he maintained, the tendency to ecstasy
and hallucination was far stronger than nowadays, and the
dark sides of mental life played a far greater role
than now in relation to reflection and self-possession. What
would now stand as abnormal could then very
well accord with normal mental life. We dare not transfer
the modern distinction between healthy and sick mental life
to the older times. Even to this day this
difference is not easy to make. The tendency to splitting of personality
was formerly far more common than now, and
sudden changes of state, visions and impulses that
stood in opposition to the individual's customary conduct
were regarded as manifestations of the Devil's activity; if
they got the upper hand, both the individual himself and others
regarded the state as a demonic possession.

Of Gadelius' many examples I cite the following,
which is taken from the Swedish nobleman Niels Gyllenstierna's
unprinted autobiography and deals with his
spiritual affliction during a visit to Stockholm, when
he was 22 years old (1732). He was transported in the ecstasy into
% --- p. 24 ---
\setcounter{footnote}{0}%
\opage{24}``the dark eternity''; ``the eternal darkness that was in his
soul and outside it in the whole black abyss'', hid
almost the customary world from him. As he then
describes what he in this state sometimes ``saw,
heard, understood, felt and experienced'', he mentions a multitude of
dark spirits, who were the authors of all vices
and misdeeds on earth. He saw the red dragon
from the Revelation of John etc. --- Gadelius remarks,
that what here showed itself to Gyllenstierna stood in immediate
connection with the religious movements that
pervaded the age of Pietism. ``He is a visionary
hysteric of the folk-psychological type that is
so frequent in the history of religion, and one must judge
every such case by a special standard, which
is adapted to the circumstances of the time and other presuppositions.
Especially it is worth noticing the
ease with which inner conviction is embodied
by sense illusions.''

Here again, then, there is a meeting of inner states
and transmitted ideas, without the individual being able
to conceive the possibility of other explanations of his
immediate experience, which here consists in the darkening of the mind.
---

A remarkable counterpart to Gyllenstierna's ``Experiences''
is formed by Swedenborg's visions, which were rather
of a bright character. ``Swedenborg too is'', says Gadelius
``an example of the hysterical development of subconscious
(visionary) faculties, and it is undoubtedly not correct
to regard him as insane in the modern clinical
sense''.\footnote[1]{According to (\emph{Kleen}: \textit{Swedenborg}. En lefnadsskildring. Stockholm 1917. I. p. XIX.) Gadelius is said later to have recognized that \textit{Swedenborg} suffered not only from hysteria, but also from paranoia (derangement). The characteristic relation between the life of moods and the nature of the visions retains its interest, and Gadelius' significant warning against applying without further ado the modern psychiatry's distinction between healthy and sick mental life to earlier times' ``folk-psychological types'' still stands firm.} It was brighter thoughts, partly arisen
% --- p. 25 ---
\setcounter{footnote}{0}%
\opage{25}under the dawning rationalism\footnote[1]{As is clearly seen from the Swedenborgian vision I have cited in my \textit{Religionsfilosofi}. § 57.}, that here determined Swedenborg's
visions. Even where he depicts Hell, this depiction is
abstract and colorless in comparison
with what the folk belief of the eighteenth century demanded.

Such a contrast shows how the interpretation
of what is experienced, even when it comes involuntarily and in
the form of vision, is conditioned by the character of the one experiencing,
his past and his imagination, and also how
interpretation and experience can immediately merge into one.\footnote[2]{Only after the composition of this chapter did I come to read \emph{Henrik Schück}'s fine portrayal of Petrus de Dacia, author of the legend of Saint Christina \textit{(Vår første Forfattere}. En själshistoria från medeltiden. Stockholm 1917). It gives classical examples of how ecstatic and hallucinatory experiences are interpreted with great inwardness, but in conformity with the ideas of the time. In Christina it is partly diabolical, partly divine forces that stand behind her extraordinary experiences. ``All her impulses were embodied and assumed shapes of angels and devils, who waged a murderous struggle for her soul. The dreams that for a modern person appear in the form of poetry, appear for her in the form of reality.'' (p. 101). --- It would be of special interest to compare her ecstatic states with the modern mystic's, of whom we shall later speak.}

\sectionhead{4. \textit{Spiritualism.}}

Of Swedenborg the man who in the nineteenth
century strove to reach a scientific basis for
belief in spirits,\footnote[11]{\emph{Frederic W. H. Myers}: \textit{Human Personality and its survival of bodily death}. (1904). I. p. 6.} says that he led philosophy back to
% --- p. 26 ---
\setcounter{footnote}{0}%
\opage{26}Heaven, in contrast to Socrates, who in his time had
led it from Heaven to earth. It was Frederic
Myers' life task to find a scientific basis for
belief in the eternity of human souls, which neither
the hitherto existing religions nor the hitherto attempted
speculations had been able to do. And the task was to be
solved through a study of the ecstatic states, a
study that was to be carried through according to the strict method
of empirical science. Since he believes himself to have arrived
at the result that the manifestations that occur
in such states cannot be explained by psychological
and physiological means, but necessarily presuppose
the continued existence of departed human beings and the possibility
of their intervention, he finds what he seeks in spiritualism,
with full consciousness that he thereby returns
to an old folk belief.

``I admit most readily'', he says (II. p. 218),
``that my true predecessors in many of the thoughts
that are maintained in this book are the poor thinkers of the Stone Age,
the believing adherents of witchcraft and shamanism.''
And he finds here a great continuity through
the ages. ``Of all subjective religious experiences ecstasy is
the one that most insistently and convincingly
asserts itself, and it is not limited to any
single religion\dots{} From the medicine men of the lowest savages
up to the Apostles John, Peter and Paul,
Buddha and Muhammad not to be forgotten, we find
accounts which in their psychological essence are alike,
however different they may be in moral and intellectual
respect. We find also at all stages that the soul
is believed to leave the body, or, if not leave it,
yet extend its range of experience (range of perception)
in a state akin to ecstasy\dots{}. The medicine man
has (when he is not a mere impostor) just as much
% --- p. 27 ---
\setcounter{footnote}{0}%
\opage{27}access to a spiritual world as Peter or Paul;
--- only the part of the world of spirits he gains access
to is different from theirs: instead of feeling himself raised up
into a brighter world, he is struck with horror by his misty
and dark visions. But for us the mere fact
that we believe in his vision helps to strengthen and support
our faith in the Apostle's vision of `the seventh heaven'.''
(II. p. 260).

The mediumistic phenomena indicate, according to Myers,
that human beings can lose their ordinary consciousness
(the supraliminal soul), and that the subconscious
(the subliminal soul) can step into its place. In the subconscious
the souls of the departed can take up their abode and
say or write what the human being in the supraliminal state
could not know or think. Or else the
subliminal soul itself can go wandering, move about in the regions
of the world of spirits and bring knowledge back.
In the first case a taking possession (possession)
takes place; in the second case there is ecstasy in the proper
sense (II. p. 259).\footnote[1]{Ecstasy means literally ``to be outside oneself''.}

When Myers, the theorist of ``scientific'' spiritualism,
thus resolutely places himself on a line with ``the Stone Age's
poor thinkers'', there is yet the significant
difference that whereas the primitive human beings' belief in spirits
rests on immediate and involuntary fusion
of experience and interpretation, of individual need and social
tradition, Myers builds his faith on what he regards
as an empirical-scientific investigation. It was his
personal need for light and consolation that led him to
try whether it might not be possible by the way of observation to
attain a knowledge of the world of spirits that neither positive
religion nor speculative thinking had been able
% --- p. 28 ---
\setcounter{footnote}{0}%
\opage{28}to bring. He also pays homage in his own way to the most important
principle of empirical science, namely that the causes
by which given phenomena are explained must themselves
be capable of being shown as phenomena within our experience [Erfaring].
Thereby he comes into opposition to the superstition of earlier times.
Only spirits whose existence has been experienced [erfaret]
independently of the experiences that are to be explained can be regarded
as causes of these. Therefore in spiritualism
there is no longer talk of angels or devils\footnote[1]{A devil is not a creature whose existence is independently known to science. (II. p. 199).}, only of
departed human beings' souls. In any case there can
be no talk of the influence of devils; for in none of the very
many observations Myers believes he has at his disposal has he
met with any connection of spirit and evil.
The old belief in the Devil melts away, the more one studies
the body of evidence that is available (the evidence before
us) (II. p. 78-200).

A rationalistic tendency shows itself here in Myers,
which exerts a critically purifying effect. It separates
him from such spiritualists for whom spiritualism is from
the first a religion, and who therefore expressly
warn against methodical inquiry in these matters.\footnote[2]{See on \emph{Stainton Moses}: \textit{Human Personality} II, 223---237.} Myers
wants to unite faith and knowledge. He and his followers have
also (like so many mystics) had to contend
with the belief that all revelations of spirits, especially such
as occur in ecstatic states, are the work of the Devil.\footnote[3]{The pastor \emph{Sanders} was so unfortunate as to have unauthorized revelations. (II, p. 17).}
Every positive religion must limit the extent and nature
of the revelations that are to be recognized
as such. And the new access to the world of spirits,
which according to Myers and his followers is now
% --- p. 29 ---
\setcounter{footnote}{0}%
\opage{29}opened, after it since the primitive human being's time
had so often been misjudged, also makes claim
to take the place of dogma and church faith, whose impotence
is held to be firmly established. ``I maintain'', says Myers
(II, 279), ``that a beginning has now been made of the application
of a method for winning divine knowledge with the
same certainty and calm assurance as that with which
we constantly make progress in the knowledge of earthly
things. The authority of observation and experiment will take
the place of the authority of dogmas and churches'' (II. p. 279).
Christ stands here as a great predecessor. But it has
gone with his discovery of the world of spirits as with the
Norsemen's discovery of America. It belongs to legend
and cannot now show us the way. The land must be discovered
anew, not by any individual, but through a great
human joint labor of interpreting observations that one
has hitherto misjudged.\footnote[1]{\textit{Fragments of Inner Life}. p. 86. (In ``Fragments of Prose and Poetry by F. W. H. Myers. Edited by his wife Eveleen Myers''. 1904).}

Faced with the objection often raised against spiritualism,
that the things the ``spirits'' communicate are as a rule
trivial and insignificant, if not foolish, Myers declares
(II. p. 202) that one must not meet with too great
expectations regarding what the spirits of the departed can
communicate. They can naturally only tell us what they know,
and we cannot expect that they have quickly been able to gather
knowledge of the new conditions in which they live. And at the same time
their communications are always conditioned by the medium's capacity of apprehension
and cannot rise high above the latter's circle of ideas.

Whereas totemism and the belief in witches are immediately dominated
by the folk belief, modern spiritualism
comes to stand in opposition to two independent, fully
% --- p. 30 ---
\setcounter{footnote}{0}%
\opage{30}organized spiritual authorities. The one is the Church's
dogmatics, the other is modern science.
Against the first spiritualism turns in principle, yet
in such a way that it seeks to show that the great personalities
to whose work the rise and continuance of the Church is due drew
their thoughts from ecstatic communion with the world of spirits.
Toward science, on the other hand, spiritualism itself
believes it does not stand in any opposition; it believes
precisely to proceed according to scientific method. As we
have seen, Myers believes he builds on the principle that the cause
of a phenomenon is to be sought in something that itself is present
within experience. The appeal to the spirits of the departed
satisfies according to his opinion this principle, for it is indeed
a fact that the departed have really lived! But
--- the question is precisely whether they still live, and about that
there is no experience. And the question is next whether
they are what they in the visions give themselves out to be, --- whether
their identity can be demonstrated. And for both spiritualism
lacks every means.

As regards the last question, the recognition of the spirits,
it already came up in the ancient
Church. While in the orthodox Church about the year 200
A.D. they had begun to settle down with the idea that revelation
had once for all taken place and been laid down in the holy books,
the Montanists maintained that new revelations continually came
through ecstatic experiences. They were then asked
how they could recognize the spirits, be assured
that it was the true God or the real
Jesus that showed himself to them. To this Tertullian,
the most outstanding representative of Montanism,
answered with the counter-question that undeniably struck
his Bible-believing opponents: ``How was
Peter on the Mount of Transfiguration able to recognize Moses
and Elijah? It must have been in the spirit (in spiritu).
% --- p. 31 ---
\setcounter{footnote}{0}%
\opage{31}From images or statues he cannot have recognized
them, for such were forbidden in the Law. He must
have been in a supersensible state (in spiritu, non in
sensu constitutus), since he recognized them!'' Thus
the recognition of what shows itself in the ecstatic
states is for Tertullian so far from being a problem,
that he sees precisely in the fact that Peter recognized Moses and Elijah
a proof that he was in ecstasy
(amentia), overshadowed by divine power (obumbratus
virtute divina).\footnote[1]{See F. C. A. \emph{Schwegler}: \textit{Der Montanismus}. Tübingen 1841, p. 25.}

But the spiritualists cannot dismiss the problem in
this way, if they wish to stand on the ground of empirical
science. In empirical science it is not
least important to be able to identify phenomena. As we
shall see, the problem of recognition turns up again in the Catholic
and in modern mysticism.

The main point, however, is this, that before one resorts
to supernatural (or at any rate supernormal) explanations,
one must demonstrate the impossibility that the content of the visions
could be due to earlier experiences which the individuals
in their ordinary states cannot call up in memory.
The study of hypnotic states and of double consciousness
has in one case after another established
that ideas that an individual in his ordinary state
does not have at his disposal can come forth in hypnosis
or in trance. On the whole, by close knowledge
of a medium's earlier history one will very often find
the key to ideas and tendencies that at the moment
stand as mysterious, because they do not seem compatible
with the existing states and situations.
Even imagination often produces, especially in its
highest forms and degrees, images and thoughts that stand
% --- p. 32 ---
\setcounter{footnote}{0}%
\opage{32}as inexplicable to the individual himself and to others. An
unconscious work has gone on, which may very well have
been able to make use of what the individual has come to know in his
earlier life. The unconscious in mental life (whatever it may
otherwise be) can thus not only preserve experiences and ideas
that the individual in the ordinary state does not acknowledge
as his own, but also produce new thoughts and images, without
there being reason to assume other forces at work
than such as history, psychology and physiology can
demonstrate. So long as the content of a vision can be explained in an
``endogenous'' way, that is to say, from the individual's earlier
history, there is no reason to appeal to the departed.
Myers himself admits indeed that the content of the vision
is conditioned by the individual's capacity of apprehension and presuppositions.
But can any limit be set to this conditionedness?
Where the endogenous explanation cannot be
completed, we stand before a problem, and the proof that
this problem can be solved only by the appearance of the departed,
rests upon spiritualism, not upon its opponents.
From the fact that the individual during the vision believes himself to be wholly
passive, nothing can be concluded. For there can be an activity
in us that does not become an object of our consciousness, but
only in its results comes up over the threshold of consciousness.

It is the Geneva psychologist Théodore Flournoy who
with the greatest energy and with the most open mind has studied
the so-called spiritualistic phenomena. He has especially
laid weight on knowledge of the mediums' whole past
and has applied this principle with great success in his masterly
study of the somnambulist Mlle Smith.\footnote[1]{ \textit{Des Indes à la Planète Mars}. Essai sur un cas de sonnambulisme avec Glossolalie. Genève 1900.} Later
he has, besides various new ``Documents d'Enquête''
% --- p. 33 ---
\setcounter{footnote}{0}%
\opage{33}given an excellent development of the psychological and epistemological
principles that are of importance with regard to
the phenomena to which spiritualism appeals, and
with regard to the interpretation it gives of them.\footnote[1]{\textit{Esprits et Médiums}. Genève 1911.} In the most recent
time he has brought forth ``a modern mystic'', who in
a later section will be examined as a characteristic
counterpart to the mysticism of earlier times.

With respect to the spiritualistic interpretation Flournoy stands decidedly
doubting. There is not a single fact
that has convinced him, though he admits that several
remain unexplained. He has here come to the same
negative result as Henry Sidgwick and William James,
who after some twenty-odd years' occupation with the
phenomena mentioned declared that they had not come
further than they were at the start.\footnote[2]{\emph{William James}: \textit{Memories and Studies}, p. 174.}

% --- p. 34 ---
\setcounter{footnote}{0}%
\opage{34}\chapterhead{II.}{PROPHETISM AND MYSTICISM}

\sectionhead{1. \textit{The Prophets of Israel.}}

The prophets of the Jewish people are remarkable in the history of religion,
not only for the power of their visions, their ethical force
and their lofty poetry, but --- and this is of special interest
for our subject --- because the relation between experience
and interpretation appears in them with greater clarity
than where belief in spirits is at once woven into the experiences.
Within the Old Testament there appears
a sharp contrast between demonic and properly
prophetic inspiration.\footnote[1]{Gustav \emph{Hølscher}: \textit{Die Profeten}. Untersuchungen zur Religionsgeschichte Israëls. Leipzig 1914, p. 94, 100, 197.} In demonic inspiration
it is an obscure power that seizes the individual.  In an unclear
way the divine being and the personality of the inspired
person melt together into one.  During the ecstasy
unintelligible sounds were stammered out, and by exalted
actions, dance and music, a captivating effect was produced
on the multitude.  Prophecy proper, as it first appears
in the written tradition with Amos, interprets the ecstasy
as the revelation of a God whose will the prophet
proclaims.  The inner experiences themselves were for the prophet
testimony to the dominion of a power different from himself
in his mind; but an obscure fusion of the two does not take
\opage{35}place.  When the prophets brought Yahweh's word to
the people, and when they appeared with authority, it was
because something had happened in their minds that was not ``of
themselves,'' and which they had perhaps even resisted.
He who speaks out of himself they declare
a false prophet.  Often the prophet has resisted
what he was driven to proclaim; in this he sees
(see Jeremiah I, 6-11) a testimony that he is working in
the service of a Higher One.  We often find again this trait of
resistance, before full self-surrender to the ecstatic
state can take place, in the history of mysticism, and often
precisely used as proof of a supernatural intervention.
Prophetism differs from the mysticism of a later age
in that for the prophets it is not first and foremost
peace and rest in the mind that are to be attained through the ecstatic
states, but these states spur them on to vigorous,
often violent action outwardly.  The prophet stands at
once in the midst of his people and raised above it.  Its conduct
and its fate condition his action.  He directs his
speech against priest and king alike.  Prophetism is
a countermovement against the stiffened worship, especially against
the elements of heathen worship that, after the establishment
of the kingdom, were mingled with Mosaism.\footnote[1]{Hermann \emph{Gunkel}: \textit{Israëls Profeter}. (Norwegian trans.) Kristiania 1916 p. 33---37.} Amos
thus speaks against sacrifices, pilgrimages, songs and feasts, and
demands in their place right and justice (V, 21—24).
In a similar way Jeremiah (Chap. 7) and Isaiah (Chap. 8).
The prophets maintain that worship did not belong to the people's
original religion and was not found in the time when
Yahweh led it through the desert.  Against what is established
the prophets maintain the original, just as against
the outward they maintain the inward.

% --- p. 36 ---
\setcounter{footnote}{0}%
\opage{36}The same blindness that made the people fall
into magical worship lulled it into security and groundless
self-confidence in the face of the dangers that threatened from the great world empires,
which step by step developed at the expense of the
small neighboring countries.  The great prophets stood
here as ``prophets of misfortune,'' again and again overwhelmed
by the dangers that threatened land and people, and warning in vain
the short-sighted people and its princes.  Precisely
under such conditions they demanded of each individual
purity and honesty of mind.  Just as the individual
prophet stepped forward in the face of the established national
authorities in church and state and thereby asserted the authority of the
individual personality to follow what his
own experiences had led him to acknowledge as
an ideal, so it was also a personal ideal that they
demanded to be fulfilled by individuals.  If there \textasciitilde{} --- as
Isaiah foretells --- through the suffering that the people will have
to undergo, a possibility of salvation can be glimpsed, then
the condition for each individual will be that in faith
and trust he bow his mind beneath the demands of purity and
righteousness.  Both in the prophet's own mind and in the
demand he makes of his people, there manifests itself a liberation
of the personality from the dominion of tradition and of outward authorities.
And it is the free relation between
experience and interpretation that makes this process of liberation
possible, whereby the prophetism of Israel marks a
great turning point in the history of the human life of the mind.
As prophetism ebbed away, worship
and temple service came again into the foreground, and the people's spiritual
leaders became the priests, not the prophets, even though the
prophetic forms still persisted for a long time.

In Amos, Isaiah and Jeremiah alike there appears in
the description of the visions a progressive clarity.  Often
presentiment and vision come overwhelmingly, without any
% --- p. 37 ---
\setcounter{footnote}{0}%
\opage{37}idea of their significance being formed.  It is
then a divine voice (audition) that helps
to understand the vision.  And on closer reflection
the significance and application of the vision can
be worked out in detail.  It follows of itself that when the experiences
were to be written down, reflection had to take the place of the immediate
being seized.  What thus holds for the
individual prophet in relation to his
experiences also appears in the history of prophecy, which
shows a progressive development from ecstasy to reflection
and to poetry.  ``The ecstatic becomes an author, the prophet
a poet''(Hølscher p. 317).  Ecstasies and hallucinations
were, after all, not in themselves testimony to truth;
they occurred in ``false prophets'' as well as in the
``right'' prophets.  The weight came then, especially from Jeremiah's time
(Hølscher p. 291 f), more and more to lie on the confirmation by experience [Erfaring]
and reflection of what had been
proclaimed after the ecstatic states.\footnote[1]{When, besides the sources themselves, use has here been made especially of the works of \emph{Hølscher} and \emph{Gunkel} on the history of religion for purposes of the psychology of religion, this has been possible in part only because scholars such as those named are themselves well versed in the psychological research of our time.}

\sectionhead{2. \textit{Primitive Christianity.}}

When the relation between experience and interpretation is to
be discussed with regard to Christianity, especially with
regard to ecstatic experiences, then, since the dispute is precisely
in this area still unsettled, it will be important
to be clear about the presuppositions that are laid
down as a basis.  They have already been stated in the foregoing, especially on
the occasion of modern spiritualism.  The essential
presupposition for all psychological and historical research
% --- p. 38 ---
\setcounter{footnote}{0}%
\opage{38}is that the interpretation which human beings give of their
experiences, not least of the most powerfully gripping ones, will always
be conditioned by the experience [Erfaring] they have had before
the experiences, and by the ideas that in various
ways have come to make up their spiritual content.
This is what has been called the endogenous explanation.
The experience and its interpretation then both stand as members in
the whole life of the human being and are understood within the whole
coherence of that life.  As regards Christianity, emphasis must here
be placed on the ideas of a coming Redeemer, who through suffering and humiliation
was to bring salvation and exaltation, ideas that
from the time of Deutero-Isaiah had developed in Jewish
prophetism and later entered into connection with Oriental
ideas of resurrection and the end of the world,
which stem especially from Persian influence.  In the light
of such ideas the individual who through
unusual inner experiences felt himself called to a great
task would naturally regard himself, and in the light of them
those who joined him would naturally come
to regard him.  As to the application in detail there
may be dispute, but an endogenous explanation must keep
to this point of view.  It does not exclude great and
perhaps unique originality; for this expresses itself
in the way in which the various historically given elements
are combined --- partly with one another, partly with one's own
experiences.  And it rests especially on the ethical and humane
ideas that are drawn forth into the light by the interpretation reached
and are maintained as decisive for life.

Confronting the New Testament and primitive Christianity there
are for research two main problems, both of which rest
on the relation between experience and interpretation.

First: how did Jesus of Nazareth himself
interpret what stirred in him and especially came to
% --- p. 39 ---
\setcounter{footnote}{0}%
\opage{39}breakthrough in him in ecstatic states?  Did he
call himself the announced Messiah, or the Son of Man,
or God made man?  He felt himself in any case
to be the one whose coming meant the near
end of the earthly existence of human life, so that
faith in him and his significance had to become a question
of life or death.  His life, as it is set out in
the legends, is a life in expectation, with a view to an
imminent glorification through suffering, a glorification
that was one with the entry of the life of the world into
new, transfigured forms.  Therefore the New Testament
ethics too is an ethics of expectation, an ethics of tension
and hope, which leaves all purely human relations
and tasks as matters of indifference, if not as hindrances;
these relations and tasks belong, after all, to the old
order of things, not to the coming kingdom.  In this
thought, which was historically prepared both in the history
of Judaism and in other Oriental religions, lay a
basis for the interpretation of the ecstatic states, which
again and again took place in Jesus himself and those
nearest to him.\footnote[1]{On the various attempts that have been made, on a psychological-historical basis, to elucidate the life and personality of Jesus, see \emph{Albert Schweitzer}: \textit{Von Reimarus zu Wrede}. Eine Geschichte der Leben-Jesu-Forschung. Tübingen 1906.}

The second question is how the first congregation
conceived Jesus of Nazareth, --- especially how
it interpreted the emotions, the ecstatic
experiences, which his life and death called forth in
the members themselves.  So much is certain, that in the New Testament
authors all gripping mental experiences [Erfaring],
all extraordinary experiences are referred to Jesus of Nazareth as
the source of life, and that all hope and all fear are in the end gathered
in the thought of him and of his return. --- Since now the
% --- p. 40 ---
\setcounter{footnote}{0}%
\opage{40}New Testament authors themselves belonged to the first congregation
and shared its presuppositions, the answer to the first
question is in the end dependent on the answer
to the second, and on this rests the great and difficult
problem that is here posed for historical research.  But
in the end the problem is and remains a psychological one and
concerns the relation between experience and interpretation.  And this problem
becomes all the more difficult in that no firm
distinction is yet made between experience and interpretation,
but they melt together very quickly.  ``The enthusiastic
disciples of Jesus said of him the highest thing that Judaism
knew how to say.  And not in order to fathom
the secret of his person, as if Jesus were
the first thing given and Christology the second, but the souls
who longed for the nearness of the Godhead and felt a need
for a Son of God come from heaven transferred the ideals of their
hearts to him.  The New Testament
Christology is thus a mighty hymn that history sings
about Jesus.''\footnote[1]{\emph{Hermann Gunkel}: \textit{Zum religionsgeschichtlichen Verständnis des neuen Testaments} (1903), p. 94.}

The creed of the Christian church is, from the psychological
standpoint, a more or less systematically carried through
interpretation of the accounts of the experiences of Jesus and the first
congregation, but in such a way that the experiences
and endeavors of later times seek to find their interpretation
in and with the interpretation of the congregation's oldest traditions.
Through a series of combinations, often mere
compromises, these attempts at interpretation advance through the
history of dogma.  A main problem, which confronts all
theology with great difficulties, is how
the wholly new culture, intellectual, emotional and material,
which has developed since the time of the first congregation,
% --- p. 41 ---
\setcounter{footnote}{0}%
\opage{41}can be brought into connection with the New Testament
circle of thought, for which all this cultural work partly did not
exist and partly could have no value.

But as far as the relation between experience and interpretation
is concerned, there arises a decisive difference between the
founding of the church and the continuation of church life
after ecclesiastical tradition and dogmatic
interpretation had become fixed.  With the great prophets, with
Jesus and still in part in the first congregation, life
could still intervene in a decisive way and set
new beginnings, so that the choice of interpretation, of models
and symbols, was conditioned by the experiences.  But
for those who acknowledge the church and the authority of the norms maintained
by the church, the experience itself is determined
by the normed interpretations.  No weight may be laid
on experiences that do not fall within the framework of the
interpretations given once for all.  One is directed
to undergo certain definite experiences [Erfaringer], to have certain experiences,
preferably in a definite order (ordo salutis).  Thereby
religious experience [Erfaring] in the ecclesiastical sense differs from the
kind of experience that otherwise in life and in science gains influence
on people's thoughts\footnote[1]{See further on this \textit{Religionsfilosofi}, § 29 and 59---60.}.

There was in the ancient church a long struggle between
new experiences and handed-down interpretations.  The ecstatic
states, visions and revelations did not cease
because certain books were acknowledged as norms.  Montanism,
mentioned above, was one of the strongest
contributions to this struggle, which for a time indeed subsided
with the consolidation of the Catholic church, but which ---
as will be mentioned in what follows --- was continued in
the Middle Ages in the opposition between mysticism and scholasticism.
It was now not only within the church that
% --- p. 42 ---
\setcounter{footnote}{0}%
\opage{42}ecstatic states appeared, also after the time of the Apostles.
Opponents of Christianity and heretics, too, could
appeal to states that were interpreted as effects of the
intervention of spirits.  A gripping and moving example
is the so-called Fourth Book of Ezra, which was written by
a pious Jew whose visions are determined by horror at the fact
that Jerusalem was destroyed before the Messiah came.  When thus
spirits spoke on both sides, there had to be
a need for a criterion to decide on whose side right
and truth lay.  Ecstasy is in and of itself essentially
of the same nature wherever it occurs, and
therefore contains in itself no guidance for deciding
which interpretation is the right one.  What mattered was to ``try
the spirits,'' and the criterion that was applied in this trial
had in the end to be the ethical effects that the
ecstatic states produced.  One concluded then, as Weinel,
who has illuminated this point very clearly, remarks,
not from what is mysterious or marvelous that
it is a holy spirit that is at work; but one concluded from
the ethical content, from the action that was called forth,
that the states were of a supernatural character\footnote[1]{\emph{Heinrich Weinel}: \textit{Die Wirkungen des Geistes und der Geister im nachapostolischen Zeitalter} (1899), p. 149.}.  It
is then not the intensity but the quality that becomes decisive.
And the intensity also gradually dwindled,
and with it the living expectation that sustained the
primitive Christian ethics.

\sectionhead{3. \textit{Neoplatonism.}}

In the history of mysticism it became of great importance that
the last great current of thought that developed on
Greek soil, Neoplatonism, laid great stress on ecstatic
experiences.

% --- p. 43 ---
\setcounter{footnote}{0}%
\opage{43}Already in Plato there is a mystical element, in
peculiar union with the intellectual and ethical
movement that stemmed from Socrates.  It seems
to have been Orphic influence that thus ran alongside
the Socratic.  Plato often refers to the mystery cults,
their expressions and their ceremonies.  Even though
he decidedly interpreted these in a symbolic way, --- so
that e.g. ``purification'' (catharsis) in him does not mean
liberation from liturgical impurity, but from the thought's and
the mind's enthrallment in base sensuality, --- it is
nevertheless characteristic of him that he feels a need for the circle of
images of the mysteries precisely when his thought reaches its
heights.  When thought, in its ascent from the
finite and sensible phenomena, comes to the concluding
concepts, the last presuppositions, and seeks to
gather them in one Idea, dialectic is replaced by enthusiasm,
and in a sudden vision the Divine is revealed
to it (Phaedrus 249, Symposium 210 E).  And then the
highest is attained; only then is life worth living.  But Plato
does not himself remain on these heights.  At once
he takes up again the work of thought, without which that vision
would have no significance, although it itself stands above all
thinking and science (Symp. 211 A).  The Highest is not
merely that which makes knowledge possible, but is at the same time that which
makes the knowing beings exist, thus the
ultimate ground of all thought and of all being (Republic VII
p. 508 E; 509 B).

In relation to what can thus appear at
the high points of the life of thought, although it itself lies beyond
all thinking, faith in gods is for Plato something subordinate.
The divinity of the gods rests only on their
participation in that which is also for human beings the highest.
The true thinker lives in ``that in which it is God's
divinity to live'' (Phaedrus 249 C); the difference

% --- p. 44 ---
\setcounter{footnote}{0}%
\opage{44}is only that what for the gods is a constant possession
is for human beings an object of constant striving (Phaedrus
278 D).

When we go from Plato to Plotinus, the
mystical element appears decidedly in greater fullness, takes up and satisfies
the craving for life in a far stronger way than
in Plato.  The method is the same: ascent along the path of thought
to the Highest, which is more than thought.  Plotinus's
biographer Porphyry says expressly of him that he
followed the path Plato had indicated in the Symposium.  But
although this path could be long, the goal that
was reached was nevertheless not a distant and otherworldly one; the goal, ``unity with
and nearness to the God who works in all,'' was always near
(Vita Plotini. Cap. 23).  This immediate presence,
which is far more valuable than science (παρουσία κρείττων. Ennead. VI 19, 4), can be
reached only through the work of thought.  It is thus at once
distant and near.  But the nearness, the immediate
unity with the Divine, presupposes a mental
state raised above the ordinary life of consciousness,
--- an ecstasy.

The ecstatic state is described by Plotinus not merely,
as by Plato, as inspiration, enthusiasm, but as
stillness and rest.  At the high point there is no longer,
as in Plato, any talk of beholding.  The soul is ``outside of
itself'' (ἔκστασις), and it enters a state that is at
once simplicity (ἅπλωσις) and unfolding (ἐπίδοσις).
The contrast between the seer and the vision falls away.  It
is as if two centers coincide, and it is as
if the human being had become another.  That is why it is so
difficult to give a description of what is experienced in
such states (Ennead. VI, 9, 10---11---11).

Of great interest is the consequence that Plotinus draws
from the circumstance that the Highest is not reached through
% --- p. 45 ---
\setcounter{footnote}{0}%
\opage{45}science, but makes itself known through immediate experience
of an indescribable character.  When we wish to communicate
such an experience to others through speech or writing, we
can only lead them to it indirectly.  As long as we
are on the way, there can be talk of instruction; but
the goal itself presupposes a personal craving, without which
the soul is not irradiated by the highest light (Ennead.
IX, 6, 41).  It is already apparent here how mysticism
leads to an emphasis on individual experiences [Erfaringer], in opposition
both to scientific research and to handed-down
dogmatics.  When one of Plotinus's friends, who was an
eager participant in the temple festivals, wanted Plotinus to
accompany him to such a festival, he answered: ``It is
the gods who must come to me, not I who must
go to the gods.''  (Vita Plotini. Cap. 10).  Plotinus agrees
with Plato that human beings can have access by their own
way to that on which the divinity of the gods
rests, or, as he expresses it (Ennead. IX, 1, 2),
``The cause of gods being gods must be a God
who stands above them in dignity.''  The concept of divinity
appears in such utterances clearly as a
determination of value.

The reason why the highest experience according to Plotinus
takes on the character of ecstasy, and cannot be an object of
description, is first that what is to be known
lies beyond all oppositions, whereas knowledge
comes about only by thought's passing from one thing
to another.  The absolute unity can be experienced [erfares] only
through rest and complete concentration.  But second,
it is because this unity itself contains in itself
an infinite fullness.  The two determinations belong closely together,
for if there were anything outside the Godhead,
it would be limited and could not be an absolute unity.
(VI, 9, 4---6).  For both reasons the Highest is
% --- p. 46 ---
\setcounter{footnote}{0}%
\opage{46}inexpressible.  For any predicate one might apply to
it, one would make it something single, determinate and
limited.  (V, 3, 13).  Plotinus here sets up a fundamental principle
of all mysticism\footnote[1]{See on this my \textit{Religionsfilosofi}, § 21 and Note 57.}.

Neoplatonism becomes foundational for the mysticism of the Middle Ages,
on which it gains influence partly through
Augustine, partly through Dionysius the Areopagite.  Whereas
medieval thought rested essentially on Aristotle's
encyclopedic production, as far as knowledge of the world
is concerned, it was Neoplatonic ideas that conditioned
the concluding spiritual content.  The relation of medieval
mysticism to scholasticism is conditioned by the fact that the
energetic absorption in, and being seized by, this spiritual content
came to form an opposition to scholastic
dialectic.  To use Plotinus's expression: the
immediate presence of the Divine (Parusi) was set
above science (Episteme) of the Godhead.

\sectionhead{4. \textit{Augustine.}}

The great Church Father too had ecstatic experiences,
and even if they were not brought about by Neoplatonic
influence, yet he has, what is in this
connection of special interest, interpreted them in Neoplatonic
forms.  From this time on, two ways of
interpretation contend, on the whole, over extraordinary, especially ecstatic
experiences, --- the animistic (in the widest sense), which
finds the explanation in the intervention of souls or personal beings,
and the Platonic, which at the high point assumes
a revelation of something that cannot find
expression in any image or any concept.  The turning point
is marked by Plotinus and Augustine.

Augustine owes, as he has himself repeatedly
% --- p. 47 ---
\setcounter{footnote}{0}%
\opage{47}declared, to Plotinus as representative of Platonism,
an extraordinary amount.  Platonism led him away from
materialistic ideas.  After a hard struggle he gave up
imagining God in bodily form.  Next
he was led by his Platonic (Neoplatonic) studies
to self-knowledge, to the study of his own inner life.  And finally
he was led in Plotinus's footsteps step by step away from
the accustomed ideas and from the world of sensation up
to the Unchangeable, the true being, which can only be
grasped in mystical beholding, --- or, as he himself expresses
it, ``in the flash of a trembling glance'' (in ictu trepidantis
aspectus).  To hold fast this vision he could not, and
when he was again in his daily state (redditus solitis),
it was only a ``loving memory'' (amans memoria),
a mood of joy and gratitude, that he retained ---
a craving for the scent of the food he had not been able
to make his own\footnote[1]{\textit{Confessiones} VII, 23 (17). --- By the expression amans memoria cannot be understood memory in the proper sense, but the mood that remains after a vision whose content cannot be held fast.}.  Indeed, after his conversion, he experienced again in
conversation with his mother on the day of her death an ecstasy,
as mother and son were gradually led in speech and
thought beyond everything sensible --- and they attained a gentle
touch of the Eternal in the full rapture of the heart
(attingimus [sapientiam] modice toto ictu cordis); but
the firstfruits of the spirit, in which they had become partakers,
they had to leave and return again to what the
ordinary language could express.  And he attempts a
translation of it which, when one looks closer, proves
to be a translation of passages in Plotinus: When
all is silent, within us and outside us, and the soul itself is silent
and steps out beyond itself by non-thinking, then
God reveals himself by himself, without voice and without sight, and
% --- p. 48 ---
\setcounter{footnote}{0}%
\opage{48}there comes a moment of understanding (momentum intelligentiæ),
in which a swift act of thought touches the eternal
wisdom. --- Augustine adds that approximately in this way
he had spoken, though not in the same words.  And he
adds further: might not that which is experienced in such a
moment be what is expressed in the words: ``Enter
into the joy of thy Lord?''\footnote[1]{\textit{Confessiones} IX, 23—26 (10). Cf. \emph{Plotinus}: \textit{Ennead}. V, 1, 2 and III, 8, 9. --- See \textit{Religionsfilosofi}, § 61.} --- This last is characteristic
of the fusion of Platonism and animism
on which all theology rests, as well as of mysticism,
as it developed under the influence of Augustine
in the medieval church.  To illuminate
the relation between experience and interpretation the sections
of Augustine's Confessions used here are of great interest.


Naturally Augustine does not identify Platonism and
Christianity.  Platonism could only see the goal, but
not go the way to the end\footnote[2]{\textit{Confessiones} VII, 27 (21).}.  In his old age
he reproached himself in his ``Retractationes'' with having
placed Platonism and Christianity too near each other.

But Plotinus has, among much else, taught him to
distinguish sharply between the immediate experience
and its interpretation.  This lesson was not wholly forgotten
in the history of mysticism, though the stricter formulation
of the dogmas under ecclesiastical authority could, for
both inner and outer reasons, make it difficult to
hold it fast.

\sectionhead{5. \textit{Medieval Mysticism.}}

Although medieval theology sought to take up the
knowledge of antiquity, especially as it had been summed up
% --- p. 49 ---
\setcounter{footnote}{0}%
\opage{49}in the works of Aristotle, Galen and Ptolemy, the
goal it strove for lay nevertheless beyond all knowledge and thought
in immediate union with the Godhead.  One has
therefore been able to say with a certain justice that not Aristotle,
but Plotinus is the true teacher of the medieval
philosophers\footnote[1]{Fr. \emph{Picavet}: \textit{Esquisse d'une Histoire comparée des philosophies médiévales}. Chap. 5.}.
The medieval organization of studies,
which begins with encyclopedia and philosophy and ends in
mysticism, is really already found in the Neoplatonic
school at Alexandria, which used the writings of Aristotle
as an introduction and preparation to the study of Plato,
whose line of thought was interpreted in a mystical direction\footnote[2]{Cf. \emph{William Norvin}: \textit{Olympiodoros fra Alexandria}. (1915). Chap. 1.}.  Mysticism
did not necessarily stand in opposition to scholasticism,
but could stand as its conclusion.  It has indeed
also been shown that the boldest of the medieval
mystics, Eckehart, himself wrote scholastic
works\footnote[3]{\emph{Denifle}: \textit{Meister Eckehart's lateinische Schriften und die Grundanschauungen seiner Lehre} (Archiv für Literatur- und Kirchengeschichte des Mittelalters) (1886). --- Even earlier (1838) \emph{Engelhardt} in his book on Richard a St. Victore and Johannes Ruysbrook (p. 97) had shown that mystics and Scholastics did not belong to two wholly different worlds.}.  But there was in mystical spirits a
tendency to build on the inner life and its experiences
and to set them far above analyzing and distinguishing
thought, a tendency that could easily lead to conflict
with the ecclesiastically sanctioned scholasticism.  And in the
school of the inner life the mystics learned, what became still
more dangerous, to distinguish between experience and interpretation.  Often
it was only in a convulsive way that they were able to
unite their personal experiences [Erfaringer] with the ecclesiastical
interpretation.  The opposition between mystics and
% --- p. 50 ---
\setcounter{footnote}{0}%
\opage{50}Scholastics rests, as I have shown in my Religionsfilosofi\footnote[1]{\textit{Religionsfilosofi} § 21 (with Note 37). --- On this point too \emph{Plotinus}, and after him \emph{Augustine}, stand as predecessors, in that they maintain that the fundamental scientific concepts (the categories) have validity only for the sensible world, not for the Godhead, in whom no distinction can be drawn between essence and attribute or between attributes among themselves. See on this \textit{Den menneskelige Tanke}. p. 143.}, in respect of epistemology on the fact that the latter maintain
the validity of the analogies by which the Godhead
alone can be thought, whereas the former deny this validity.
The opposition between mysticism and scholasticism thus depends
not only on the opposition between experience
and interpretation, but also on the opposition between
criticism and dogmatism.

A few single examples from the history of medieval
mysticism shall here be brought forward to illuminate our
subject.

By Hugo a St. Victore (at the beginning of the 12th century)
ecstasy is described in the following way.  After
having passed through various stages the soul goes as it were
out of itself and is carried away, without knowing whither.
``Then I become,'' it says, ``suddenly renewed and transformed
and feel a joy that I cannot describe.  The conscience
is cleared; all sorrow over suffering endured is forgotten;
the heart rejoices; the understanding becomes clear, the mind
illumined.  Longing satisfied;\dots{} and there is something
I hold fast with inward love's arm,
without knowing what it is.''  To this the answer is: ``It is thy
Beloved (dilectus tuus) who visits thee; but he comes
invisible, hidden and incomprehensible.  He comes
to touch thee, not to be seen by thee.''  Clearly a distinction is here drawn,
as in Plotinus, between ecstatic experience and visionary
beholding.  And the ecclesiastical doctrine of a personal
God is used, with allusions to the erotic
% --- p. 51 ---
\setcounter{footnote}{0}%
\opage{51}language of ``the Song of Songs,'' to interpret the indescribable that is experienced. --- As
stages leading up to the high point where the ecstatic
experience becomes possible, Hugo names four: reading
(for devotion) (lectio); reflection (meditatio); prayer
(oratio); good deeds (operatio).  The high point is called
contemplation (contemplatio), which is one with ecstasy;
it is different from reflection and signifies
nothing intellectual\footnote[1]{A. \emph{Liebner}: \textit{Hugo von St. Victor und die theologischen Richtungen seiner Zeit} (1832) p. 270—276. --- The various stages that Hugo indicates are already found in similar characterization among the mystics of earlier times, not only among the Neoplatonists, but also in Indian mysticism (Yoga and Buddhism). \emph{Friederich Heiler} \textit{(Das Gebet}. 1918 p. 253—258) gives a (probably schematic and detached-quotation-based) comparison (Tafel der Gebets- und Versenkungsstufen), which shows how the mystics of all times know approximately the same stages of absorption. Much probably rests here on tradition; but that the tradition has become more than dead transmission can be explained only by new, independent experience. And this becomes clear when one goes back from the identical expressions and utterances to the way in which they are motivated in the individual mystics.}. --- This series of stages points back through
many intermediate links to Plato's Symposium and
recurs with various nuances in other mystics.
Thus especially in Richard a St. Victore (in the middle
of the 12th century), who not only describes the stages that precede
``contemplation,'' but also the highest
stage itself more thoroughly and with more distinctions.
Of interest for our subject is the difference that is
drawn between ``speculation,'' in which we see through a
mirror, and ``contemplation,'' in which we behold the truth
without a veil.  But it is precisely this immediacy
that makes the state of contemplation into ecstasy\footnote[2]{H. O. \emph{Taylor} \textit{(The Mediaeval Mind. A History of the development of Thought and Emotion in the Middle Ages} (1911 II p. 363) aptly speaks of the ``irrational immediacy of conviction,'' by which the mystic feels and knows his object.}.  This
% --- p. 52 ---
\setcounter{footnote}{0}%
\opage{52}state appears after longing, anxiety
and repentance have been felt, and it has the character of a peaceful stillness,
a gentle breeze, in which the human being is carried beyond
himself and sees the light of the highest wisdom without a veil.  But
what has thus been beheld in the state of rapture
is afterwards to be considered and thought through, so that it
can be understood, and so that by means of reasons and comparisons
it can be brought to general understanding.  And
the mysteries of faith, which are known by revelation
or taught by theologians, these one is to bring into contemplation,
so that they may become objects of admiration
and enjoyment.\footnote[1]{\emph{Engelhardt}: \textit{Richard von St. Victor und Johannes Ruysbrook}. p. 81.}

In these last sentences the relation between experience and interpretation,
fateful for all mysticism, is hinted at.
The mystic is convinced that it is God he has
been in contact with during the ecstasy; but when no
analogies suffice, how then explain it to others,
and how recognize the relation between what is experienced
and what the church teaches?  Richard indeed remedies
the difficulty by urging that Bible and church doctrine be
put into the ecstatic state itself,
in order to admire the doctrine so much the more (in that
it thus enjoys the benefit of the state of enthusiasm).  But
this reading in can mean nothing other than a hasty
interpretation.  In the ecstasy itself one is not able to
put anything in, although the ideas in which one otherwise
lives will naturally come to determine the visions
that can attach themselves to the ecstasy.

With Master Eckehart (about the year 1300) the difference
between experience and interpretation recedes, not because
he follows Richard's counsel to put theology into the
immediate experience, but precisely the reverse, because he
% --- p. 53 ---
\setcounter{footnote}{0}%
\opage{53}at once and decisively interprets the dogmatic ideas
from the mystical experience.  As long as the soul is occupied
with what is distant in time and space, and as long as differences
and oppositions in general make themselves felt
to it, so long it is not as it ought to be, not
raised above all sorrow.  What matters is to experience in an eternal Now
the whole fullness of existence.  That God sends his Son into
the world is not a distant event, but something
that happens in the innermost of the soul, where God's being and man's
own being are one.  Blessedness rests neither on
knowledge nor on love, but on that from which
both knowledge and love spring, and where there is
neither a past nor a future.  For
him who has even once glimpsed this depth, the
greatest material goods lose their value, and there is on the whole
no external end for action, not even God
or eternal blessedness; one acts out of the soul's
innermost depth, where God dwells --- nearer to the soul than
the soul is to itself.  One is not here to ask whether it
is through grace or through nature that God works in
the soul, both are one.  The proof that God works in the soul
consists in the happiness and rest that arise.  There is
an immediate way to unity with God, and the masters in
the schools wrongly make themselves great scruples over
how it can be possible for man to become
one with God.  It demands renunciation and abstinence; but
God is in us; it is only we who are not at home with
ourselves.\footnote[1]{In this characterization I have used the translation of selected writings of Eckehart that \emph{Aage Marcus} has given. (\textit{Mester Eckehart; Prædikener og Traktater}. Oversat og indledet. 1917).}.)

It was no wonder that the church became uneasy
about such utterances.  Eckehart was accused
of heresy, and 28 propositions that were extracted from his
% --- p. 54 ---
\setcounter{footnote}{0}%
\opage{54}writings were condemned.  Eckehart himself, however, did not live to see
the condemnation.  In the Catholic world Eckehart is not highly
esteemed.  Denifle seeks to reduce him to a
preacher for pious ladies and derives his mysticism
from his ignorance of dogmatics.

Eckehart is influenced by Augustine and (through Dionysius
and Erigena) by Neoplatonism.  But clearly there lies
at the bottom an original mental experience.  He
has had moments filled with that peace and rest
which is possible only through absorption in one's inner life and in
the life that can there involuntarily well up.  And as far
as one can see, these moments have for a substantial part
been conditioned by an intellectual interest in working
oneself out beyond the differences and oppositions with which
the life of thought directed against experience [Erfaringen] has to contend.
But this intellectual interest has for Eckehart gone
hand in hand with a striving after attachment to, or rather union
with, something that is raised above strife and division.
And when he without further ado uses the word ``God'' of that
which stirs innermost in him, the concept of God
appears clearly as a predicate concept.  The experience is of
the highest value; therefore it is to be called divine;
and therefore it must signify God's indwelling.  We meet
here the Neoplatonic thought which above (p. 46) was
called the foundation of all religious mysticism.

Moreover it will always turn out that when the question
of what God is is seriously raised, it is always the
inner or outer experience that human beings have that dictates
the answer\footnote[1]{Cf. \textit{Upanishad}. Übers. v, Deussen. p. 204 f, 395, 799 f. --- \emph{Augustine}: \textit{Confessiones}. Book 10. --- \emph{Luther}: \textit{Catechismus Major}. (The explanation of the first Mosaic commandment). (Libri Symbol. ed. Hase p. 403). --- \emph{Spinoza}: \textit{Tract. theol. polit}. Cap. 14. (Deus: veræ vitæ exemplar).}.  But from the involuntarily given answer there is again

% --- p. 55 ---
\setcounter{footnote}{0}%
\opage{55}a distance to the answer that the religious tradition, within
which the individual moves, will give.

What the church condemned in Eckehart was the strict
carrying through of the thought of unity on the basis of what
had been experienced in ecstatic states about freedom from all
differences and all barriers.  His disciple Suso also incurred
an ecclesiastical condemnation; but since he was
less bold in his thinking, and more careful
to remain in agreement with church doctrine, he avoided
greater trouble.  And he gains precisely through this striving
for ecclesiastical interpretation interest for our subject,
in that he is thereby led to distinguish sharply between experience
and interpretation.

In the 54th chapter of his Life he describes ecstatic
states which for a time occurred every morning and every
evening.  ``He then sank so wholly into God, the eternal
Wisdom, that he could not speak of it.  Sometimes
he had a loving conversation with God; at other
times he fell into wailing and sighing or into longing
weeping or into quiet laughter, often it was to him as
if he hovered in the air or swam between time
and eternity in the deep sea of God's unfathomable wonders.
Sometimes his heart was so full that he laid his hand
on it and said: Ah, dear heart, what is to
become of thee!''  --- This ecstasy has rather the character
of a vision.  In an earlier ecstasy, the first he
had, the visionary recedes more before the emotional.
In the 3rd chapter of his Life he relates that as
he, suffering and disconsolate, stood alone one day in the church,
his soul was rapt --- whether it was in the body
or outside the body --- and he saw and heard what
no tongue can utter: ``It was without form and shape
and yet caused as great a delight as all forms
and shapes together cause.  The heart yearned and was
% --- p. 56 ---
\setcounter{footnote}{0}%
\opage{56}yet satisfied; the mind was glad and content; every wish
and all desire was satisfied.  He was wholly absorbed in the gazing
toward a radiant light, which made him forget himself
and all things.  Was it not night?  that he did not
know.  It was a lovely breaking forth of the eternal life in the present,
quiet receiving.  He then said: If this is not
Heaven, then I do not know what Heaven is!\dots{}  When
he came to himself again, he felt like a person
who comes from another world.''  After-effects
of this experience persisted long, as a ``heavenly
taste,'' just as a box in which one has kept
a precious substance retains the taste of it.  It is Augustine's
``loving memory'' (amans memoria) (see above
p. 47) that we here meet again.

Like Augustine, Suso experienced a great contrast
between the ecstatic moments and the ordinary conditions.
And there was a great difference between what he wrote
in his ordinary state and what was given him
in the ecstatic states.  But even this latter he subjected
to a thorough criticism, so that nothing should
be found in it that departed from the utterances of the holy Fathers,
but that what the divine Wisdom imparted to him
through spiritual illumination could be supported by the authority of Holy
Scripture, preferably in such a way that it was both in word
and in meaning one with its utterance\footnote[1]{Et quum ea, quæ data erant, ad tabulam concepisset, ad se reversus studiose rimatus est, ut nihil in iis esset, qvod a dictis sanctorum Patrum discreparet, sed hoc, qood divina tradente Sapientia sibi per spiritualem illustrationem infundebatur, auctoritate sacræ scripturæ\dots{} fulciretur, quum idem in sensu, aut etiam idem in verbis et sensu\dots{} reperiretur esse. (\textit{Heinrich Suso's Leben und Schriften}. Herausg. von Diepenbrock 1829. p. VIII f.)}.

This utterance (which I have already used in my
Religionsfilosofi § 62) is of classical significance as regards the
% --- p. 57 ---
\setcounter{footnote}{0}%
\opage{57}relation between personal experience and ecclesiastical authority.
It clearly illuminates the difference between the psychological
and the sociological conception of the essence of religion.

Socially and ecclesiastically emphasis is laid on regularity
and definite, recurring forms, and here
the ecstatic experiences are threatened with danger if they are not carefully
interpreted according to ecclesiastical dogmatics or at least
expressed in accordance with ecclesiastical usage.
There stirs in such states a restless, and
in any case a decidedly individual element.  The church
must hold fast that all experiences of life can come under
its once-for-all formed circle of images and ideas,
and it must see a danger in the fact that the most valuable experiences
should not be capable of being interpreted exhaustively by means of
it. ---

I now pass to the main subject of this essay,
a comparison between mysticism, as it in the
sixteenth century appears in a Spanish nun,
and mysticism, as it in the twentieth century appears
in a Protestant woman teacher.  The foregoing
characterizations were to serve only as an introduction to it.

% --- p. 58 ---
\setcounter{footnote}{0}%
\opage{58}\chapterhead{III.}{SAINT TERESA}

\sectionhead{1. \textit{The First Stages.}}

Teresa de Cepeda was born in 1515 in the city of Avila in
Spain, of a noble family. In her autobiography, which she wrote
at the command of her confessor, she tells of her
religious development. Religious romanticism stirred in her
very early; at seven years of age she wanted, together with
her brother, to go to the Moors in order to become a martyr
and thereby reach Heaven the sooner. She belonged to a
strictly Catholic family; yet she confesses that she
had a great desire to adorn herself and a great desire to
read romances of chivalry. From her fourteenth year she was
brought up in a convent, but felt at the beginning great reluctance
toward the life of the cloister, and she chose it only because the privations
and sufferings it entailed were less than those
Purgatory would bring. But by degrees she came into
a more inward relation to the life that was lived in the convent,
and when she was consecrated a nun, nineteen years old,
she felt herself as if irradiated by a divine
light. Very early the contrast between her
inner life and the ecclesiastical representations of her confessors
came to the fore. She had a long struggle to wage here, for her
ecstatic and visionary experiences were often taken by her confessors
to be inspirations of the Devil, who
ought to be brought before the court of the Inquisition.  But apart from
% --- p. 59 ---
\setcounter{footnote}{0}%
\opage{59}the times of trial she felt herself so certain that her
inner experiences were of divine origin,
that she would not be shaken by the warnings she received in
the confessional. She kept changing confessors until
she found one who understood her. She speaks sharp words
of half-learned confessors, ignorant of the inner life.
A confessor she regarded as necessary, for she was
a Catholic in her inmost soul and wished to live and die in accordance
with the doctrine of the one Church that brings salvation.
And she felt her lack of learning when it was a question
of finding the right interpretation of what she experienced.
There was here, she thought, a danger that might threaten women
in particular. They must have an authority to lean on in order
not to go astray, but it must be men
who from their own experience [Erfaring] know ``the spiritual ways.'' And
one should therefore choose one's spiritual guide with care
and after trial. It was among the Jesuits that she found
what she sought. Their Master had himself had ecstasies
and visions, and yet no one had done as much as he for
the fortification of the Catholic Church.Teresa's Jesuit confessor
was convinced that her remarkable experiences were of
divine origin. Only he thought that she lacked
self-knowledge and humility, and he had her go through
St.\ Ignatius' ``spiritual exercises.'' This school
trained the imagination in picturing the bliss of Heaven
and the horror of Hell and in living itself into the life of Jesus in humiliation
and exaltation, but aimed above all at strengthening
obedience to the holy Church and the will to follow
the goal of life that it set\footnote[1]{A thorough characterization of Ignatius' \textit{Exercitia spiritualia} is given by \emph{Boehmer}: \textit{Les Jésuites}. (Trad. frane. par Monod. p. 25---38). Cf.\ also \emph{Henri Joly}: \textit{St. Ignace de Loyola}. (1899).}.

It was through prayer that Teresa gained access to
the spiritual world, which unfolded ever more richly
% --- p. 60 ---
\setcounter{footnote}{0}%
\opage{60}and more inwardly for her, until in the course of 2 or
3 years (from her 44th year) she lived through a series
of ecstasies and visions, which undoubtedly mark the
culmination of her spiritual life.

There are three main sources for Teresa's inner experiences,
--- her autobiography, which she wrote down and rewrote
in the years from 1562 to 1566\footnote[1]{\textit{Vie de Sainte Térèse} écrite par elle même. Trad. Brouix 14. éd. Paris 1896.}, her letters\footnote[2]{\textit{Lettres de Sainte Térèse}. Trad. Brouix. 2 éd.} and her
writings\footnote[3]{\textit{Les Oeuvres de Sainte Térèse}. Trad. Arnauld d'Andilly. Paris 1670.}, especially the book on the ``Castle of the Soul'' (Château de
l'âme 1577). ---

Teresa confesses that in her first time in the convent
she had a great desire for distraction, especially for conversation with
visitors. She attended to her exercises and said the prescribed
prayers, but it was only the mouth, not the
heart, that prayed, and when she lay in prayer in the church, seemingly
full of devotion, she yearned nevertheless in her inmost self
to hear the stroke of the bell that announced the end
of the hour of prayer. She thus belonged in this period to those of whom
she later says in her book on the ``Castle of the Soul'' that
they knew only the castle that their own soul is, in the same
way as the guard who makes his round outside the castle,
but never comes into it. The door to the castle of the soul is
the prayer that does not consist merely in words, but in which thought
is present, and in which the whole mind is present: oraison mentale,
as opposed to oraison verbale.

The first dwelling in the castle of the soul is the home of
self-knowledge. But one learns to know oneself best by
knowing God; his greatness makes man recognize
his own smallness. In the next two dwellings God speaks to
man, not only through prayer, but also through
the words of good people, through books and through
% --- p. 61 ---
\setcounter{footnote}{0}%
\opage{61}sickness and sorrow. The need for penance arises, and love
of one's neighbor awakens. To avoid false security
even those who do not enter a convent ought to have a person
to whose will they submit themselves, so that they do not act
according to their own will, but learn full humility.
And only then comes --- in the fourth dwelling of the castle --- the
stage at which prayer truly takes on an inward character and
through the prayer of reflection (oraison de méditation) passes
over into becoming the prayer of quiet (oraison de quiétude). So long as
reflection still rules, full collection and
peace are not attained. But in the midst of contemplation or reading
it may happen that the soul feels itself gripped by the presence of God,
sinks into him as into a divine
abyss. In such absorption neither
thought nor imagination is at work; memory is almost vanished;
the will is absorbed in love, and the understanding is seized
by the greatness that reveals itself.

To a state like this one cannot raise
oneself. One must wait until one is drawn upward. Therein consists
precisely the difference between the prayer of quiet and the prayer of meditation.
Meditation is to be compared with the watering
of a garden that is effected by the water being led through pipes in over
the garden's plants, or by one's laboriously
drawing the water up out of a well. The spring from which the water
is fetched is God himself. But in the prayer of quiet the spiritual
watering takes place immediately, as when a sweet rain streams
down over our garden. It takes place quite quietly; but
when our own reflection and striving cooperate, it
does not pass off without noise. ---

In the book on the ``Castle of the Soul'' Teresa says that she now
sees these things more clearly than 15 years before (that is
to say, than when she began to write her autobiography).
She is thinking presumably above all of the psychological description
of her remarkable states. She had earlier grieved
% --- p. 62 ---
\setcounter{footnote}{0}%
\opage{62}that even in states in which all her faculties
ought to have been concentrated on the thought of God, her
ideas could flutter about in highly diverse regions,
often among the unclean and poisonous animals that romped outside
the castle, in whose inner dwellings she believed that her
soul had now taken up residence. She now understood that there were
two kinds of thought. The one goes from subject to subject as
the imagination goes from image to image; the other
consists in concentration on a single point and is at bottom
one with the will to hold fast a definite direction
in longing and striving. All thought has therefore not, as she
had formerly said, vanished in the ``prayer of quiet.'' It is
the difference between discursive thinking and instinctive concentration
that now stands more clearly before her. After
she had had this confirmed by ``a very learned man,''
she no longer let herself be troubled by the fact that there might be
noise and unrest in the inner and outer regions of the soul, while
in her inmost longing she gathered her mind about the one,
great, divine thought.

Other psychological experiences, which also seem to
originate from a later time than the autobiography, are likewise
of interest for the psychology of the mystical states.

It does not help to practice self-compulsion in order to bring about
the state that is called the ``prayer of quiet.'' This remark
is directed against certain ``Contemplatives'' who maintain
that one should turn with attention toward what
God wills to work in us, while one must not let reflection
work. Such preparation, according to Teresa, helps
not at all. Before we have turned the mind in the direction of
God, God may already have led us into the castle of the soul,
without our knowing how we have come in there. And if
preparation does not help, why then not use our
reflection for something useful, instead of exerting oneself
to refrain from thinking! Use your faculties in the best
% --- p. 63 ---
\setcounter{footnote}{0}%
\opage{63}way, until God himself brings forward such lofty subjects before
you that thought sinks away in holy seizure! --- By
these remarks a greater difference is evidently set between the prayer of
reflection (oraison de méditation)
and the prayer of quiet (oraison de quiétude) than she had
set before. This point is of all the greater importance
in that the mystics and ecstatics of all times saw, precisely
in the sudden, unprepared and involuntary manner in
which their highest experiences arose, a proof
that they had a supernatural origin, --- whether
God or the Devil was this supernatural
cause.

The ``prayer of quiet,'' despite its sweetness, is not the
highest spiritual state. And it may contain a
danger, in that it may leave behind a period of passivity, of
sinking away, which has a particular attraction especially for melancholy
natures, so that they revel in such states
and seek to bring them about and to maintain them by
long-continued prayer, vigils and chastisement. They thereby think
to attain a continuation of the state of the ``prayer of quiet,''
which is the work of God, and are perhaps even inclined
to confuse what they experience in this way
with the highest (not yet mentioned) ecstatic states.
The enjoyment of a natural state is confused with the joy
of a supernatural state. And the Devil surely
has a finger in the game, in that he tempts weak natures
to attain through chastisement what must come
as a gift. Chastisements with this aim ought therefore to be
forbidden by the confessor or the prioress.

\sectionhead{2. \textit{The Highest Stage.}}

The highest form of prayer is the prayer of union (oraison
d'union). In the prayer of quiet the soul's longing was indeed
% --- p. 64 ---
\setcounter{footnote}{0}%
\opage{64}concentrated on a single subject, but still there was yet
the possibility that memory and reflection might romp
about. This possibility has vanished in the prayer of union.
The feeling of union here becomes ecstasy, which is soon described
as elevation (élévation), soon as rapture (ravissement).
Ecstasy is related to the mere feeling of union
as a bonfire is related to a gentle fire. In the
gentle fire the iron still lies as a black mass, but
in the bonfire it becomes glowing. --- In the book on the ``Castle of the Soul''
ecstasy is described in a manner similar to that in
the autobiography; it belongs there in the fifth
dwelling.

The description of this leads Teresa to some
interesting remarks. One notices that her attention
to nuances and transitions has been sharpened
in the school of experiences.

a. At its sudden onset, as well as at its
culmination, the ecstatic state is quite simple
and plain. The soul has no feeling and no will; and
it knows not what it is filled with, --- knows not what
it loves. An immediate presence makes itself known,
without anything being seen or heard. The state is therefore
indescribable. And yet it is a state of immediate
certainty. If one wishes to ask how there can be
certainty when nothing is either seen or heard, the answer is
that it is one of the secrets of Omnipotence. It is like
a nourishment that we notice in ourselves without having eaten it.
Afterward, when the strongest ecstasy is over, there may
arise visions and auditions, which can be the subject
of description. And later again one may try, through
express comparison, to make the experience clear to oneself.

The simple certainty of a presence may be called a
vision, but one must then distinguish between spiritual vision
(vision intellectuelle) (and audition), which is without images
% --- p. 65 ---
\setcounter{footnote}{0}%
\opage{65}and without words\footnote[1]{In a letter of 1576 Teresa describes vision intellectuelle as the apprehension of a presence, without any inner or outer vision being connected with it. One notices on which side the object is more distinctly than if one saw it. The immediate recognition too is mentioned in this letter. \textit{Lettres de Sainte Térèse}\textsuperscript{2}. Trad. par Brouix. 1882. I p. 366.}, and the vision proper (visions imaginatives)
(and audition). Auditory hallucinations often arise
before visual hallucinations. When Teresa asked who it
was that made himself known to her in an intellectual vision
(hallucination of presence), Jesus made himself known by
speaking to her. She had, however, beforehand (by the mere
presence) been certain that it was he.

b. The visions proper (the imaginative ones), in which something
is seen or heard, often develop successively. Thus
Jesus, after having first acted by his mere, indeterminate
existence, let her see first his hands, and then
the rest of his body.

As to these pictorial visions themselves, Teresa distinguishes
between those in which she saw Jesus himself,
and those in which she saw only an image of him.
--- This difference seems to correspond to that which has been made
between hallucination and pseudohallucination. In pseudohallucination
the object stands as it were with a symbolic
character, without direct reality. ---

The successive arising of visions is of great interest,
because --- as already remarked above ---
it opens greater access for ideas that the individual
has had earlier to gain influence on the content
of the visions. Thereby it is understood why visions on the whole
agree with the whole circle of experience and
ideas of the visionaries. Already where there is no pictorial vision,
but where, as with Teresa's ``intellectual
vision,'' the immediate sensation or feeling of
% --- p. 66 ---
\setcounter{footnote}{0}%
\opage{66}a presence fills the whole state, there begins an
interpretation in the form of recognition: Teresa was certain
that it was Jesus who was present, though she as yet
heard and saw nothing! --- We shall later see an example of
the immediate sensation or feeling of a presence
that retains its indeterminate form, without any recognition.
It depends on how near and how firm
the interpreting ideas lie.

c. The enthusiastic and enthralling rapture may
be connected with pain, because it comes so
overwhelmingly, and because the divine influence is
like a consuming fire. During the very refreshment
pain may be suffered --- also because the longing is so deep that
it asserts itself again and again. Once, when Teresa
had sunk away, seized by privation and longing, she saw
at her side a cherub, who pierced her with
a glowing iron arrow, which was thrust through her heart
down into her entrails, which seemed to follow when the
arrow was drawn out. ``The pain of this wound was,''
she says, ``so vivid that it drew forth moaning;
but this unutterable martyrdom made me at the same time
feel the sweetest joy, and I could not wish that it
should cease. It is not a bodily suffering, but
wholly spiritual, though the body has a share in it.''

It would surely be wrong to take this description
as testimony to the erotic-sexual origin of religious ecstasies.
The Freudian school cannot without more ado
claim it as support. The relation is rather
that the violent being-seized has as it were streamed out
over the central organs and has taken possession of the whole body,
so that ``the entrails'' too have come along. In my Psykologi
I have called this phenomenon the diffusion of the feelings. In such a diffusion, which can be observed in
all strong emotions of an inward and pervasive
% --- p. 67 ---
\setcounter{footnote}{0}%
\opage{67}kind, lies one of the reasons why the erotic-sexual
language lies so near at hand when the mystics seek expression
for their experiences\footnote[1]{Cf.\ \textit{Religionsfilosofi} § 62.}. The question is, however, complicated, and
another mystical phenomenon will in a later section give
occasion to return to it.

\sectionhead{3. \textit{Teresa and the Doctrine of the Church.}}

Only on a few occasions does Teresa hint at a difference between
experience and interpretation. As a rule there was connected with
rapture an immediate certainty, which stood before
her herself afterward as a wonder. And likewise, even if nothing was seen or heard, an immediate
recognition asserted itself, which could be further unfolded and
motivated through auditions and visions. When she
comes to raise the question of what stirred
in her in the ecstatic and visionary states,
it was a psychological rather than a theological problem
that she posed. It concerned the inner experiences in the soul in
such states (les sentiments intérieurs dans cet
état), --- what the soul did while they lasted
(ce que l'âme pouvait faire pendant ce temps). The question
presented itself to her when, after the ecstasy
had ceased and she had received communion, she prepared
to write about it. But she could not become clear about it
and declared it a divine secret.
Especially the experience of a presence without vision
and audition was enigmatic. But --- she adds --- something
she must nevertheless have been able to grasp from these states,
otherwise she could not even have expressed herself about them
by means of comparisons.

To express how the soul can sometimes
raise itself above itself without ceasing to be itself
% --- p. 68 ---
\setcounter{footnote}{0}%
\opage{68}she says that it is like a hearth fire, which flaming
rises into the air: the more violent the fire is, the
higher the flame rises above the hearth, without however
ceasing to be the hearth's fire. This image, which is meant to illustrate
the relation between the ``soul'' and the ``spirit,''
bears, as it appears in the eighth chapter
of the autobiography, the stamp of an involuntary symbol. By contrast
a more arbitrary comparison appears in
the image that is carried out both in the book on the ``Castle of the Soul''
and in the autobiography, of the garden and the different ways
in which it can be watered. It is a widely spun-out
allegory. Another allegory is that expressed in
the title of the book on the ``Castle of the Soul.'' The soul's different
stages of development in the direction of absorption in and union with
the Godhead are compared with the different dwellings
in the castle. Seven such dwellings are distinguished, and
peace and deep mood may reign in one room, though
noise and unrest are heard from another.

As allegory is related to symbol, the arbitrary
comparison to the involuntary one, so
is --- when she, so to speak, passes over from psychology
to theology --- the conscious theological interpretation of
the experiences related to the immediate assurance of the presence or revelation of divine
beings in ecstasy. Although she humbly admits that she
is no theologian, she is nevertheless convinced that in some of
her visions she has beheld the mystery of the Trinity ---
and precisely as the Athanasian Creed
sets it forth! Everything stood clear and vivid before her.
Jesus showed her how one God can appear in
three Persons. But she was (unfortunately) not able
afterward to give a closer explanation of what she
had beheld. The higher enlightenment that in such cases
was granted her comes sometimes through
% --- p. 69 ---
\setcounter{footnote}{0}%
\opage{69}words she hears, but sometimes also ``in an incomprehensible
manner, which without words gave her an even
clearer understanding than many learned men together
would have been able to impart to her.'' Thus she saw
in a vision how all things are in God; but she
cannot describe how this was shown to her. ``I
did not see them in their particular shapes, and yet
the sight I had was of an extraordinary clarity.''

When she theologized with full consciousness, she
sometimes drew conclusions from her experiences that scandalized
her confessors. That God is present in all
beings became clear to her in that, in the ``prayer of union,''
she discovered God's presence in her soul.
``Some people who were not learned'' told her that
God becomes present in us only through grace; but
a learned Dominican explained to her that God was in reality
present in all things, which was
a great consolation to her. --- The ``learned Dominican'' Master Eckehart
also held this belief, which with him was closely connected
with very heretical views (see above p.
53), which would surely have given offense to Teresa.

Teresa felt, as one sees, certain of her cause, but
it was nevertheless a need for her to obtain an ecclesiastical confirmation
of her interpretations. By this double relation
to religious authority her relation to her
confessors was especially determined, a relation already touched on above, but
deserving a closer exposition.

For 20 years she had, she says, sought in vain for a
confessor who understood her. A confessor must have
good judgment, personal experience and thorough learning.
There is great danger in being guided by a man
who is ignorant of the doctrine of the Church. Better to give up spiritual
prayer than to go astray from the beginning!
But the confessors who indeed have theological learning, but
% --- p. 70 ---
\setcounter{footnote}{0}%
\opage{70}do not know the inner experiences, are inclined to think
that ecstasies and visions come from the Devil. Teresa
herself was for a time regarded by all, with the exception of her
Jesuit confessor, as a victim of the Devil's
inspirations. The confessor was advised to
distrust her, so that she feared that in the end she
would be unable to obtain any confessor. Confessors must
be humble, must not imagine that they understand what is
after all a closed world to them. Teresa has the
greatest pity for ecstatic souls who suffer under
ignorant confessors who do not understand them. And she
complains that when one reports visions that the Devil
has inserted, the confessors have no doubt, whereas
they are at once hesitant when one reports
visions that have God as their author. (Fondation de Medine
des Champs. chap. 8). She was often threatened with the Inquisition
on account of her visions.

As far as one can see from Teresa's accounts, the
confessors' scruples did not originate from the content
of her visions. She was also too Catholic in her stamp
and disposition from childhood on to engage in
heresies. She says of herself: ``She is so unshakably
bound to the Creed that all conceivable revelations,
even if they showed her the heavens open, would not
shake her faith, not even in the smallest of
the articles the Church teaches. If it happens that the soul
wavers in its faith regarding a doctrine of the Church, or if
it only for a little moment stops at the thought
that when God has said something to it, this may be
just as true as what he has said to the saints, --- then
this wavering or this thought would originate from the Devil\dots{}
When a person does not find this manly
strength in his faith, and when the inwardness of pious moods
and the visions that are granted him do not
% --- p. 71 ---
\setcounter{footnote}{0}%
\opage{71}contribute to strengthening faith, then he must not regard
them as trustworthy\dots{} I see clearly, and I know from
experience, that one ought not to be sure that anything
comes from the Spirit of God, except in so far as it agrees
with Holy Scripture. Were there the slightest deviation,
I would believe that the visions came from the Father of Lies.''
(Autobiography. Chap. 25)\footnote[1]{Teresa sometimes had visions of the Devil, in some cases an ``intellectual vision,'' in other cases an ``imaginative vision''; moreover she had attacks of emptiness, frivolity and indifference, which she was certain were directly aroused by the Devil. Hell also sometimes appeared to her, quite in agreement with the accounts she had read and heard of it. Indeed, she saw in such a vision even the place in Hell that stood prepared for her, if she was not to be saved! (Vie de Sainte Térèse. Chap. 30—31).}.

The scruples must have originated in the fact that she
had at all another access to entering into a relation
to God than through the Church, pastors and the sacraments.
She herself too is sometimes uneasy about the
immediate relation that has been granted her. From an
ecclesiastical standpoint there lies a danger in every such
immediate relation, in which the individual becomes independent
of tradition and worship. This danger the Church sought
to avert by its condemnation of two such eminent
mystics as Eckehart and Suso; but in their case it was
not only the source, but the very content drawn from
the source, that was aimed at. In a mystic of the twentieth
century we shall see the contrast between experience
and tradition emerge very sharply. ---

A doubt about the objective, real character of what
stirred in ecstasy and was made known in the visions
is not to be traced in Teresa, any more than in her
scrupulous confessors. The possibility that the great experiences
should be due merely to her own imagination
is rejected at once, partly because of the resistance that the mind
% --- p. 72 ---
\setcounter{footnote}{0}%
\opage{72}tries to make against surrendering to them, partly finally
because of their effects in bringing the soul
peace and awakening a need for works of love. No
Either---Or between imagination and reality
asserts itself, but rather an Either---Or between two
realities, one that originates from God, and another
that originates from the Devil. The criterion for deciding between
these two origins was for Teresa the effects that the
great experiences led to: joy, courage, aversion to the world,
a feeling of one's own nothingness. Teresa stands here at the same
stage as the early Church (see above p. 42).

\sectionhead{4. \textit{Autobiography and System.}}

Hitherto in this exposition we have been able to use
the autobiography and the ``Castle of the Soul'' side by side. But at
the point we have now reached, a difference asserts itself
between the biographical and the systematic
presentation.

In the autobiography Teresa tells how, after the
high experiences she has described, there arose an unspeakably
painful suffering, from which she cannot free herself, and which
she has great difficulty in describing. She adds:
``It is this state of unspeakable suffering in which
I now find myself. It lies, then, after the period
when I, faithful in prayer, received such great joys and
tokens of grace from the Lord. Now and then he is so gracious
as to lavish such upon me; but my soul's most usual
state is to feel itself crushed by the pain of which I
shall now speak.'' The description that now follows rests
first on a comparison with the pain that
--- because of the soul's infinite longing and slight capacity
--- was felt in the midst of the rapture itself. The difference
is that now the soul feels itself forsaken --- ``like the lonely
sparrow on the roof.'' God seems to be far away. And
% --- p. 73 ---
\setcounter{footnote}{0}%
\opage{73}if a glimpse of glory is granted the soul, it
serves only to make the forsakenness be felt so much the
more strongly. The soul asks itself: Where is now your God?
(Autobiography Chap. 20—21).

This depiction corresponds to what in the ``Castle of the Soul''
is given as characteristic of the sixth dwelling. The sojourn
here means a state of great sufferings. From without comes
misjudgment, and God is distant. And the pain of being
separated from God is the greater, the more glorious the visions were
in former times. What is to be done, then? The
high experiences cannot be forced, and therefore one ought in the
dark times to immerse oneself in contemplations of the humanity of Jesus,
of the Madonna and the saints, --- contemplations
that one must not cast aside because they do not reach up to
the heights of contemplation. We are after all not angels,
in whom love is never extinguished! --- It takes great
courage to endure such a state.

In the ``Castle of the Soul'' this stage, which in the autobiography
is the concluding one, is the introduction to a last stage, the
seventh dwelling, of which nothing is reported in the autobiography,
which indeed is also 15 years older than the more systematic
work.

Two things are said to be characteristic of the sojourn
in the seventh dwelling, namely first a more inward
absorption in the religious ideas, no longer
mere understanding, but spiritual beholding of the mystery,
--- and second a greater continuity between the
different states, which in the sixth dwelling stood in so
painful an opposition to one another; ecstasy and flight of the spirit
are rare, but the spiritual dryness is over;
when God sometimes leaves the soul to itself, it is
only in order that it may so much the more appreciate
what it means to be in communion with God. ---

Of these two peculiarities it is now only the
% --- p. 74 ---
\setcounter{footnote}{0}%
\opage{74}latter that is decisive. For, as mentioned above (p. 68),
Teresa could already at earlier stages --- 15 years
before the composition of the ``Castle of the Soul'' --- grasp in one vision the mystery of the
Trinity. In a letter from the years 1560—1562
(the same years in which the first draft of the autobiography
was drawn up) Teresa communicates to her confessor
an imaginative vision of the Trinity, which agrees
with what she had heard learned men say. And this
vision is evidently the same as that which is related
in the last chapter of the autobiography, and as that whose effects
on the mind are described in the seventh dwelling of the ``Castle of the Soul''\footnote[1]{\textit{Lettres de Sainte Térèse}. I p. 37 ff. (with the Editor's Notes).}. It is therefore not correct when it has been held
that Teresa first in the last period of her life was able
to unite the God of the Bible and of mystical experience.
This capacity she had --- as the whole preceding
exposition shows --- already earlier possessed,
and at none of her stages did she have the slightest
doubt of the truth of the doctrine of the Church, so impregnated
as she was from first to last with the spirit of Catholicism.

The description of the seventh dwelling in the ``Castle of the Soul''
concerns, then, not a stage at which were attained
intuitions that were not earlier supposed to have been possible.
Rather, it reproduces a period in which the great contrasts
in the domain of the inner life had fallen away
and been replaced by a continuity, within which there could indeed
arise rising and falling, but where yet
the image of the lonely sparrow on the roof did not at
any point come to apply. The overwhelming ecstasies no longer occur,
nor the painful
longing, the burning thirst to feel oneself one
with God. A balance prevailed throughout, ---
a balance that the interruptions taught one to appreciate so
% --- p. 75 ---
\setcounter{footnote}{0}%
\opage{75}much the more. This characteristic distinguishes the
last of the states described in the ``Castle of the Soul'' from
the state with which the autobiography ends.

In an interesting discussion of Teresa's mysticism in the
French Philosophical Society Delacroix maintained that ecstasy
does not mark the culmination of Teresa's mystical life,
but that this is to be sought in the seventh dwelling of the ``Castle of the Soul,''
where the contrasts are overcome, and where a calm and
continuous harmony reigns, --- a synthesis or systematization
of the contrasts of earlier periods. ``We see,''
he says\footnote[1]{\textit{Bulletin de la Société Française de Philosophie}. 26. Oktober 1905. (Le développement des états mystiques chez Sainte Térèse).}, ``Saint Teresa reach the seventh
dwelling, this last and definitive state, a synthesis of
the preceding ones. It is the equivalent of ecstasy, in that
it spreads over the whole soul and over the whole of life that
which ecstasy had gathered in one eternal moment, which, so to
speak, was lived high above life. The seventh dwelling is
the apotheosis of life.'' --- A similar conception has been advanced from the Catholic
side by Abbé Pacheu in a book in which
he seeks to unite psychological method and scholastic
theology in his theory of mysticism. The difference between
Pacheu and Delacroix is this, that whereas Delacroix wishes to
carry through a purely psychological method at all points,
Pacheu maintains that where natural causes do not suffice,
an appeal must be made to the intervention of God, so that the
psychological mechanism becomes an instrument on which
God plays, --- and he makes the transition very quickly
from the instrument to the divine organist. But
Pacheu agrees with Delacroix in this, that ecstasy is not
the culmination in the development of Teresa's mysticism; the culmination
is a beholding, which brings peace of mind, in that
the mystical experiences and the theological doctrine melt
% --- p. 76 ---
\setcounter{footnote}{0}%
\opage{76}together in a kind of crystallization\footnote[1]{\textit{L'expérience mystique et l'activité subconsciente}. (1911). p. 76.}. But in order to substantiate
this theory the Catholic author must do violence to
psychology. He asserts, namely, that there is a continuity of memory
between ecstasy and the later reflection
that leads to interpretation; by this ecstasy is supposed to differ from
hypnosis. When one has come to oneself
after ecstasy, the mystic, Pacheu thinks, has a
perfect knowledge of what has taken place\footnote[2]{Après être revenu à soi ou sait parfaitement ce qui c'est passé. \textit{L'expérience mystique et l'activité subconsciente}. p. 142. cfr. p. 176.}. But
that this should always be possible conflicts with the express
testimonies of Paul, Plotinus, Suso and Teresa.
Ecstasy at its culmination excludes every interpretation,
also every proper beholding. As we have seen, Teresa is
referred to involuntary and arbitrary comparisons,
to symbols and allegories, in order to express what she
has experienced. And even the imaginative visions unfold
as a rule successively, so that there is room for
the involuntary influence of ideas won earlier.
Finally, Teresa, like Suso, is very careful
to exclude all non-orthodox elements from the interpretation
that she holds to. A critical labor goes on here
that does not build on an analysis of the experiences themselves, but
is to secure agreement with tradition. ---

To return to Teresa herself, ecstasies undoubtedly stand
for her as the highest states, and
she exults over the glory that in them has been
granted her, little as she can make it clear
to others. With melancholy she recognizes in the autobiography that
their time is past, and that she now lives in the time of longing and
of spiritual oscillations --- corresponding to what she
15 years later calls living in the sixth dwelling. Nevertheless
the depiction of life in the seventh dwelling in the ``Castle of the Soul''
% --- p. 77 ---
\setcounter{footnote}{0}%
\opage{77}is scarcely a pure construction. It rests on the fact
that Teresa has reconciled herself to the fact that the time of the great experiences,
of the mighty visions, is past, and that she now
appreciates the harmony and continuity that her
life has gained in her old age. Whether there is, as
Delacroix maintains, an equivalence between the
mighty momentary movement of the ecstasies and the later state's
subdued harmony, is a very difficult question
to decide. It is at bottom the same problem as
the question whether in the form of disposition there can come full
compensation for what was experienced in the
form of emotion. In my Psykologi I have discussed this problem.
But as far as Teresa is concerned, she maintains no
equivalence here. She makes no comparison
between the stage at which she wrote the autobiography,
and that at which she composed the ``Castle of the Soul.'' Otherwise
it could not have escaped her that she had already
in what corresponded to the fifth dwelling been able to
behold the theological mysteries in full clarity, so that it
was no peculiarity of ``the seventh dwelling.''
She says herself several times that when practical duties
had interrupted her exposition, she does not read through
what precedes before she continues.

In a letter from May 1884 (the year before her death) she
gives an account of her spiritual state. This description
corresponds to the depiction of ``the seventh dwelling,'' which she
composed four years before. There is, however, a certain melancholy present,
and the contrast to the preceding periods is stressed
much more strongly than in the ``Castle of the Soul.'' She feels herself
now secure and has no longer any need of strong bodily
chastisements. Imaginative visions have ceased, while
intellectual visions can still come; the Trinity
dwells constantly in her soul. Such flaming
movements as she could formerly be put into
% --- p. 78 ---
\setcounter{footnote}{0}%
\opage{78}did not now arise, but neither such strong pains. It
is a calm and reconciled mood of old age that reigns here,
--- but no ``higher unity,'' no equivalence in
relation to what the time of the great experiences had brought.
And that it is ``the seventh dwelling'' that she now believes she
lives in appears from a letter of November of the same year. She asks
here a woman friend to read the chapter on ``the seventh dwelling''
to Father Rodrigo and tell him that ``the one he knows
well'' has reached the state described and enjoys the
perfect repose and peace of which the said
section speaks\footnote[1]{\textit{Lettres de Sainte-Térèse}. III p. 367.}. ---

Alas, old age may bring its peace and its resignation;
but it does not always bring, indeed very
seldom, full compensation for what in its time stirred
mightily in us. Even if with this more vigorous life
there came also unrest, doubt and tension, yet even
this was a testimony to the force of life. Where one believes one has
obtained full equivalence, not to say attained a
higher unity, it often rests on an illusion. --- Yet
Saint Teresa had a happier old age than
the two Danish pastors who in their old age confessed
their poverty to one another\footnote[2]{See below my little study \textit{To Præster}.}.

The time when life stirs vigorously, and when the
real work is done, shall have its due, and keep
it. ---

If one compares Teresa with the mystics of the Middle Ages,
one gets the impression of greater psychological ability
in the description of the mystical experiences in their successive
development, more purely historical in the autobiography,
more systematic and influenced by the calm of old age in
other writings (especially the ``Castle of the Soul''). Thereby it is for an
% --- p. 79 ---
\setcounter{footnote}{0}%
\opage{79}investigating consideration easier to discover the relation
between experience and interpretation. In her own case
this relation has naturally not come to clear consciousness;
for that she is from the outset too much lived into the ecclesiastical
tradition, although her independence shows itself in
the critical attitude toward the confessors. Even in
the most indeterminate, inarticulate and unusual states
immediate recognition and confident
explanation appear. With modern mystics it was to be presumed
that on this point it would be otherwise, and there
is a most interesting example that it is also
the case.

% --- p. 80 ---
\setcounter{footnote}{0}%

\opage{80}\chapterhead{IV.}{A MODERN MYSTIC}

\sectionhead{1. \textit{Psychology and Mysticism.}}

In December 1910 the distinguished psychologist Théodore Flournoy of Geneva,
whom we have already had occasion to mention above (p. 32f),
was visited by the headmistress of an educational institution for young girls,
conducted in an orthodox Protestant spirit, Mlle Cécile Vé, as he
calls her in his account\footnote[1]{\textit{Une Mystique Moderne}. Dcouments pour la Psychologie religieuse. (Archives de Psychologie. Ed. Flournoy et Claparède. Tome XV. Mai 1915.) --- It is Cécile's own experiences and communications that are presented here, for the most part from the time before she knew that Flournoy was thinking of publishing them. In what follows I have attempted, on the basis of them, a characterization of her remarkable experiences, and I compare it later with the conception Flournoy has put forward.}. She came to ask
his advice about her whole mental condition. This was
a peculiar kind of double consciousness. She felt herself divided between
two selves. The one was sensual and passionate and expressed
itself in strong erotic impulses, which took on a demonic
character and made her nights restless and painful.
The other appeared in her striving to work
in her calling as a teacher, inspired by positively religious ideas,
such as prevailed in the circle within
which she worked. Her condition, however, did not fall
under genuine double consciousness, whether the
% --- p. 81 ---
\setcounter{footnote}{0}%
\opage{81}simultaneous or the successive, for she experienced both
states as her own, and her two ``selves'' had
mutual consciousness of each other. The split was not
complete, and just thereby it became so much the
more painful. ``Je suis à la fois A et B, et l'une aussi sincèrement
que l'autre.''

The two ``selves'' point back to the oppositions between
which she had been tossed about in her earlier
life. She was 47 years old when she turned to Flournoy.
At the age of 18 she had suddenly been torn out of
her childlike simplicity about sexual matters through
brutal rape. This episode she could never forget.
The memory-image surfaced again and again with
a mixture of horror, disgust, and voluptuousness. Later
she had felt an inclination for a married man, but had
broken with him. The relationship seems to have had a playful
and frivolous character. But the fact that she had not
obtained satisfaction for her erotic craving, and
not least for her craving to experience the relation of motherhood,
brought on violent crises in her, in which she
thought that demonic forces overwhelmed her. What
she experienced she herself described as ``nocturnal genital
hallucinations'', ``nightmares'', ``half-somnambulisms'',
``possession by a lower personality''. Flournoy rightly says
that if she had lived in earlier times,
she would have been regarded as possessed. The ``half-somnambulism''
from which she suffered came to stand in a
so much more glaring and more repellent light for her
as her religiosity, after her return from a
stay abroad, had settled to rest in the traditions of orthodox
Protestantism. In accordance with
these she worked with earnestness and energy as head and
teacher at an institute for young girls. The opposition
between what stirred in her in the times of ``possession''
% --- p. 82 ---
\setcounter{footnote}{0}%
\opage{82}and what was both her personal faith and
animated her work of teaching became at last
intolerable to her.

It is significant that she went to a psychologist and not
to a pastor. She not only had a psychological sense,
but was also not unacquainted with modern psychology.
She looked, on the whole, from the outset differently upon her
experiences than she would have done if she had lived
in earlier times or had met with other presuppositions.
She was not disappointed in turning to
Flournoy. He not only represented to her that even very
believing people did not escape spiritual trials, and that
even Paul from his own experience knew a ``Satan's angel'',
but he also impressed upon her that sexual sins
are not of such essential importance as sin against justice
and love. And above all he explained to her
that the sexual drive has so deep a ground in the nature
of all that lives that it again and again rises up against the barriers
which social life sets for it, while on the
other hand it is so closely bound up with energy in the conduct of life
that its extirpation would bring not purity
but dullness. Finally he advised her to exact
observation and investigation of her states, since the study
of our restless and unruly states is often the
best means of diminishing their influence. These
instructions he supplemented with a prudent use of hypnosis
to fortify the insight and strengthen the will to apply
it. --- Flournoy's way of taking up this whole
matter recalls to a high degree the means which Spinoza
(Ethica V. 1---20) prescribes against the ``affects'' (remedia
affectuum).
Cécile herself has stated how beneficial the adoption
of the psychological standpoint was for her. It freed
her from the many self-accusations and from
% --- p. 83 ---
\setcounter{footnote}{0}%
``\opage{83}Protestant scruples''.
``My real divine inspiration'', she says, ``came to me the day when I resolved
to regard the whole endless struggle from
a medical standpoint, and when psychological explanations
set me afloat (m'ont mise au large), in that they greatly
lightened my moral responsibility\dots{} The morbid
states I can and ought to fight against, but they do not attack
my moral purity''. ``That I have been set afloat
shines out in my whole life. If in a certain sense I
have become more indulgent toward myself, then I can say that I have also become so
toward others, and that I am able to love them in a better way''. ---

The struggle, however, was not over. But a good year and a half
after she had turned to Flournoy, she received an
inner help which made it possible for her to do without hypnosis.
Whereas the half-somnambulistic experiences which she
called demonic had a tendency to personify
themselves as uncanny beings, there now appeared as a
counterweight also a beginning personification of the opposing
tendencies in her nature. She often had an
indefinite sensation of the presence of an invisible friend.
Everything that went on in her and tormented her she could
confide to him. No visual sensations were connected
with the ``presence'' she perceived. And when
the ``Friend'' spoke, she could not decide whether the words sounded
in her inner self or came to her from outside. One of these
experiences she describes thus: ``I felt that he
came from the other side of the room and approached
my bed, so that I had the sensation that I
could reach him if I stretched out my arm. But the
material sensation played the smallest part in
this revelation (apparition); the prevailing impression
was that of friendly presence, and presence of something
that was mightier than I myself. Throughout my whole life
% --- p. 84 ---
\setcounter{footnote}{0}%
\opage{84}I have felt a strong craving to be supported, to be protected,
indeed, to obey a will that is stronger than my
own. And it is this craving, so often felt, that has now
personified itself''. At times the invisible guest could
remind her of her father, to whom she had looked up greatly
and whom she had greatly missed. Perhaps her
hypnotist and psychological confessor too has contributed
elements to the figure.

When she spoke to the ``Friend'', she did not always receive an answer;
but she could sense whether he approved or disapproved of
what she said. When he answered, the answer came
slowly, automatically, and the content was only what she
had often thought of herself, but was not thinking of at the moment.
A poem which she had learned in her youth and was very fond
of, but of which she could remember only the first
verse, he recited to her almost completely.

The phenomenon had developed gradually, as, according to
Flournoy's experience, it often does in somnambulists.
Cécile could give no definite date for its
first appearance, but she could recognize it when
it came, even before it had assumed a more definite
shape. At times she could sense the guest before
he actually came; it seemed to her then as if
beyond the threshold of her consciousness there were something that
could not press in, because each of her moments
was occupied by her work or its preparation.
She could then try to make room for the hidden
experience (l'expérience latente), much as one tries
to call up something in one's memory that has fled
away. ---

The consciousness of a presence which Cécile here describes
we already know from Teresa. It is what the
Spanish nun called intellectual vision (as opposed
to imaginative vision). It is a phenomenon which is
% --- p. 85 ---
\setcounter{footnote}{0}%
\opage{85}well known, also outside the psychology of religious life.
And likewise there are examples outside this
branch of psychology of the sense of presence passing
over into visual hallucination, as appeared above
in Teresa\footnote[1]{Cf. \emph{William James}: \textit{Varieties of Religious Experience}. Chap. 3.}. Flournoy describes the sense of presence itself
as a hallucination in the domain of the vital sensations
(hallucination cénéstétique), as a rule
connected with organic well-being (euphoria),
calm and rest, and sometimes connected with muscular numbness
(engourdissement musculaire) and with cold legs and
arms, --- symptoms which he has often found in the examination
of somnambulists (mediums).

In Cécile the hallucination of presence did not develop
further, as it did in Teresa. She never came to see
her Friend. At most she had, in addition to the sensations of presence, what she herself
calls an illusion of touch,
but in such a way that she was perfectly clear that nothing
touched her. She adds that nevertheless the conviction
of a presence was beyond doubt; she knew when
the ``Friend'' came, and when he went, and the visit always produced
a feeling of moral help, of calm and of purity.
She is loath to admit that the ``Friend'' should be merely
a product of her imagination, and rests her case
on the fact that she is not master of his coming and
going. But she is clear that it is a deep craving
in her being that here finds expression and satisfaction.
She would wish that she were not so sure
that there was only a doubling of something in her own
personality --- just as other elements in her nature
tended toward personification through sexual hallucinations.
She would, she thinks, take the phenomenon
more seriously if she were not so sure of its
% --- p. 86 ---
\setcounter{footnote}{0}%
\opage{86}mystical origin; now she sees all too clearly the threads
by which the puppet is guided (je vois trop les ficelles).

The most significant effect of the visits of the ``Friend''
is, according to Cécile, that coherence comes into
her divided being. Although his visits are an illusion,
they nevertheless help her forward; she believes she can establish
greater unity in her nature (plus d'unification)\footnote[1]{Cf. with this \emph{Hannah Joël}'s description of how the transition, in her case, from a state of soul-sickness to one of soul-health was experienced as a struggle for unity and coherence. (See my \textit{Mindre Arbejder}. III. p. 81.)}.
This unity could work its way forward in her with
greater ease than if she had suffered from genuine double consciousness.
It was not necessary (as in Morton
Prince's patient B A C) to cultivate a new personality
that could displace both of the two warring selves. The
self that stirred in her in both of the opposed states
was clarified and strengthened with the help of the ``Friend''.

It is probably a reaction against the demonic
side of her nature that works its way forward in
the ``Friend''. It is her ideals from different times
and in different domains that here form a peculiar
complex. Cécile looks, despite the strength with which she is seized
and despite the peace she owes to the phenomenon,
at it in a purely psychological way. The ``Friend'' is identified neither
with Jesus nor with any other Biblical personality,
nor with her deceased father or with her
psychological confessor. She does not believe, as a spiritualist
would, that she is visited by the spirit of one deceased. She does not fear
that the phenomenon, by a purely psychological interpretation,
would lose its value for her.

\sectionhead{2. \textit{A Modern Mystic's Ecstasies.}}

After Cécile for half a year had had her
spiritual abode with the ``Friend'', there fell to her a new
% --- p. 87 ---
\setcounter{footnote}{0}%
\opage{87}kind of experience through ecstatic states. Whereas
in the revelations of the ``Friend'' she had had a sensory
basis in the hallucination of presence, in ecstasy
every such basis dropped away. There was nothing to
describe. In her accounts, therefore, the moments
in which the ecstasy sets in are marked only by dashes.
Whereas Teresa's highest states (see above
p. 64) were hallucinations of presence, Cécile has known
states that were still more indeterminate. Teresa could
immediately recognize what was present to her,
and to the mere vision of presence (the ``intellectual
vision'') there soon joined other visions (``imaginative''),
which made a more detailed description possible. What stood
so enigmatically before Teresa was merely the soul's own state
during the visions; it is for this state that she
looks for symbols and allegories.

Just as for Teresa the vision of presence could pass
continuously into imaginative vision, so Cécile maintains
that there was --- however different the experiences with
the ``Friend'' seemed to be from the ecstatic experiences ---
a continuous transition between them. The transition
took place through the receding of the tendency to personification
which expressed itself in the ``Friend''; it set a
limit which there was a need to transgress. With the ``Friend''
she could still talk; but she sensed a
deeper craving, a craving for something that could make everything
in her fall silent. From this sprang what she called the
divine experience (l'expérienec divine).

If the difference between experience and interpretation made itself
felt already in Cécile's case with the revelations of the ``Friend'',
it had to --- for her intellectual honesty
and psychological eye --- make itself felt still more
in the face of the ecstatic experiences, which left no real
memories at all.

% --- p. 88 ---
\setcounter{footnote}{0}%
\opage{88}In the accounts Cécile sent to her psychological confessor,
the ecstasies appear as lacunae. Only the preceding
and the following states could be described.
Beforehand there may come longing and expectation with involuntary
concentration, palpitation, muscular numbness and
coldness in hands and feet. When the ecstatic state
is about to break forth, a certain resistance may arise
against it; the soul clings, as it were, to its self and will
not surrender itself, not become wholly passive, --- all the while
longing to be gripped by a mighty,
surging wave. What was experienced in the ecstasy itself
stood afterwards before her as an experience of an inwardness,
a depth, an infinity and a simplicity, which set
it in opposition to all other experiences, also to
the half-somnambulism in which the ``Friend'' appeared,
and which now seemed to her to have moved only in
more superficial layers of her personality. The ecstasy
came each time as something original, elementary; only
later could she recognize it as ecstasy, as something
she was not experiencing for the first time.

Only after the end of the ecstasy were reflection, analysis
and judgment possible. And if it is asked how it was possible
to think about it when it had vanished and could not
be recalled, the answer is that after awakening
a flood-wave of wonder and glad enthusiasm went through
the soul. ``I felt'', she says on one occasion, ``that I had been
permeated by the Divine, --- that I had
been beyond the world of the senses and of reflection, beyond
all time, and that I was now myself again, ignorant
of what had taken place, but preserving a kind of radiation
(rayonnement), which made life easier for me.''
It is what I have called the expansion of feeling that
here asserts itself. It asserts itself without interruption in
higher or lower degrees, in that our state of mind is not
% --- p. 89 ---
\setcounter{footnote}{0}%
\opage{89}merely determined by what we experience in the moment, but
also by what we, more or less immediately beforehand,
have experienced, even if we no longer think of it or
can think of it. After ecstasies and visions that cannot
be recalled, it is the only demonstrable after-effect.
It is what Augustine (see above p. 47) called amans
memoria, the only thing that remained after the great
vision. In the last canto of ``Paradise'' Dante says, after
the last great vision has ended, that while the memories of the vision
have vanished, ``the sweet effect'' endures. What
holds of visions holds still more of ecstasy in the narrower
sense.

But Cécile calls the experience ``divine'', and
expérience divine becomes in the following time a technical
expression for her. It is characteristic of her acuteness,
as well as of her honesty, that she
expressly seeks to become clear why she calls
the experience divine. It is a determination of value,
but it is nevertheless a predicate. And with what
right is this predicate used here? She at once makes the remark
that she had conceived the Divine
differently. She was accustomed to imagine it in personal
form, hence as an object. When she now
used the word divine, it was because that which had
asserted itself in her during the ecstatic experience
stood before her as a norm, as that which ought to
be, and which one must unconditionally acknowledge. Or, as
she explains it more closely: ``This experience acts on
me as a force that drives me forward, and which at the same time
keeps me up at a level of thoughts and presuppositions
that is higher than that which I formerly occupied. I notice
its influence even in domains that do not seem
to me to have anything to do with the region in which the experience
unfolds. That is why, I believe, it is
% --- p. 90 ---
\setcounter{footnote}{0}%
divine \opage{90}for me, i.e. permeating my whole being and convincing me
of a reality that stands above that in which my daily life goes on.'' At times
she calls the higher life that enfolds her
in the ecstatic moments straight out God: cette vie-là, c'est
dieu! --- Faust's confession of faith, which indeed likewise appears
as an outcome of an enthusiastic state, fits well
her description, although she does not refer to it. ---
Other expressions too from the circle of her accustomed religious ideas
she finds use for, e.g. grace.
(Ce n'est que maintenant que je comprends le sens du
mot ``Grâce''). ---

It is characteristic that the first and most important predicate
that presents itself to the modern mystic is a
determination of value, an expression of an evaluation of feeling. The mystic
describes here something that we can all know in
other forms and degrees, and which life cannot do without.
Our experiences do not stand on an equal level in personal value.
We distinguish between our best moments, the high points
of our life, and the ordinary course, which takes up
the greater part of life, apart from the moments we call
dry, low, perhaps bestial. And we regard the former
as ``norm'', i.e. as the reality to which we wish to give as
large a place as possible in our life. We find in the
ordinary and low moments --- as soon as we can raise ourselves
above them --- a consolation in having had experiences of the
high kind. Thus the ecstatic moments stand for the mystics.
I disregard here the psychological and physiological
conditions of the high moments, which we too
rightly call a grace to have, because they come
involuntarily, even though we can strive to make ourselves worthy
of having them. These conditions seem more enigmatic
in ecstasy than in the high moments of ordinary people; and yet it is a question whether we know more
% --- p. 91 ---
\setcounter{footnote}{0}%
\opage{91}about these than about those. --- A question in itself is
whether the evaluation that lies in expressions like ``divine''
or ``high'' can be rationally grounded. This question
there is here no reason to enter into\footnote[1]{See on this \textit{Den menneskelige Tanke} § 145—146. --- \textit{Totalitet som Kategori}. (Vid. Selsk. Skrifter 1917). Chap. 5.}. The
evaluation itself, which finds expression in these predicates, is a
fact, and it is, as said, significant that the modern
mystic sees clearly that a predicate of value is the first
predicate that can be used of the strange experiences. ---

In Teresa's attempts to express her ecstatic experiences
there is really nothing that corresponds to Cécile's
explanation of why she calls her experience divine.
She does not know pure determinations of value.
On the other hand there is something kindred between the two female
mystics in their use of analogies to express
what they have experienced. As in Teresa, so in Cécile
too an involuntary comparison springs up.
She too needs symbols. There is an image that
again and again presses itself forward in Cécile, and which is due to
a memory from a stay on the Scottish coast in
her young days. She had there one evening seen a black
plank lying at the water's edge. The waves constantly rocked
it, until at last it was swept along and carried out
to sea, where it was shone through by the splendor of the evening sun.
This memory-image came forth involuntarily in
her immediately after the first ecstasy: ``It seems to
me that I too have been seized by the current and carried
out into the sun, by which I am now wholly shone through''.
Later this image often comes forth after her ecstasies,
and it stands for her as the only thing that to any
degree corresponds to what she believed she had experienced.
Once the image is changed so that a small
% --- p. 92 ---
\setcounter{footnote}{0}%
\opage{92}boat appears instead of the plank, and it also appears
before the onset of an ecstasy, where the resistance of the boat or of the plank
against the waves' attempt to sweep it
along symbolizes the resistance that was involuntarily made
against being overwhelmed by the ecstatic seizure.

In Teresa the image of the iron corresponds to this symbol,
which in a gentle fire lies as a black mass, but in a blaze
becomes glowing (above p. 64), as the plank is shone through
by the evening sun. In the involuntary use of a
kindred image the similarity between the experiences of the two mystics
comes to light, despite all the difference that
otherwise appears between them\footnote[1]{An example of how one and the same image at highly different stages of culture can be used as an expression for the same religious type I have given in my Religionsfilosofi § 36 (the fish on dry land).}. ---

Another circle of symbols presses forward in
Cécile, images which, however, have significance not merely as
analogies, but according to her conviction point to
a real connection. Ecstasy is illuminated by comparison
with sexual self-surrender. What is experienced
during this can no more than the ecstatic
experience be described. Cécile finds the kinship between
the erotic and religion to consist in this, that in both domains
the craving for inward and immediate connection
and complete interpenetration is satisfied. Both are expressions
of life's strong craving, and that the great forces of life
are expressed in the same language, whether it is a matter
of divine or human life, has for her
nothing marvelous. The mystics again and again use the
language of human love to express their high
experiences (Cf. above p. 66f).

It was, however, for Cécile more than figurative speech and
psychological analogy that was here at issue. She had indeed herself
to experience time after time how deeply the
% --- p. 93 ---
\setcounter{footnote}{0}%
\opage{93}sexual instinct is rooted in our nature, and her psychological
confessor had indeed also sought to strengthen and console
her in her struggles with nightmare by reference
to the unavoidable natural drives we carry with us,
even if we rise ever so high in spiritual culture. From this
it follows, as Cécile has clearly seen, that the more deeply an experience,
even of an entirely different kind, takes hold of our mental life,
the more will also the sexual side of our nature
be set in motion. Often she had to experience that ecstatic
experiences incited the sexual drives. What
for reflection came to stand in the most glaring contrast,
thus had a real, subterranean connection.
The very opposition presupposes forces that work in the
same plane. After having again experienced the glaring
opposition between the high and the low, Cécile says:
``J'ai peine à réaliser que je suis la même qu'il y a 24
heures, alors que mon être tout entier vibrait encore de
t'attouchement du divin\dots{} Pourtant c'est bon d'arriver
à voir quelques-unes des ficelles de notre pantin,
comme cette connexion, qui semble si étrange, entre
l'émotion religieuse très pure, très haute, et l'excitation
sexuelle''.

The relation between sexuality and religion is now not
merely this, that the emotion aroused by religious motives,
by way of organic diffusion, propagates itself to the sexual
domain, but also the reverse, that sexual craving
can undergo a metamorphosis, in that the energy it has
released is transferred to psychical processes of an entirely different
kind; the possibility of such a metamorphosis lies in the
moment of self-surrender which is present as soon as the sexual craving
and the sexual relation have more than a momentary and
isolated character. The turning point is there when the object
that is to satisfy the craving stands as more
than a mere means, so that an image can form of
% --- p. 94 ---
\setcounter{footnote}{0}%
\opage{94}it that points beyond the single moment and the isolated
enjoyment. Drive and interest then begin to be
bound to something that has value beyond the momentary
satisfaction, and which perhaps can even demand
the sacrifice of this. It is such a metamorphosis, or
series of metamorphoses, that Plato has Socrates depict
in the ``Symposium''. By this road the erotic can pass over into
religion, especially when the whole process of idealization is supported
by the exalted character of religious ideas and
images that meet the individual through tradition
and worship. The erotic cannot from the first have produced
religion, but the craving that stirs in the former can ---
especially through disappointment --- be transformed into religious craving,
in that self-surrender is now directed toward something more comprehensive,
perhaps an Infinite. Energy can be conserved through the displacements.

As for Cécile, it was ``the craving to
be loved and to become a mother'' that demanded satisfaction,
and she herself thought that if this craving had been satisfied,
she would have found rest within the circle
that had then formed around her. She could
be revolted that such happiness was denied her. But
then precisely the struggle with unsatisfied craving
into which she was led had given her life greater value. ``Would
the human soul'', she asks, ``without suffering sense
the thirst that can torment, but which conditions its royal
dignity? What do they know at bottom, those who have not
suffered?'' It is through suffering that the craving is idealized,
even if dregs remain behind. And here
ecstasy, ``the divine experience'', has brought her
positive help: ``Each time the divine experience
has unfolded in me, the glowing iron has touched
my soul and has purified it.'' Even the fact that she could
experience desireless states, which nevertheless captivated and inspired
% --- p. 95 ---
\setcounter{footnote}{0}%
\opage{95}her, must have been of great effect. And what she experienced
in the ecstasies was for her a continuation of
what she had experienced in the half-somnambulistic states,
in which she sensed the presence of the ``Friend''. In the
first of her ecstasies the two experiences passed directly
into each other. Also in the half-somnambulistic states
she had found that there were other elements in her
nature than the unruly craving.

Under all this the memory-image of her father too has
exerted great influence. He had been an energetic,
profound and independent man, who stood for her as
norm and model, so that it was a craving for her to
feel herself ``her father's daughter'' through self-control
inwardly and independence outwardly. Flournoy has, however,
perhaps laid too great weight on the effect of this ``imago
paternelle'', partly because he is too strongly influenced by
the theories of the Freudian school. I suppose, moreover,
that Flournoy is not right when he holds that the ideal
of a man Cécile saw in her father was the cause of
Cécile's not marrying. There were suitors enough; but none
could have measured up to the paternal ideal.
It is, however, improbable that Cécile with her craving
``to be loved by a man and to have a child in her
arms'' should not have taken up this side of life,
if the right one, i.e. the one she could wish to be
loved by, had really come, and there is not a trace
of it in her notes. But the Eros that could
have filled her did not come. ---

The whole side of life which --- not merely as natural
symbol --- but partly as resistance that had to be
overcome, partly as contrast-effect that could humble,
partly as motive that could be metamorphosed, plays so
great a role in Cécile, seems almost entirely
to be lacking in Teresa. This can perhaps be explained by the fact
% --- p. 96 ---
\setcounter{footnote}{0}%
\opage{96}that she came so early under the discipline of convent life, so
that the natural tendencies were given no time to unfold
themselves. They have led a more subterranean existence, and
when they stirred more strongly, confession in the confessional
has formed a sufficient counterweight. Neither in the autobiography
nor in the writings does one notice them directly.
But in the symbolism they can be traced, for erotic images
presuppose prior experience or consciousness
of what the image means. Only in a single vivid
description does the organic diffusion mentioned above appear
distinctly. The possibility, however, always remains that
Teresa, despite her great capacity for observation and analysis,
has not been able to match Cécile as regards this
capacity. Cécile thus often notes announcements and
omens of the onset of ecstasy, whereas with
Teresa it mostly comes unprepared and unexpected.

The most important difference between the Spanish nun of
the sixteenth century and the modern mystic
of the twentieth century is, however, that the former very quickly,
and without noticing the transition, passes over from experience
to theological interpretation, whereas Cécile for the time being remains
standing at involuntary symbolism.

\sectionhead{3. \textit{Ecstasy and Religion.}}

One might wonder that Cécile, who paid homage to a
pronounced confessional religiosity, and who had undertaken
the religious direction of an educational institution,
did not at once in the powerfully seizing experiences that
fell to her lot see announcements from the God of the Bible and of the
Church or revelations of Christ. As we have
seen, there stirred in her first and foremost a craving
for simple appreciation of what was experienced in its significance
for her inner life, and next she involuntarily sought
for images to make this significance intuitive. An
% --- p. 97 ---
\setcounter{footnote}{0}%
\opage{97}image like that with the plank in the sunlit billowing
sea is a good image for the poetry of life that is an
immediate expression of what has been experienced about the relation
between the dark and the light, what weighs down and
what lifts up within mental life. But any dogmatic
application lies far from her. Whereas Teresa immediately
recognizes Jesus in the figure that appears to
her, first as merely present, later as object
of hearing and vision, Cécile expressly declares that she
has never conceived the ``Friend'' as one with Jesus, and that
she also, when she prayed, had no impression that
it was the Friend she came into relation with. The Biblical
ideas were not at all brought into connection
with that which procured her the best consolation and the
deepest peace. She takes the experience as it comes.
In the same way as toward the Biblical ideas she placed herself
toward the spiritualist interpretations.
It would indeed, she says, have been a quite simple
explanation to assume spirits as origin, and it would
have been a great joy to her if she could have
believed that it was her father's spirit that visited
her. But she had no occasion for it.

The difference between Teresa and Cécile in this respect
is connected with the fact that in Cécile the
involuntary plays a far greater role than in Teresa, and
ecstasy stands from the first in definite opposition to
her conscious religiosity; the significance that ecstasy acquires
for her is perhaps, as we shall see, conditioned
by the fact that her conscious religiosity at bottom did not
satisfy her. In Teresa, by contrast, ecstasy was
prepared through prayer. ``The door to the castle of the soul'' is
prayer, both that in words and that in the inner self, accompanied
by the turning of the mind (l'oraison tant vocale tant mentale,
accompagnée d'attention). The various
% --- p. 98 ---
\setcounter{footnote}{0}%
\opage{98}forms of prayer are so many steps up to ecstasy. After the
voluntary turning through prayer the
highest state then comes involuntarily; ``before we think of
raising our thought to God, the powers of our soul are already
inside the castle, without knowing how they have come in there.'' The
voluntary thus belongs only at the
lower steps; later it is precisely a matter of not doing violence
to oneself; but voluntary preparation nonetheless plays a
role that has no counterpart in Cécile.

Just as Teresa --- in defiance of some of her confessors
--- found in her mystical experiences confirmation
of the thought of God's indwelling in all things, so that it is not
merely through grace but in a still more inward
way that he works in the human soul, so too
Cécile is convinced that ``the divine experience''
does not presuppose a breach of the laws of nature, but is an outcome
of a natural unfolding of the inner life. But
for Cécile this conviction acquired --- at any rate for a period --- a
significance of principle that Teresa could not
attribute to it. The point was that Cécile in the period
in which ``the divine experience'' was of decisive
significance for her could not use the Biblical
and churchly ideas, however much she
might wish to.

``I wished'', she says, ``that I could call forth
my old Christian assurance again in its original
form, take refuge in the thought of the atonement of my sins
and of a redemption procured outside myself;
but I no longer found strength or consolation
therein; it was as if it no longer reached
me\dots{} I have formerly been able to `feel God' in my life
or in my prayer, but it was a God I was familiar
with, whose history I so to speak knew, and whose influence
was more or less as I expected it. It
% --- p. 99 ---
\setcounter{footnote}{0}%
\opage{99}was a God who had stepped forth in history, revealed
through the life and teaching of Jesus. But that with which I have now
come into contact surpasses every expression and
all history, --- and yet it is not merely an abstraction.
\dots{} Now every formulated prayer, or
at any rate every prayer that is formed beforehand to be
performed at a particular moment, seems to me to be sacrilege. I
probably go too far here, but it is now my
conviction, against which I cannot act without
having the feeling of violating my conscience\dots{} In the last analysis it must nevertheless be of little significance
for me where the divine experience has its place
among the realities that are registered under a
particular label. It exists, and its value for my spiritual
life is as beyond doubt as this spiritual life
itself\dots{} For a person who, like me, has been brought up
with piety and theology, it is difficult to
accept a wholly new form of revelation; it does not happen
without wavering. But one must learn not to be afraid of
the truth, whatever it may turn out to be!'' ---

In ``the divine experience'' there lay for Cécile
the germ of a new concept of God. ``Only now'', she says,
``does it really dawn upon me how narrow, dogmatic and
human my concept of the Divine was.
Led by an imperious craving everywhere to find again the character
of something personal, I had evidently come
to forge for myself a God who was wholly contained in the
moral sanction and wholly and entirely revealed in
the Father whom the Christ of the Gospels proclaimed\dots{} I feel
myself on the border of a great, promised, but as yet
unexplored land\dots{} In loftiness and in immediacy
the divine experience surpasses everything I have hitherto
been able to imagine\dots{} I feel such a spiritual
renewal after the brief moments of `divine
% --- p. 100 ---
\setcounter{footnote}{0}%
\opage{100}touch', that it seems to me childish and meaningless
to go on seeking edification by hearing
sermons and attending meetings\dots{} What does the whole
Christian preaching and the immense edifice of dogmas mean in face of
this direct revelation to a human soul?\dots{}
It seems to me now that I stand on the
shore of an immense ocean, and that God passes by in the
mighty waves and in the rays of the setting sun.
But it is not the small-cut and modeled
God. It is the highest Real, the final end
of all, the living and triumphant force, in which
I finally find the fundamental being that is raised above
time and space and above good and evil!'' ---

If one wished to say that Cécile here glides over from description
to interpretation and even reads a whole pantheistic
metaphysics into her experiences, it must nevertheless be answered
that she has really only stated that what she has experienced
does not allow her to let any opposition hold
within that which in her highest, most strengthening moments
has appeared to her. She asserts a docta
ignorantia, which leads her to look with humor at
dogmas and theological edifices of doctrine. ``If this really
is the way in which God reveals himself to
a human soul, how could the theologians then
have written so many thick books about it?'' ``For the future
the thought of God imposes on me above all silence.''
--- She has learned from experience that ``our craving, especially our
spiritual craving, always surpasses the expression we can give
it.''

From what operated in ``the divine experience'',
disharmonies, misfortunes and crimes, as
ordinary experience shows them to us, could not be derived.
God cannot be the cause of evil. She found
consolation in this conviction in face of the disharmonies
% --- p. 101 ---
\setcounter{footnote}{0}%
\opage{101}she had to experience in her own inner self. She had moments
when she rebelled that she was created with such strong
instincts, when they were not to find satisfaction.
But especially when the World War broke out, the thought
became a consolation to her that ``God is not in this world''.\footnote[1]{Probably she has found this sentence in a posthumous work by Mme Flournoy, where it is maintained with great inwardness and emphasis. On this work see \textit{Den store Humor} p. 91.}
``Si je me sentais encore tenue de croire qu'il joue un
rôle dans ce drame, j'arriverais au blasphème.'' ---

The turn in Cécile's religious development that set in
with ``the divine experience'' was conditioned
by the opposition to the preceding period of her
life, in which she had more and more entered upon
a religiosity that had its expression in stiff dogmatic
rubrics. ``The experience'' sprang then, she thinks, from
a deep craving that this dogmatism could not satisfy,
and from the spiritual distress which her
confessional faith had not been able to help her out
of. Never had the life of feeling stirred so strongly and
inwardly in her. The inwardization which her religious
life thereby won radiated over her whole conduct of life.
Although she could not communicate the experience itself
to others, it made it nevertheless possible for her to work with
greater force and love. Since what she experienced was
so highly different from what others experienced, she felt
herself more independent than before and acquired more faith in herself
than before. Now she had at last become a person
(quelqu'un), and not merely a complex of the thoughts of others.
She understood herself better, and therefore also understood
others better. With greater patience and more inward
sympathy she now stood in face of the sufferings of others, even if
they were of another kind than her own. She had hitherto
not had sufficient respect for her own spiritual
% --- p. 102 ---
\setcounter{footnote}{0}%
\opage{102}development when it did not proceed in the accustomed dogmatic
rubrics. She saw that she had no right to maltreat
her soul in order to get it into a definite schema.
Even the most authentic Christianity is --- she now sees
--- not a spiritual medicine for all. A sermon
she hears only increases her confusion. She felt a desire
to become a Catholic in order to be able to kneel in silence in a
church without being disturbed by what people say\footnote[1]{It is surely not a few who, like Cécile, find that the liturgical part of worship can captivate them, while the sermon only disturbs. In a note from 1847 Kierkegaard says: ``Unconditionally, the singing of hymns occupies me most of the whole service.'' William James told me that he did not belong to any confession, but liked to hear hymn-singing; sermons he could not bear to hear. J. N. Madvig says in a posthumous note that organ playing and singing awaken deep and lively emotion, longing, love and courage, --- ``but then comes the sermon with its so often halting and contradictory dogmatics, its obscurity and crooked inferences.'' A deceased friend of mine, a famous natural scientist, who was a constant churchgoer, was greatly absorbed by the hymn-singing, but as a rule behaved predominantly critically toward the sermon.}.

\sectionhead{4. \textit{The Principle of Personality and Love of Humankind.}}

The strange states that our modern mystic had
experienced seemed bound to lead to a decisive change
in her spiritual life, namely to a deepened understanding
of herself and of others. She now felt herself as
``her father's daughter''. Her father was an outstanding pedagogue,
who in religious respect was strongly influenced by
Vinet and in connection with him maintained the right of the individual
personality in face of all confessions and dogmas.
His watchword to his children was: Follow your conscience,
and immerse yourselves in earnest in the life and work of Jesus.
Cécile had for a time denied this counsel. After her
return from abroad, where she had felt herself
% --- p. 103 ---
\setcounter{footnote}{0}%
\opage{103}much tossed about, she had been strongly influenced by the
orthodoxy and pietism prevailing in French-speaking Switzerland
and had with zeal thrown herself into work for
this tendency. She had then not been content with her father's
standpoint, but thought to raise herself above it by headlong
(à corps perdu) throwing herself into ``this whole
somewhat precipitate and superficial, but intoxicating Christian
activity''. Her zeal brought it about that she was appointed
as teacher and headmistress of an educational institution
for young girls conducted in an orthodox
spirit. She found the orthodox basis convenient for exercising spiritual
influence on her pupils. But now she felt
in increasing degree the opposition between her inner life and
the outer forms and formulas. Ecstasy had the significance
of strengthening her in faith in the right of the inner life.
She could now calmly let the whole dogmatics collapse.
Her orthodox surroundings noticed, to be sure, that she
no longer as before was absorbed in the forms and formulas of faith;
but she was allowed to continue her work, since the
heightened energy of the inner life brought new inwardness and
force into it.

But all the while this process of spiritual liberation
was going on, another current too asserted itself
in her. She feared that the intensive inner life
and its experiences would isolate her and gradually
weaken the interest in influencing her pupils. ``In a certain
sense'', she says, ``the life of my own soul is more valuable to me
than the souls that surround me and need my
love''. The ``divine experience'' had now become
the center and sun of her life. She said to herself that
if her spiritual force and wealth had been
somewhat greater, she would be able to let both the inner and
the outer life come to their full right. But here she felt
her poverty. She could not feel satisfied
% --- p. 104 ---
\setcounter{footnote}{0}%
\opage{104}with something she had to keep to herself, and
which could not be communicated to others. She constantly feels
how much she owes to ``the divine experience'',
and she will not throw stones into the spring from which she
has slaked her thirst. She has won new force, and her
religious ideas have been freed from narrowness,
from dogma-faith and theology. But she felt her new wealth,
as if she had received a 100 frcs. note she
could not get changed, although it was precisely small change
she now needed in dealing with others. ---

This conflict is not what interests least
in this modern mystic. From a standpoint in the psychology of religion,
ecstasy and mysticism can be regarded as
strongly marked (in the eyes of many half-pathological) forms
of what in every rich inner life can be experienced
of high moments, bright and warming hours, which stand
in sharp contrast to the scatteredness and externality of daily life.
And the difficulty in which Cécile felt herself placed
is therefore not something peculiar to the mystic, but
can arise for anyone in whom the difference between
an inner and an outer life makes itself felt at all.
At bottom completely direct communication is impossible.
Cécile has rightly said that the fullness of the inner life
can far surpass the means of communication that are at
hand. For those who know Danish philosophical literature,
no closer pointing out is needed that indirect
communication is the only bridge that can be built
between independent and peculiar personalities.
At any rate it is doing violence to personality
to demand the same outer formulas acknowledged and the same
mental development gone through by all. Cécile complains
of the slight respect of religious circles for personality.
One wants to push through one's dogmas, however it may go
with the individual souls. And she herself learned, when
% --- p. 105 ---
\setcounter{footnote}{0}%
\opage{105}she with great caution discussed her problem about
the relation between dogmas and inner experiences with a
pastor (``Pastor Z.''), with what suspicion one
also in Protestant ecclesiastical circles, regards experiences
of the kind that had fallen to her lot. Pastor
Z. established two things, first, that every experience
that unfolds in the border regions of consciousness is of
a lower, fleeting and illusory kind, --- second, that every
experience that does not lead to the assumption of a personal
God is morbid and cannot be made the basis of
the conduct of life. In opposition to this she placed --- in her
mind, for she would not express herself openly to a pastor
--- the experience [Erfaring] that nothing had taken hold so strengtheningly and
purifyingly in her life as ``the divine
experience'', although it could not be brought under the
accepted dogmatic forms (clichés très connus et
catalogués).

It was her own love of humankind and conscientiousness
that caused her problem. She demanded
that the plank she had so often seen rocked in the sunlit
sea should become a lifeboat also for
others than herself. And one cannot help
asking whether she was not in reality on the way
to build such a lifeboat of the plank. She holds
constantly that despite the unutterability of the inner experiences,
they had enabled her ``to spread a little
more spiritual force and divine love around
her''. Can one at bottom get further? Must one necessarily
lower one's demands on one's most inward experiences and thereby
degrade oneself in order to satisfy the demands of
love of humankind? Will not precisely a full and thorough
acknowledgment of the principle of personality in others as
in oneself lead to a truer love of humankind than
% --- p. 106 ---
\setcounter{footnote}{0}%
\opage{106}that which speaks to others the ``catalogued'' language,
but a wholly different language to itself?

Sometimes Cécile had the impression that the moment
might come when she had no right to keep
to herself that which she must be able to communicate,
namely, that it was at all possible ``to come into
contact with the Divine without any intermediary whatever
of speculation, dogma or form of faith.'' But
she dared not draw all the consequences of it. She
glimpses the principle of personality, but she lets it fall
again. The craving for fellowship, and for the security that fellowship
gives, was too strong. ``Shall we get to see the time'',
she asks, ``when true religious freedom will prevail? Sometimes,
e.g. today, I wish it, but at other times
I fear it and discover that I am still very
tradition-bound. I am afraid of going too far. I
let the dove out through the window of the Ark, but hasten
then to shut the window again\footnote[1]{One of the women of the 19th century who has done a great work in the service of love of humankind, and who was at the same time an independent thinker, \textit{Florence Nigthingale}, sent a dove out into the world, and she did not shut the hatch against it. On the contrary she maintained with religious enthusiasm the conviction that it was every human being's task freely to bring the best possible out of his personality. She could not conceive God as an individual being, but as a universal being, a law. Of the churchly customs she used only what she needed. She was ``very sure of God'', but very sparing in defining his nature. (\emph{Edwad Cooke}: \textit{Florence Nigthingale}. Danish trans. 1917. I. p. 349—441; II. p. 188). --- A following essay in this book will lead us to speak again of this remarkable woman. (See ``Features of the History of Prayer'').}. To be sure, there is
cramped room in the Ark; one gets stiff limbs, and one
breathes an overheated air without sufficient oxygen; --- but
then one is nevertheless in the Ark, sure of being in
% --- p. 107 ---
\setcounter{footnote}{0}%
\opage{107}safety, while outside the waves and the strong winds
rage.''

Cécile was clear that if she spoke out completely about
her inner experiences, she would lose her half-clerical
position. It was surely not from fear of such a
break, however, that she did not take the decisive step. Besides
the craving to be something for others, which she thought
was best fulfilled on the basis of the confessional
ideas, there stirred more and more also a
craving for definite images and personifications, a
craving that so much the more easily gained the upper hand as
the time came when the ecstasies ceased. This leads
us over to the last step in that part of her development
into which the documents give us insight.

\sectionhead{5. \textit{Cécile's ``Seventh Mansion''.}}

It was not merely love of humankind, the craving
to be something for others, that led Cécile anew back
to churchly ideas. There stirred more and
more in her a dissatisfaction with the unutterable,
with that which cannot be expressed in any concept,
or rather in any image. Her imagination has, to
use Schleiermacher's expression in one of his letters, been
``determined by a shudder at the obscurity of the indeterminately thought,
not by a shudder at the contradiction that arises
from lending the Infinite the forms of the finite.'' As
I have pointed out (Religionsfilosofi § 95),
there stirs here rather an opposition between
the tendency to hold fast to a certain definite kind of images
and the tendency to let experience and mood involuntarily
call forth the images that can symbolize them.
This latter tendency can be connected with consciousness
of the limitation of every image, and the mystic
will be able to experience states that do not permit any
% --- p. 108 ---
\setcounter{footnote}{0}%
\opage{108}image to emerge, before after-effect and reflection assert themselves.
When Cécile more and more needed a certain
definite circle of images, it was perhaps because she had
been right when she declared that she was not really a mystic.
More correctly it must be said, however, that she did not remain
one. More and more she attempts to translate
the experience in a personifying manner. The experience continued
to be the same, and she does not efface the difference
between experience and interpretation. The personifying interpretation
did not succeed, either, immediately after an experience. ``This morning'',
she reports, ``I tried, from the powerful
religious emotions I had experienced, to derive
a worship, to translate it into formulated prayer, to centralize
it around my love for the person of Christ by
attributing to him an intervention in my inner life. It did not succeed
much. It seemed to me to weaken and
diminish the experience when I wanted to make it one with
that which constitutes or constituted my Christian faith.'' Yet
the translation succeeded better and better under the influence
of a craving of which she became more and more conscious.

I do not think one understands this process of development
in Cécile rightly if one does not notice that
at the same time the ecstasies became rarer and rarer and ceased
finally altogether, just as also her strong dreams
(of which in a single case ecstasy developed)
gradually disappeared. Possibly the connection
of mysticism and sexuality returns here, since she was now at
the age when menstruation, which often seemed to favor
ecstasy, ceased. She herself does not think that age
has had any influence. She thinks that the ecstasies
ceased because they no longer satisfied
her. The relation is rather the reverse: since ``the divine
experience'', which according to her own saying
gave rise to high ideals and comprehensive horizons, for one
% --- p. 109 ---
\setcounter{footnote}{0}%
\opage{109}or another reason fell away, the standard for
religious ideas that it had afforded her no longer made itself
felt. --- It went, as I have tried to show,
in a similar way with Teresa. ---

Added to this was that ecstasy --- in the period when it
was possible --- only occurred now and then, and that it could not
be called forth by resolution, although a certain involuntary
composure and concentration could facilitate its onset.
It was a comet, but there was need of a fixed
star, which one could always find on one's horizon.
A steady faith and trust is perhaps less marvelous,
but is at bottom --- thinks Cécile now --- a better state
than the brilliant but fleeting rays of the experiences.
And she now sets her will to fixing her gaze
on the Pole Star during life's voyage. Thereby
her return to confessional religiosity stands in a
sharper opposition to ``the divine experience'',
than she is herself aware of. The transition presupposes resignation.
The relation is as when a strongly flaring-up
emotion is replaced by a steady disposition. But it
takes place in Cécile through a choice, a voluntary act,
not through involuntary equalization. In the time of
``the divine experience'' she had breathed deeply and
fully, and she will never again come to do so.
But she takes her part and gradually finds calm and
peace thereby. The dreadful World War, which, besides
intervening in her economic circumstances, brought about
concern and anxiety for the friends she had on
both sides of the warring parties, contributed also to her
needing a fixed abode that was common to her and
her surroundings. ---

Cécile's mysticism has thus, as she herself says, been only
a passing phase in her spiritual life. Perhaps
it is so with all mystics. The ecstatic
% --- p. 110 ---
\setcounter{footnote}{0}%
\opage{110}moments have been few, although they have been sufficient
to stamp the contribution the mystics have made to the
development of spiritual life. Concerning the causes of the ecstasies
being tied to certain definite years, hardly
anything can be said on account of lacking specific information.
Cécile has, by her own and Flournoy's opinion,
come out strengthened and enriched from the whole process, even if the
transition to calmer waters after the inner storms
has not occurred without resignation. ---

The psychology of religion stands deeply in debt to Flournoy
for his publication of this remarkable episode in
the life of a gifted woman, as well as for --- what
was a precondition for this document's coming into
being --- the trust that his personality and his research
have been able to awaken. He has sought to see this episode in
its connection with Cécile's whole life and development,
and he has attached illuminating remarks to the account.

It seems to me that he acknowledges too much
the Freudian theory of the erotic-sexual
origin of the ecstasies. As we have seen, and as Flournoy himself emphasizes,
the ecstatic experience precisely stands in sharp
opposition to the state of being sexually seized. The inner
connection between the different forces, which here
perhaps draw nourishment from a common root, is more intricate
than the Freudians in their naive psychology think.
And to this is added that the Freudians have given the
sexual expressions (libido e.g.) a meaning so far extended
that the whole theory is finally washed out. If
Plato too has been a Freudian, then we are at bottom
all of us; but then the special character of the theory drops
away.

A question of more significance in respect of the psychology of religion
is whether Flournoy is right in regarding
% --- p. 111 ---
\setcounter{footnote}{0}%
\opage{111}the development Cécile goes through as harmonious and
progressive. ``She comes back'', he says, ``to
her starting point; but as after an ascent in a spiral
it is a higher stage to which she now returns''. In her
return to the Christian ideas he sees
a symbolic expression of the victory of love and self-surrender
over the spiritual egoism that was called forth
by the self-absorption and self-preoccupation that ecstasy brought about.
(Il y a eu abandon progressif de l'expérience
divine, privilège égoiste et involontaire, pour en venir
au don personnel de soi-même à autrui).

I cannot but see that Flournoy brings too much system
into Cécile's course of development, in that
he above all seeks to demonstrate the conclusion of the development
in a higher unity. He commits here the same
error as Delacroix and Pacheu in regard to Sainte-Térèse.

The spiral I cannot catch sight of. Cécile herself has
undoubtedly found peace of mind, and quite especially harmony
with her surroundings, by her return to her
starting point. But it is an adaptation, which does not in
and of itself denote progress. Yes, if being
a teacher of religion for young ladies were an ideal that
ought to be held fast at any price, it would certainly be a matter
of working oneself out of everything that could hinder
full self-surrender to it. Hopefully there are, however, other
ways for love of humankind, which make it possible to be
something for others without sharing their faith at the price of denial
of one's own deepest, most penetrating experiences.

Cécile has undoubtedly also later preserved much of
what her mystical period has taught her about the
independence of the inner life in relation to all words and all outer
forms. But her personal solution of the problem
of the opposition between inner experience and outer activity
cannot be said to be the one right one. It would
% --- p. 112 ---
\setcounter{footnote}{0}%
\opage{112}in any case be sad if it were so. She
admits herself that she closed the hatch of the Ark,
after she had for a moment been so bold as to let
the dove out. And she admits that it was stuffy inside
the Ark. Has the air been better now that she had
settled down inside it? She has surely been able to contribute
to an improvement of the air on the basis of her
experiences. But would it not have been better to let the hatch
stand open and then patiently wait to see whether the dove
would come back --- perhaps with an olive leaf? A
new principle was dawning here, but tradition was
stronger than the memory of the experience that had let the
principle dawn. Only thereby did it become possible that a wholly
different and opposite principle could be followed: that love of
humankind is to show itself in adopting people's opinions.

% --- p. 113 ---
\setcounter{footnote}{0}%

\opage{113}\conclhead{CONCLUSION}

The aim of this essay is to illuminate
the relation between experience and interpretation by an investigation
of the religious significance of ecstatic states at
different times and under different individual and
social conditions. This is naturally only a single example
of the significance experience can have for the
religious ideas, but it is an important example.

The relation between experience and interpretation raises both
an epistemological and a psychological problem. For
epistemology, when it does not fall back
into dogmatism, all science must in the end stand as a great
interpretation of experiences, each of which is described
as exactly as possible, preferably through number, measure and
weight. Also from the epistemological standpoint
it is a matter of not confusing experience and interpretation. On
interpretation it depends how the single experience
enters as a member into the world of our thoughts. Through
interpretation the single experience first comes to
its full right. But it must again and again be tested whether
interpretation exhausts the whole experience, both its quality
and its quantity, both its successive unfolding
and its typical peculiarity. And therefore
it must be tested whether one can from the interpretation derive again
the experience with its special character\footnote[1]{\textit{Totalitet som Kategori}. (Vid. Selsk. Skr. 1917). p. 43---45.}.
% --- p. 114 ---
\setcounter{footnote}{0}%
\opage{114}Psychologically a more comprehensive concept of interpretation is used.
Psychology investigates not merely the interpretations
that lead to a well-connected structure of thought,
hence to science, not merely the clearly conscious and logically
formed interpretations, but also, and especially, the involuntary
and half-conscious way in which an experience
can be brought into connection with the human being's remaining
circle of ideas. The differences that here make themselves
felt for a comparative view, from totemism
through belief in witches to spiritualism, or from medieval
mysticism to modern experience of a ``divine''
kind, are of great cultural-historical interest. When one
follows the series of examples, there appears an ever
greater difference between experience and interpretation, and an
ever greater independence for experience itself, --- in
analogy with the difference the critical philosophy makes
between observation and its scientific interpretation. ---

As regards our special example, the ecstatic experiences,
they lie at the basis of a significant
conflict between the inner life and the traditions,
under whose influence this inner life is led. In
mysticism personality has for long times had its refuge
against a firmly hammered-out dogmatics, which demanded to
determine what was to be recognized as religious
experience [Erfaring]. It is seen in the Israelite prophets, in the mystics of the Middle Ages
and in particularly marked form in
the modern mystic whose acquaintance we owe to Flournoy.
There goes on here a quiet development right
up to the time when it was proclaimed (here in the North
by E. G. Geiger, Sibbern and Kierkegaard) that in the
religious domain nothing has value that does not spring from
personal experience, and when even the individual
was declared to be nowadays ``the only living thing
in religion.'' It is not least the direct
% --- p. 115 ---
\setcounter{footnote}{0}%
\opage{115}incommunicability of mystical experiences that is of significance here,
in that it makes the individual himself and others aware
that the most important thing in life takes shape peculiarly in
every serious individual, and that everything common is of secondary
rank. Many ironists --- i.e. representatives of
the great irony --- approach in their inner life the
mystics, insofar as they, like these, notice that there
is an incommunicable fullness that unfolds in them. It
holds, as I have sought to show, of Socrates. Also
the great humor is conditioned for a substantial part by
the incommensurability between the inner and the outer\footnote[1]{\textit{Den store Humor}. § 16—17; 19; 25.}.
However one then conceives or explains ecstatic
states, --- whether one calls them pathological or
not, --- whether one regards them as pent-up
sexuality or ascribes to them a nobler origin,
--- it is a fact that there have been human beings,
and not the least among them, who had in them the most important,
the most decisive experiences, and this is, for the psychology of religion,
the main point.

For the mystics, as for the totemists, the
ecstatic experiences denoted the high points of their life. From
these experiences radiates a splendor that determines
the order of rank among all other experiences. It lies,
as I have sought to show, in the nature of feeling that, when
it is aroused by a single experience, it seeks to spread
over the whole mental life and to color all other elements in
this, whether they stand in connection with that event
or not. Inner states that in and of themselves
have nothing to do with one another can thereby acquire
a common stamp, and a whole can arise out of the
series of inner experiences, whether this be concluded
or not\footnote[2]{See, besides my Psykologi, \textit{Religionsfilosofi} § 27.}. It is as determining elements in a
% --- p. 116 ---
\setcounter{footnote}{0}%
\opage{116}total feeling that the ecstatic experiences acquire their significance.
Even if the time of such experiences should be
past, the problem that presented itself to the mystics
will always present itself to us. For there will constantly make itself felt a difference,
perhaps an opposition, between moments
in which we experience what for us acquires the highest
value, and the ordinary, everyday experiences, or even
the low and animal states, and it then depends on
what result the conflict between these different
elements leads to, and what judgment about life is consequently
passed. It is of great importance that the experiences
are allowed to exert their involuntary influence,
and that the interpretation is undertaken with full freedom. The history of
mysticism has taught us that self-control and thinking for oneself
are not thereby excluded; they will only be involuntarily determined
by the element that appears as the most valuable.
There will here, especially for marked personalities,
be much that does not fit into any concluding concept
and also cannot be ``catalogued'' in any dogmatics,
and therefore, as we have a beautiful example
of in Cécile, a poetic expression will be the only possible one.
All poetry of life has here its origin. Cécile's favorite image
of the plank, which is rocked and shone through in the sea glowing with
the evening sun, has its great value not merely
through its poetry, but also in that it is self-produced,
immediately sprung from experience.

% --- p. 117 ---
\setcounter{footnote}{0}%

\opage{117}\essayhead{SOCIOLOGY AND THE PHILOSOPHY OF RELIGION}

\essaysource{LES FORMES ELÉMENTAIRES DE LA VIE RELIGIEUSE \\ Par Émile Durkheim. (Paris. Alcan). 1912.}

In this work by the leader of the French sociological school
a path has been opened, to a higher degree and with greater clarity than in
any earlier work from that school, toward an understanding
between sociology and the other parts of philosophy. Through the study of a special form of religious
life the author not only casts light upon the religious
problem itself, but also shows how one can
acknowledge the special sociological standpoint for human
life of the mind and yet maintain the legitimacy of other
standpoints. Sociology has hitherto often, in Durkheim's
own presentations too, laid claim to being the
only science of the mind. Partly directly, partly also, and especially,
indirectly, this claim is given up in the significant
work now before us. True, social life, society
and its institutions and traditions, is maintained to be the source of all human
life of the mind; but through the closer development
of the supporting ideas the critical reader is led, by the author's
own hand, beyond the exclusively sociological
standpoint. And what is significant is that sociology
does not lose but precisely gains thereby. Just as
% --- p. 118 ---
\setcounter{footnote}{0}%
\opage{118}society proves its force precisely by allowing the individual
persons to move with freedom and to develop their
peculiarity each for himself, so too the significance of the science of society
will rise the more the different
disciplines that investigate the life of the mind from different standpoints
assert their independence. Comte himself
in his last period set ethics up as an independent discipline
beside sociology. And in the same way
psychology, epistemology and the philosophy of religion will show
themselves to mark independent points of departure in relation
to sociology, however much they may learn from it.
I know well that Durkheim will not acknowledge this fully;
but I shall try to show that, besides all that we can positively
learn from the work before us, which in its own
special field marks a great advance,
we learn more fully than before to see how sociology partly
in fact makes use of psychological and epistemological
presuppositions, and partly in its consequences leads to
the acknowledgment of the importance of those other disciplines of the mind.
It is from this side that a philosopher naturally
comes to regard a work such as the one before us.
For a sociological specialist it will certainly have a significance
that I am not competent to give expression to.
But in my conviction the work has a significance
that extends far beyond the purely sociological field.
For my own part I have in any case
not only learned sociology from it, but have found enlightenment
and confirmation on other points as well. Ideas and
standpoints at which I had arrived by a psychological route
I have found, through the study of this thorough and masterfully
executed work, both correction and confirmation
of. By way of a characterization and a
critique attached to it I shall try to substantiate this
more closely.

% --- p. 119 ---
\setcounter{footnote}{0}%
\opage{119}\romanhead{I}

While the historian and the ethnographer study what
has taken place at definite times, and what has been
or is custom and usage in definite places, the sociologist, according to Durkheim,
seeks to get to the bottom of things. He wants to find
the forces that are at work in human life today, the
realities that lie behind the events of history and
the customs of peoples. As regards religion, it is
then a matter not merely of describing dogmas and ceremonies,
but of finding the real that lies behind them and gives
them their proper significance, --- the eternal and human
in religion. The sociologist proceeds from the view that religion
is connected with the reality of life, with
a human need that stirs more or less
at all times.

Most theories of the nature and foundation of religion
stand powerless before the fact that religion,
despite all resistance and criticism, constantly maintains itself. It must
then correspond to a definite need, a definite force in human
nature. If it were, as the animistic
theory supposes, merely due to dream-ideas that are
taken for reality, or if it were, as the naturistic
theory supposes, merely due to a false explanation of nature, its
time would long since have passed. But it is not
a need to understand, it is a need to live, to
attain strength and security in the struggle of life, that
lies at the bottom. Durkheim thus declares himself decidedly
against the intellectualist conception of religion.

To find the foundation of religion one must first and
foremost be clear that the religious consciousness is not
itself able to interpret the need and impulse that
work in it and drive it forward. It is inclined to
seek the cause of what moves within it outside
% --- p. 120 ---
\setcounter{footnote}{0}%
\opage{120}itself. If one wants to become properly clear about this, it is expedient
to seek out the simplest, most elementary forms
of religion. One cannot, to be sure, lay down as a general
rule that the more elementary form of religion
explains the more developed, any more than
one can lay down the opposite rule. But the more
elementary forms offer the advantage that the popular
and priestly imagination has not yet allowed a luxuriant
growth of myths and dogmas to arise. In the religious as in the
social respect only the strictly necessary is found.
And in particular no individual differences
make themselves felt in the religious respect. All consciousnesses
move in the same wake. The individual and the
social type coincide. Such an elementary religion
Durkheim finds in the religion of central Australia,
which in recent years has been so excellently described
especially by Spencer and Gillen and by Strehlow. It is the
simplest religion known. It has no proper
concept of God. What makes it a religion is that it
is dominated by the great opposition between the holy and the
profane. And this opposition proves here to stand in
the closest connection with the whole social organization.
Like the religion, the social system too is exceedingly
simple. The members of the clan are bound together neither
by blood relationship nor by common dwellings; it
is worship itself that unites them. They have the same
name and emblem (totem) and perform the same
ceremonies, hold the same conception of what is
holy and what is profane.

It is the opposition between the holy and the profane that
has to be explained. Whence comes the feeling of the
sublime that man already at this stage has
toward certain things? How does the opposition arise between
a higher and a lower in and outside man?
% --- p. 121 ---
\setcounter{footnote}{0}%
\opage{121}How is man led to acknowledge something that
is greater, more comprehensive than himself?

Both the worship of souls and the worship of nature presuppose
this opposition. And the religious system we here confront,
so-called totemism, lies further back
than the worship of souls and the worship of nature.

Totem is a name and an emblem that is common to
the members of the clan and obliges them to mutual help,
to avenge one another, to mourn with one another, and to observe
certain ceremonies and certain prohibitions of marriage. Totem
may be an animal, a plant or an inanimate object.
Man feels himself one with his totem, and through it
with his ancestor, from whom it descends, or who
perhaps has himself changed into it. Just as little as
the Australian distinguishes between the image and the thing, just as
little does he distinguish between his totem and himself or his
ancestor. And he feels himself thereby solidary with
all who have the same totem. Through a peculiar, in part
severe and painful process of initiation the young
come to share in the clan's worship of the totem and thereby receive participation
in the holy. A holy implement (churinga or
tjurunga), which he receives at this initiation, connects
him with his totem and his ancestor, and
there issues from it a strengthening and encouraging force, for which reason
the loss of it is a great misfortune.

The difference between the holy and the profane is maintained through
worship. Negative worship enjoins, through fasting,
wakefulness, pain, solitude, silence, that for initiation
into participation in the holy the former man must
be laid aside, so that new qualities and forces may be attained.
Positive worship brings man under the
influence of new forces that radiate from the holy
objects, where they are contained as in a kind of permanent
reserve. Through worship not only man
% --- p. 122 ---
\setcounter{footnote}{0}%
\opage{122}himself is strengthened, but the totem object itself is nourished thereby. It is
a kind of sacrament or communion, whereby man
at once comes to share in the force of the holy and by
his offering strengthens and increases this force. That man
needs such strengthening comes from his dependence
on the powers of nature. The life of nature undergoes changes,
rises and sinks in force and fertility, and
man needs to have his confidence in nature strengthened,
while he also believes that through the exact
observance of the customs of worship he can himself help nature.

Behind the opposition between the holy and the profane
there lies hidden, according to Durkheim, the opposition between
society and individual. Totem is from the first a
social concept, an expression of man's solidarity with
the clan and of the solidarity within the clan. Not fear,
but confident adherence, and at the high points enthusiastic adherence,
to what can bring salvation and security, a
salvation and security that only life in the clan can bring,
is the foundation upon which the opposition between the
holy and the profane unfolds. The presupposition of
the worship of totemism is that there is a fund of force which
the individuals can contribute to preserve and
to increase. It is then not merely a play of association of ideas and
hallucinations that is at work in totemism and in the powers
and symbols that attach to it. The source of religion
lies manifestly in other regions of mental life than the
purely intellectual. Through self-surrender to the holy
and the observance of the holy customs man feels himself
in relation to something whose reach extends far beyond
the existence of individuals and of single generations.
Totem is in the simplest religion we know
the characteristic symbol of this, and it is the clan's continued existence
and welfare that is the real foundation of this
symbol. Through worship is periodically strengthened and re-created
% --- p. 123 ---
\setcounter{footnote}{0}%
\opage{123}society itself. At the festivals celebrated in common and in the forms handed down
from the fathers, particular interests
and daily wishes and cares fall away.
The individual feels himself gripped by something that is greater than
himself. Through the customs of worship, which often have a
captivating and exciting character, the mind is set in a
motion that ordinary life cannot cause. It
feels itself reborn, bathed in the source to which all life is owed.

The clan's continued existence and welfare are in many ways
bound up with the course of nature and conditioned by its rhythm,
depend on whether the fruit ripens, whether rain comes,
whether the hunt succeeds. And through worship man believes
that he can exert influence upon this; therefore it is of so
great importance that the holy customs be strictly observed.
This faith is an outcome of the strong inner being seized
that participation in the common worship causes. The periodic
recurrence of the great festivals corresponds to the rhythmic
course of nature, something that can be traced right up into the highest religions.
In the intervals between the festivals the members of the clan are scattered
and struggle, singly or in small bands,
for their self-preservation. Only through the totem images
do they then hold fast the memories of the great moments of
being seized and enthusiasm, and they feel the opposition
sharply. In their own inner life the opposition between
the holy and the profane makes itself felt, corresponding to the opposition
between participation in the common life through worship
and the isolated existence in the intervening times. The living
need to attain those great moments again and again
is the foundation of religion. So too in the
higher religions one can see men, among whom there prevails
a more or less conscious doubt about the significance of religious ideas
and customs, take part again and again in
the old worship. The point is that the real
justification of religious customs is not to be sought in the apparent
% --- p. 124 ---
\setcounter{footnote}{0}%
\opage{124}purposes they serve, but in the invisible effect they exert
on the mind, the way in which they attune the whole
inner state.

When society in those great moments asserts itself
as a power in the inner life of individuals, an involuntary
idealization of it takes place. The totem symbol signifies something
that reaches beyond what exists. Naturally
man is not himself conscious what the ultimate
ground is for the opposition between the holy and the profane,
--- what in the end awakens the feelings that overwhelm
him. The cause of these feelings he finds outside himself ---
just as he straightforwardly transfers the qualities of his sense-perceptions
to the objects themselves. ---

Now religion and religious ideas, already
at the most primitive stage, concern not merely human life,
but acquire also cosmological significance. The force
by which man feels himself seized in his worship, and which strengthens
and nourishes him, he involuntarily finds again in the whole of nature.
The opposition between the holy and the profane is extended
to all natural objects that come within range of sight.
It is a kind of contagion that is at work here. All that
has any contact whatever with the holy or
with the profane becomes itself holy or profane. The
primitive man does not know our distinctions between
different kinds of objects. Just as he
finds no contradiction in conceiving himself and his
ancestor as one with the totem object, so he conceives
everything whatever as one, by which he is affected
in the same way. This extraordinary contagiousness
of the religious forces (l'extrême contagiosité des forces
religieuses) is due not merely to associations of ideas. The character of holiness
rests, after all, first and foremost on the strong
feeling of dependence and of being strengthened that is produced
in the mind, and the more strongly this feeling expresses itself, the
% --- p. 125 ---
\setcounter{footnote}{0}%
\opage{125}more will it also spread to everything that, especially
in moments of being seized, comes into contact with man.
Hereby one understands the ease with which the
primitive man identifies the most dissimilar
things, however much they appear to the senses with
different qualities.

Thereby a religion becomes not merely a view of life,
but also a world-view, not merely a psychology,
but also a metaphysics. ---

The religious need is thus one with the social
need of self-preservation, which stirs in all and binds
them into a whole. It now leads involuntarily to faith
in spirits. It is a misunderstanding to think that faith
in spirits should lie at the bottom of religion, which
has, as we have seen, a quite different foundation. But there
is in fact no positive religion without faith in spirits.
In the case of the Australians the concept of
soul appears at once with the stamp of holiness. Through his totem
the individual is connected with the ancestor, is
a fragment of the great force that through the worship of the totem
animates the entire clan. There is only a certain number
of souls in existence, and the souls of newborn children are
direct reincarnations of the souls of forefathers, or they
are due to germs that forefathers have sown. The soul is
nothing other than the incarnation of the totemic principle in a
single individual. --- Dreams and similar phenomena
have given support to faith in souls, but do not contain
the real explanation of its origin. Since the clan
continues its life constantly in the same way through the
new people who are born, it stands as a quite natural
idea that the new souls descend from souls
that have lived before. Faith in immortality has as yet here
no moral significance. Both faith in souls and faith
in mighty spirits, who in the past laid the foundation of
% --- p. 126 ---
\setcounter{footnote}{0}%
\opage{126}all good, have developed out of totemism. Proper gods
the central Australians did not know. There is nevertheless a
development in the direction of an ever greater distance between
the great ones of the past, of whom myths and legends speak, and
the other forefathers. Especially at the festivals common to the whole tribe,
not merely to the single clan, the mythical
beings acquired a more exalted and universal character.
From totemism the development could proceed step by step
toward ideas such as we find in higher
religions. ---

The study of elementary religion, of which central Australia
gives so classic an example, is not merely illuminating
with regard to the properly religious life itself.
It also demonstrates in an instructive way how religion
stands from the first as the only form of spiritual
culture; in particular it contains the origin of science.
And since society is the mother of religion, it is
also the mother of science. The authority of scientific
thoughts derives from their being a work of collective
experience, of the common labor of society. It is the observation
and reflection of countless individuals and generations
that lies at the bottom of what individuals can work out
in clear concepts. Religion in particular has rendered thought
a service by making possible the first idea
of what relationships of kinship between things mean. It has taught men to place together
what the senses separate. To the world of ideas the way
could only be paved by an exertion of intellectual
forces which first became possible in and through society, particularly
through its worship. The first attempts at classification
were made with the division of the tribes into clans
as a model; thereby the foundation of all logic was
given. The origin of the fundamental concepts, the so-called categories,
finds its explanation neither through apriorism
% --- p. 127 ---
\setcounter{footnote}{0}%
\opage{127}nor through empiricism; it is due neither to original
inner dispositions nor to sense-impression from without. It is social
states and relations that contain the explanation,
just as already with classification. The idea of time
is developed thus through the divisions of time that correspond
to the rhythm of festivals and ceremonies. The division of space
is originally connected with the different
religious significances of the different directions and places.
The encampment of the clans within the tribe's common
meeting place also has furnished a scheme. The principle of causation,
specifically the principle that like must be produced
by like, is due to the strict demand of the tradition of worship
that the ceremonies be performed in a quite definite
way; otherwise the desired result could not
be attained. Here, then, were norms according to which all had to
form their ideas. The necessary demands
of collective action were expressed in categorical forms. The authority of
reason is of social origin. But society itself
is a part of nature and has only the advantage over
other regions of nature that the fundamental relations appear
more clearly and lie closer to our lives.
Therefore ideas that are worked out after the pattern
of social things (sur le modéle des choses sociales) have been able to
help in conceiving things of an entirely different
nature.

As regards the future of religion, two considerations
are of importance for Durkheim.

Religion, as it can be studied historically,
has hitherto been essentially a social phenomenon. Religion
and society belong closely together historically. Through society's
effect upon individuals --- within them by
countless unconscious influences, from without by countless commands,
prohibitions, habits and traditions --- there arises in and for
them the great opposition between the holy and the profane,
% --- p. 128 ---
\setcounter{footnote}{0}%
\opage{128}between a higher and a lower, between the ideal and
the sensuous. And through the great common acts at
the festivals there arises the enthusiasm, the high flight, that
is a constant source of religious feeling, and without which
it cannot in the long run subsist. But nowadays it has
become different. We live in a time of transition. Individualism
has spread together with the division of labor. We
can no longer celebrate with a whole heart the festivals that
inspired our fathers. We go each our own way. When
it was characteristic of all religion hitherto, right from its
most elementary forms, that it is common to a
group which gets its vital force nourished and renewed through
common festivals, then it stands as a great riddle
how religion can subsist in the future, and
how a new religion is to be able to form. If this
riddle is to find its solution, we can in any case
form no idea of this solution.
There is, to be sure, a religious idealism in the form
of inner, subjective religiosity, which takes shape differently
in the different individuals. But sociology has only
established and realized facts to deal with,
and it is a question for the future whether religious idealism
can come to acquire such a character of fact
that it can become the object of objective study.
In sociology we must keep to what religion has
been in the past.

Besides the relation between the social and the individual
the relation between religion and science also
is of importance for the prospects of the future. Here it must
be taken as settled that religion has long since lost its intellectual
significance. Its cosmology can no longer be used.
In this respect science is a complete
equivalent, --- science, which has itself developed out of
it. But it is also only in a purely intellectual
% --- p. 129 ---
\setcounter{footnote}{0}%
\opage{129}respect that science is an equivalent here. It cannot
replace religion in so far as this is a force that
keeps man alive spiritually, a ground of courage
and confidence. Life cannot wait till science becomes
finished. Indeed, it never becomes finished. And life will therefore
ever drive thought beyond the limits of science.
It will therefore rather be a transformation than an annihilation
of religion that the future will bring.

\romanhead{II}
Only an imperfect idea has the foregoing presentation
of the contents of the work before us been able to give
of the rich instruction that can be reaped from it. It is
built upon a very special subject, the religion of the central Australian
tribes, and this subject is discussed with
thorough expertise and energetic analysis. I have gone back from
Durkheim's book to some of the source writings,
in particular to Strehlow's reports, and have found full
confirmation of the book's conception. By its concise
presentation it distinguishes itself from other writings on
similar subjects, which weary by their breadth and repetitions.
Frazer's famous and excellent ``Golden Bough'' thus often makes
the reader stop from weariness on account of
the repetitions and the lack of concentration. Moreover
there is in Durkheim an instructive criticism both
of particulars and of hasty generalizations.
But all the same the special subject interests him
only as a basis for an understanding of the general nature of religion.
He believes that in the religion that is the
simplest of all we know he finds the essential in
all religion. And by digging as deep as he does, ---
by carrying the analysis through to the end with the energy that
distinguishes him generally as investigator and author,
he has also, as far as I can see, come into contact
% --- p. 130 ---
\setcounter{footnote}{0}%
\opage{130}with the central thing in religion. It has in any case
been to me of great interest to see how here
from a sociological standpoint a conception has been worked out
that I in my ``Religionsfilosofi'' have reached by a psychological
route. But in order that a closer understanding here
may be possible, there is an account to be settled between
sociology and psychology.

From the great opposition between the holy and the profane,
in which Durkheim finds the essential in all religion,
and with which all religion begins, in that it
is fundamental in relation to all worship, all myths and
all dogmas, he goes back to society, its continued existence
and its welfare. But, as he himself hints, there
is still a question to be asked: why is it precisely
society that, without man becoming conscious of it,
is the source of all that is holy? Durkheim
gives the answer himself: It is because society is the source of
the greatest values for man, --- because it is the
great support in the struggle for subsistence, and not least,
because it is the origin of all that makes man
more than an animal. Durkheim's thought here is analogous to
that which lay at the bottom of his statements in the French
philosophical society on ``Morality as a Fact'' (le fait
moral).\footnote[1]{\textit{Bulletin de la Société Française de Philosophie}. 1906. --- Cf. my book \textit{Den menneskelige Tanke} p. 353---357.} He there traced all morality back to society
and based this on the ground that society is the source, preserver
and transmitter of all civilization, especially of the
goods of the spiritual life. It then depends on whether society
also really fills this place. And here
it turns out that there are presuppositions that lie
deeper than the continued existence of society. When the holy presupposes
the social, the social presupposes in turn the
valuable. Society is, after all, to be the source of values,
% --- p. 131 ---
\setcounter{footnote}{0}%
\opage{131}and only thereby can it lie at the bottom of the quality of holiness.
By this route I come from Durkheim's
definition of religion back to the one I have tried to
carry through in my ``Religionsfilosofi'', namely that all religion
rests on the presupposition of the existence of value.
Only so long and in so far as the existence of society stands as
a part of or means to the existence of value can it acquire
a religious character. --- I quite agree with Durkheim
in his sharp distinction between the ultimate presupposition
for or the ultimate motive of religion and the
way in which men conceive the cause
of what, as a consequence of this presupposition or
this motive, stirs in their minds. The analogy with the qualities of sense
is here completely apt. ---

But from this it follows again that society's opposition to
the individual and the latter's dependence on the former is not an absolute condition
for religion. The significance of community and of association
cannot be depicted more aptly than Durkheim
does it, when he shows how society asserts
itself in the innermost inner life of the individual, without its
being itself conscious of it. But it does not follow from this that all experience
of values and all faith in the existence of value is
bound up with society.

The consciousness of society (la conscience collective) is,
according to Durkheim, a synthesis, unique of its kind, of
single consciousnesses. And the same may be said of society
itself: la société est une synthèse des consciences
humaines\footnote[1]{\textit{Les Formes Élémentaires de la Vie Religieuse}. p. 605; 615.}. But everywhere where in the field of
the sciences of the mind there is talk of synthesis, we stand before a
peculiar antinomy. Synthesis is a gathering together and
presupposes to that extent a manifold of elements
that are gathered together, a manifold that is given before
% --- p. 132 ---
\setcounter{footnote}{0}%
\opage{132}the gathering together. But in the spiritual field
we know the elements only as members of the connection
that is due to the synthesis. Psychologically we thus know
sensations, ideas, feelings and impulses only
as members of the whole life of consciousness: the synthesis itself is the
constant presupposition, not merely a result. And in
a similar way in sociology: we know individuals only
as socially determined and conditioned, as members of a society.
There is here, as indeed in all fields, a
constant tension between continuity and discontinuity,
which we can never get beyond. We have no right
to lay exclusive weight on the one or the
other of the members of this relation of opposition. The relation
between society and individuals is only a special example
of the relation between continuity and discontinuity.
If society is a synthesis of individuals, these must
nevertheless at every point have a certain independence
in relation to it. Without such independence the value of
the life of society could not be experienced by the single individuals
or assert itself in their inner life, which nevertheless
according to Durkheim is the presupposition of religion. The historical
connection into which the individual lives himself is itself
come into being as the result of the cooperation through
the ages of manifold individual
(though constantly socially conditioned) forces.

Durkheim answers to this that individualism is secondary.
Thus, e.g., individual totemism,
where a single person has his own peculiar totem, is derived in relation
to the social. This holds, however, only when one
thinks of a conscious and principled opposition between
the individual and the existing society. But the talk is
here merely of the individual forces in whose interaction
society itself properly consists. These forces
cannot at any point be excluded from consideration.
% --- p. 133 ---
\setcounter{footnote}{0}%
\opage{133}Every new element in the forms and traditions of society,
holy and profane, --- every new idea
or custom, --- must be due to one or more single individuals,
after whom it has been taken up by the others. There
is a constant tension between the individual and the
common. Durkheim finds the study of so elementary
a form of religion as the central Australian especially
favorably situated, because here individual differences
do not make themselves felt. But it is the now existing,
fully and completely developed central Australian religion
that is described in the sources. How can it have come
into being? It must surely have had its history of development
like everything else in the world, and here the sporadic,
the beginning from several different points of departure
may just as well have made itself felt as in
organic growth and in the development of species. That one now
cannot definitely demonstrate the sporadic and the individual
does not justify assuming that it has been lacking
before or should be wholly lacking now. We do not know who
composed our old folk songs, but there must nevertheless
have been single persons who first conceived
them, be it in ever so imperfect a form. Later
they have been repeated, altered and recomposed in many
ways. Individual forces, partly receptive, partly productive,
cannot have been lacking at any point, however
far back we go. The very ``human need''
(besoin humain), to which religion owes its origin and its
continued existence, stirs scarcely in the same degree and in the same
way in all. And when it has led to formation
in worship, myth and legend, this must have happened through
experiences, combinations, displacements and postponements,
thus through a series of psychological processes
that must have taken place in the inner life
of single individuals. Just as the gods nourish themselves on men's
% --- p. 134 ---
\setcounter{footnote}{0}%
\opage{134}prayers and sacrifices, so society nourishes itself on the single
contributions to the consciousness of society (conscience collective).
Such contributions are made, partly receptively, through adherence
to and participation in the worship established once for all,
partly actively, through new experiences, impulses and
wishes, which --- through the need, also noted by Durkheim,
to communicate and spread one's faith ---
seek to acquire social significance. Already the origin of individual
totemism presupposes something of the kind. And when
totemism among some Indian tribes in America appears
more fully formed, in part also nearer dissolution,
than in the elementary form known from the central Australians,
this presupposes a series of changes
that must have begun as individual variations,
where in turn the cause of these variations would
have to be found.

It is not even certain that individual deviations from
common faith could not be discovered already at a
stage like that of the Australians. Frazer mentions thus skeptics
among the Papuans on the island of Kiwai, who did not
share the tribe's faith in ``the land of the dead'', whither the souls
go after death. Likewise he mentions
that within one tribe doubts could be cast on the correctness
of another tribe's faith\footnote[1]{\emph{Frazer}: \textit{The Belief in Immortality and the Worship of the Dead}. I. London 1913. p. 213; 239.}. Had we more and
closer knowledge of the spiritual life of the more primitive peoples,
more examples of individual deviations
and initiatives would probably come to light. They cannot ever
have wholly been lacking.

What thus must not be denied, even in the case of the most
primitive societies, naturally comes more and more to the fore in the course
of development, more in the religions that we --- precisely for this reason --- call
% --- p. 135 ---
\setcounter{footnote}{0}%
\opage{135}the higher than in those we call the lower. But then it is
also not merely justified, but a necessary task
for the science of religion to follow the religious life
into the inner life of the single individualities, where it has its
real existence. It is an important material for the study of religion
that can be obtained from the way in
which especially outstanding personalities, consciously or
unconsciously, have shaped the religious life, even if they
in the end were governed by historically transmitted forms.
It is on such documents that I have chiefly relied
in my philosophy of religion\footnote[1]{In my book \textit{Den menneskelige Tanke} I have briefly stated at p. 385—392 my view of the relation between history of religion and philosophy of religion.}. I cannot but
find it interesting that by this route results have been reached
to which Durkheim's investigations also lead in their
consequences. Because of the constant reciprocal relation
between the individual and the social one could
not expect otherwise than that from the two
points of departure one would come to kindred results.

According to Durkheim psychology presupposes sociology.
Thus the opposition between sensation and understanding,
as well as that between egoism and duty, is to find
its explanation in the socially conditioned opposition
between the holy and the profane. But psychology is not
content with this. It asks how then the opposition
between the profane and the holy itself arises, and
further, how the opposition between the social
and the individual arises. And psychology has here an
ally in biology. There must be an involuntary, instinctive
foundation for this latter opposition. It has been
fateful for sociology that its founder
Auguste Comte eliminated psychology from the hierarchy of the
% --- p. 136 ---
\setcounter{footnote}{0}%
\opage{136}sciences and distributed the psychological phenomena
between biology and sociology. When sociology now
appears on its own account, outside the systematic
connection with other sciences on which Comte laid
weight, it easily comes to overlook the psychological
conditions that biology already leads one to enjoin.
Comte himself was in reality led, as is seen especially from
statements in his letters to Stuart Mill\footnote[1]{\textit{Lettres d'Auguste Comte à Stuart Mill}, Paris 1877. p. 275.}, to his doctrine
of man's social nature through the study of Hume
and Adam Smith. Here is seen sociology's debt to psychology.
It is likewise through psychological analysis of the
new documents that now stood at disposal that Durkheim
has reached the results he presents in his book.
The documents give in and for themselves only external facts. Only
by a psychological route can one get back from the intellectualist
conception of religion to the ``human
need'' in which Durkheim rightly finds the real
foundation of religion. It is, finally, only for a psychological
conception that the distinction, so important for the philosophy of religion,
between experience (religious experience) and interpretation
can be carried through in so consistent a manner as
it has been in ``Les Formes Élémentaires de la Vie Religieuse''.

In several places Durkheim builds directly upon definite
psychological laws. Thus he explains the cosmological
application of totemism by the fact that ``the religious forces,
which are due to ideas and feelings awakened by
the influence of social life through worship, by virtue
of their intensity have a tendency to spread to
all that comes to hand or before the eyes of man,
while those forces prevail.'' ``Contagiousness'' (contagiosité)
Durkheim calls this. It is one with what
% --- p. 137 ---
\setcounter{footnote}{0}%
\opage{137}I in my ``Psykologi'' have called ``the expansion of feeling'',
which was first described by Hume and Beneke. Durkheim
rightly remarks that the explanation of how any
connection whatever can suffice to transfer the character of holiness
to new objects cannot be sought in
the laws of the mutual association of ideas; it
is a more deeply lying law to which one must here go back.
It is again a psychological law on which Durkheim builds,
when he points out a tendency to identify
all things that awaken one and the same feeling.

By paying attention to these psychological laws
Durkheim casts light upon what Comte called the theological
stage of knowledge, and in general upon the nature of
primitive knowledge. He also corrects Levy-Bruhl's
assumption that primitive men should have a wholly
different logic from the civilized, because with them a ``law
of participation'' (loi de participation) took the place of
the ``law of contradiction''. For one understands, e.g., that
an Australian can say of himself that he is a kangaroo,
in that he identifies what affects him
alike: the idea of his totem (the kangaroo) and
the idea of his own person. The identification of
being and signifying extends through the whole history of religion.
It was the point of contention between Luther and
Zwingli in the doctrine of the Lord's Supper, and it is still a point of contention
in the conflict of our own day between dogmatism and
symbolism.

When the difference between sensation and understanding
is also to be traced back to religious-social influence, it must
be remarked that this difference does not have the sharp
character Durkheim attributes to it. Sensation contains
understanding already in germ, in so far as it is the ability
to distinguish and to recognize, and there are many gradually
different intermediate links between this elementary thinking contained in sensation
% --- p. 138 ---
\setcounter{footnote}{0}%
\opage{138}and what in a narrower sense one
calls thinking. Psychology seeks to find a
continuity here, as in the other fields of mental life.
In the discussion of Durkheim's theory in the French
philosophical society Delacroix already pointed
to this\footnote[1]{\textit{Bulletin de la Sotiété Française de Philosophie}. 1913. p. 79.}. --- In particular Durkheim supposes that the presupposition
that like is produced by like, which marks
an essential side of the principle of causation, should be
due to the enjoining that the holy acts must always
be performed in the same way, in order that they may produce
the desired effect. But there has otherwise, in
life, in the intervals between the ritual gatherings
that are supposed to have so decisively determined men's
way of thinking, been occasion to have that presupposition
practically enjoined. When thus a man,
of whom Strehlow\footnote[2]{\textit{Die Arunda und Loritja-Stämme in Zentral-Austrálien}. III, 2. Frankfurt 1911. p. 69.} tells, threshed the plants
he had gathered with a stick, cleaned the seed by
sifting it, shook it, roasted it and crushed it on stones
and then ate it, --- here there is a whole series of
actions that aimed at a definite end, and that
had to be done in a definite way in order that this end
could be attained. There is other practice than the religious,
and worship is not the only occasion to learn that
one must go quite definite ways to reach definite
ends. Life gives manifold occasion to
stand before the necessity that first meets man
in the form of the fixed relation between means
and end.

Durkheim makes a peculiar attempt to derive
the basic forms of human thought (the categories)
from religious-social influence. Naturally the development
% --- p. 139 ---
\setcounter{footnote}{0}%
\opage{139}of thought goes on under the constant influence of the
discussion that is possible only in intercourse and communal life.
And the forms and rubrics of social life, in particular the division
into tribes and clans, have quite naturally furnished
rubrics that were involuntarily applied when it was a matter
of ordering and classifying observations and experiences. By
way of analogy (sur le modèle des choses humaines)
the social rubrics were transferred to all things that one
came to deal with in experience. As little here as at so
many other points is a distinction made from the first between
analogy and identity. But this implies nothing other and nothing
more than that social phenomena were the first objects
to which comparison and division were applied
with special interest. And the development that the categories
have undergone, after that first analogy ceased
to be decisive, is far more significant
than that which took place in the period of clan and totem ideas.
The history of religion can, to be sure, yield an
interesting contribution to the history of the categories, but a contribution
that cannot in value measure up to that which can be
yielded by the history of the sciences. It is through the development
of the sciences in interaction with experience
that the categories gradually receive their rational formulation,
their justification and their limitation.

Even if the derivation of the categories from religious-social
influences were historically correct, it would nevertheless not
furnish any justification of their validity. Validity
and authority by no means coincide. The first examples
of the application of a form of thought have no advantage
over the later. The decisive thing is the law that prevails
in the whole series of examples. Religious thought
must naturally present the same basic type as the
scientific, since both are human thoughts.
That the former to some extent precedes the latter gives it no
% --- p. 140 ---
\setcounter{footnote}{0}%
\opage{140}special interest beyond the historical interest. There
is an analogy between social and physical relations; but
in this very analogy the peculiar character of human
thought comes to light. For full
consciousness of the essence of the categories a long
series of examples from the most various fields is needed.
Here there is constantly new work to be done.

As to the future of religion, or rather of the religious problem,
finally, I see it in a similar way
to Durkheim. I am convinced that religion,
as we know it through history, is connected
with an essential element, a deep need in human
nature, and that it therefore cannot disappear
without an equivalent, if the energy and concentration of the life of the mind
are not to be weakened. The spiritual division of labor,
the distinction between science, art, morality and
faith, has made religion a problem. The difference between
Durkheim and me is especially this, that Durkheim expects
the new spiritual form of life from vigorous social associations
and the enthusiasm they would be able to
arouse, whereas I have the quiet hope that already
now in the individual ways of living and views of life there are
germs present that will one day be able to develop and take
shape as models, represented by mighty
personalities, such as the legends portray at
the founding of the higher popular religions.

Already 70 years ago a Swedish thinker,
Erik Gustaf Geijer, said that the individual was now the only
living thing in religion. If that is true, the possibility
of a concentrated life of the mind in the future ---
whether one will call it religion or not ---
must also rest on the principle of personality. One cannot get
around this. But out of it, through the purgatory
which every truth that has significance for personal
% --- p. 141 ---
\setcounter{footnote}{0}%
\opage{141}life must pass through of unrest, doubt and seeking, there will
be able to emerge new formations of the highest value. Whoever
believes in ``the existence of value'' must in any case hope
for this.

I part from Durkheim's book with the feeling of how
much I have learned from it, both where my thoughts
approach his, and where I have had to take
distance. In my conviction it marks a great
step forward, perhaps a turning point in research
in the philosophy of religion.

% --- p. 142 ---
\setcounter{footnote}{0}%

\opage{142}\essayhead{ON CONSOLATION}

\romanhead{I}

To go through what has consoled people at different times
would be to write the whole history of religion, for consolation is always an element in religion,
though an element that can occupy a very different
place in relation to the other elements of the individual religion.
And it is precisely this difference between the element of consolation
and the other elements of religion that is of great importance for the
characterization of a particular, given religion. In general
it may be said that the element of consolation has, in the course
of the history of religion, become more and more predominant.
More and more human life is absorbed
in the work on material and spiritual culture, and it is
the progress or decline of this work, the limits
it meets, or the new possibilities that are discovered for
it, that determine the character and direction
of human feelings and wills. Only when the mind feels its impotence
and must let its hands sink does religion come into force
and can then provide the satisfaction or refreshment that
was sought in vain in the work itself. With the work itself
religion has nothing to do directly; its tasks are set
from a quite different region from the one in which religion is
at home.

In the classical periods of religion it is otherwise. In them
% --- p. 143 ---
\setcounter{footnote}{0}%
\opage{143}the element of consolation does not come to the fore, but religion
is a positive factor in life, itself sets the decisive
tasks, in the end the only task that is acknowledged. To
the classical periods belong first and foremost the great periods of beginning
(as we know them in the higher religions).
In original Buddhism and in primitive Christianity
religion is the sum and substance of all spiritual culture,
gathers everything about a single goal. To the classical periods
belong also the great periods of organization, when religion
is indeed not the only form of culture, but when it draws
all spiritual culture in under itself, in that it stands
as the highest measure of value for all human striving.

In the period of beginning, consolation is given in the very adherence
to the one thing needful. There is only one
need in question, and if it can be satisfied, all is as it should
be. In the periods of organization there are several values to
be safeguarded, and these will never be wholly conditioned by religion
alone. There are many needs besides
what one still holds oneself entitled to call the one thing
needful. Medieval Christianity (in the 13th
century) is the greatest example of such a
period of organization. But already science and
political life, indeed the external ordering of church life itself, emerged as
tendencies that were subject to their own laws and had
their own destinies. Modern Christianity cannot
be called any form of organization, although in the Romanticism of the 19th
century one tried to conceive it so. It
acknowledges, more or less, the cleavage between the
different spheres of life. It is then not so easy to find
consolation; if one is consoled with regard to what is essential,
there is yet so much that was also a matter of the heart, and
that continues to stand undecided and unfinished. In a correspondence
between two Danish pastors around the year 1900 this
emerges clearly; they long for the return of Christ, but
% --- p. 144 ---
\setcounter{footnote}{0}%
\opage{144}if the great event now soon occurs, what
becomes then of the ecclesiastical and popular development
from which a glorious fruit of human life was to issue?
(See in what follows the essay ``Two Pastors'').

What thus appears in the history of religion
holds also of human life generally. There is a consolation
that is so inseparably bound up with what is life's
main interest and main activity that one
can really no longer speak of consolation, so long as this interest
can be fulfilled and this activity continued.
It will then depend on how inwardly all other interests
and strivings that stir in individuals
can become members or elements of the main striving. The
more inward the relation is, the less will the individual partial interests
require a special consolation. Consolation then signifies
everywhere a rest, a refreshment, that is a condition
for continued work. The necessity of consolation signifies
the necessity of gathering and discharging force
before one goes further. Consolation has then, one may say, an
active, not a passive character. What is received is
the basis for new self-activity. Consolation here stands only
as an end, because it can become a necessary means.
It cannot become the last end, only a member in
life's great rhythm. Even the most energetic activity of life
cannot run its course along one single straight line,
but will be like a curve now rising, now sinking.

By a strange contradiction, passive consolation nevertheless stands,
even in very energetic views of life and ways of living, as the highest, the definitive, the goal toward which
one looks. In the Beatitudes in the Gospels this tendency
prevails, more pronounced in Luke than in
Matthew. ``The Kingdom of Heaven'' brings eternal consolation
after the sufferings of the present time. In the popular
% --- p. 145 ---
\setcounter{footnote}{0}%
\opage{145}religiosity, where the expectation of the great breakthrough
does not determine the whole of life, and where the
Last Judgment, which brings the decision, has been forgotten, ``eternal
rest'' has become the current expression for the life beyond.
There the thought of most people stops. Even for Kierkegaard,
that man of struggle, it stood as the highest that the
whole struggle had vanished at a stroke. And in more genial
forms of belief in a hereafter people picture the other
life as a continuation of the present life's holidays
and meetings of friends. Some of the greatest spirits of modern times
have spoken out against this kind of consolation. Already Leibniz
said that if blessedness did not consist in progress,
the blessed would become dull and torpid. (Nisi beatitudo
in progressu consisteret, stuperent beati). Kant's and
Goethe's belief in immortality went in the same direction. On
such a conception consolation can never become definitive.
For where activity continues, new resistance is continually possible
(to be active is indeed precisely to overcome resistance), and new
destinies present themselves according to the relation between force and
resistance. If what matters is absolute consolation and rest, then
Buddhism, with its Nirvana and its striving for
liberation from the world of deeds (Karma), is the most consistent.

But if it is so that true consolation is obtained only when
activity can be discharged, so that only he who can be
self-active can be consoled, then there lies here the possibility of a
tragic disproportion. For then there is no true consolation for
those for whom work is an impossibility. The infirm,
in whom the memory of the time of activity has not
wholly vanished, can then be receptive only to purely
passive consolation. Or he must in any case content himself
with a more elementary and meager activity,
lead his life on a limited stage, perhaps in glaring
contrast to the great arena of his earlier life. Such
% --- p. 146 ---
\setcounter{footnote}{0}%
\opage{146}pain and disharmony often arise when life is checked
without yet being wholly dissolved. Pain is indeed everywhere
a sign that there is still vital force present, though not
enough for normal functioning. In the face of complete infirmity
the individual does not need consolation, since he no
longer has the strength for sorrow and want; there it is the others
who need consolation, especially because they can bring
no help.

Even those who hold that consolation keeps, and ought to keep,
the last word can yet insist that one is not
to be consoled in every case, and that consolation may contain
danger. The author of De imitatione Christi (On the
Imitation of Christ) distinguishes between external consolation, which e.g.
is obtained through conversation with friends, through travel or through
music, and internal consolation, which consists in thinking about
heavenly things (speculari coelestia) and in inward jubilation
(jubilare interius). In his opinion one reaches
far too quickly for external consolation. ``When a man begins
to grow dull, he fears a little exertion and
gladly reaches for an external consolation.'' Internal consolation thus presupposes
work. But the old mystic goes one
step further. ``It is'', he says, ``very great to be able to
do without both human and divine consolation and
for the sake of God's honor willingly to put up with spiritual exile
(exilium cordis).'' This is surely to be understood in the sense
that the ``exile'' afterwards proves to be
``to the honor of God'', not in the sense that he who does without all
consolation should at the same time be clear that it is ``to the
honor of God''; for in that case he does not do without all
consolation, is not placed in full ``exile''. And even where
man is not in exile but feels the divine
love, this feeling must not be made a
mere means to his own consolation. One is to love Jesus for
his own sake, not for the sake of one's own consolation. The
% --- p. 147 ---
\setcounter{footnote}{0}%
\opage{147}mystic here indicates the possibility of a displacement in a
sentimental and egoistic direction, which can be precisely favored
when the need for passive consolation gains power, and which
is characteristic of many forms of modern religiosity.
Often one expressly acknowledges that one accepts
religious assumptions only because one cannot do without
the consolation they contain. The thought of the Middle Ages
is here more grand. The Divine puts
its demands without regard to the individual's private well-being
and sympathy. Thomas Aquinas could thus
very well imagine that a single human being might become
blessed while all others were damned, and that his blessedness
would not be darkened by their fate. For blessedness consists
in ``having the fullness of one's perfection
in God'', --- and in face of this fullness every thought
of finite beings is impossible\footnote[1]{See my \textit{Religionsfilosofi}. Note 105.}.

Ethically considered, the task will be to hold fast to the active
type as the ideal, and to see to it that passive
consolation is not needed. The will is to be trained, as I have expressed
it in my Etik (X, 3), to stand on its own feet.

\romanhead{II}
The boy in ``Flugten til Amerika'' had had a series
of sorrows, but he ``drowned his sorrow and found his
consolation at the bottom of the sago soup'', and Jeppe thought himself, in
face of his great misery, entitled to ``use
the means nature gives us to drive away sorrow''.
But when consolation is not to have this passive character,
but is to consist in the discharge of new activity,
then it presupposes a spiritual growth. There must be
possibility of a new release of force, now a more voluntary,
now a more involuntary one. Perhaps with full
% --- p. 148 ---
\setcounter{footnote}{0}%
\opage{148}consciousness a work may be chosen that has in itself
nothing to do with the sorrow and its object. Consolation
then becomes, as with the boy and with Jeppe, consolation on the side.
The sorrow is diverted by a concentration of energy in
an entirely different direction, and the expansion of sorrow in the mind
is checked or hindered. It may be travel or the
tasks and duties of daily life that claim
interest and energy. Or consolation is found in art, science
or religion. The experience is sought to be displaced by new
experiences of an outer or inner kind. One goes around
the wounded spot. But it is also possible that consolation
is won by directing reflection precisely at the loss
and the sorrow themselves. Then one makes the value of what is lost
vivid to oneself, far from wishing to diminish it by self-surrender
to other values\footnote[1]{\emph{Seneca} exhorts Marcia, whom he wishes to console over the death of her son, to lend an ear to the renown and remembrance of her son. (Apertos aures præbe ad nomen memoriamque filii tui. \textit{Consolatio ad Marciam}. V, 3. (Dialogi ed. Gertz p. 170).}. Sorrow is deepened and pervades
the whole mind, comes to run its line out to the end. It is not
that one wants to enjoy sorrow (what Petrarch called dolendi
voluptas, and Amiel called douceur mélancolique),
nor does one find an honor in being especially unhappy
beyond others\footnote[2]{\emph{Seneca} warns Marcia against the reprehensible honor of wishing to stand as the most unhappy of all. (perversissima gloria, infelicissima videri). \textit{Consolatio ad Marciam}. V,4. (ed. Gertz p. 170).}. Comparison with others need
play no part at all, whether to give
one's own sorrow a glorifying background, or to
impress that others have had greater sorrows, or to
see sorrow as a common fate. What prevails is a need
to immerse oneself in what has happened. Only thereby
can perhaps the petrifaction that is the first effect of loss
be dissolved. Full discharge for the state is sought,
% --- p. 149 ---
\setcounter{footnote}{0}%
\opage{149}a way is opened for the diffusion, the being seized of the whole
organism, that may be necessary. In the Edda poem
``The First Song of Gudrun'' it is described in a moving
manner how noble women, ``the noble wives of earls'',
sat by Gudrun and sought to console her with
telling of their sorrows. ``Each of them told her own
sorrow, --- the bitterest she had borne\dots{}
Yet Gudrun could not weep, so weary was she in spirit
over the dead husband, so full of pain at Sigurd's
funeral procession.'' ``Gudrun sat over the dead Sigurd. She wept
not like other women, but she was near to bursting
with grief.'' But then ``Gjuki's daughter'' reproaches the
other women that they did not know how ``to console the young wife'';
she uncovers Sigurd's corpse and lays Gudrun's head
against his cheek. And then Gudrun wept.

Not only in the face of a loss, but also in the face of what one
never possessed but inwardly wished to win, the
diffusion can be a need. This is what Goethe depicts
as ``Trost in Tränen'' and ``Wonne der Wehmuth''.

If both the value of what is lost and the reality of the loss itself
have been fully lived through, some natures may
feel a need to think the two together and at the same time to think
them into life's great whole. What one has
wished for and striven for, what one has suffered of loss or
disappointment, is then seen as belonging as natural, perhaps
necessary, members of the great whole of existence. The total image,
the ``synthetic intuition'', that thereby becomes
possible can take hold of the mind and fill it with the satisfaction
that spiritual activity brings where
it succeeds. Spinoza and Goethe went, each in his
own way, this road. Life is here converted into an understanding
that will let both the value and its loss come into
their own by making them elements in a whole.
Spinoza emphasizes especially the joy that
% --- p. 150 ---
\setcounter{footnote}{0}%
\opage{150}understanding itself gives, and this joy of knowledge stands with him
as the equivalent for the sorrow that fades away through
consolation. Such an intellectual or aesthetic consolation, as
Spinoza has depicted it with genius and quiet enthusiasm\footnote[1]{Cf. \textit{Spinoza's Ethica}. Analyse og Karakteristik. (Videnskabernes Selskabs Skrifter 1918). p. 73—77; 87---89.},
and as Goethe\footnote[2]{Cf. my \textit{Mindre Arbejder}. Third Series. p. 162---164; 166—169.} has given expression to it in his great art,
will nevertheless, although it does not go around sorrow, be consolation
on the side. The world of reflection and of fantasy is after all
not the whole world of the mind. And even if through
ever such profound mental fusions
and organizations a total feeling forms itself as the result of life,
to which both the great sorrows and the great
joys have contributed, there is yet no
certainty that this whole cannot be burst asunder again,
when fate knocks at the gate, without this
knocking (as with Beethoven) being the introduction to a symphony\footnote[3]{\textit{Den store Humor}. p. 104,}.
A definitive consolation can never be established,
no more by the religious than by the philosophical road.

The fact is that the individual who is consoled is always a part
of a greater whole, whose laws do not coincide with
the laws under which his existence is conducted. Of every
kind of consolation, --- the passive and the active, the voluntarily
chosen and the involuntarily grown, the melancholy
or the despairing, the intellectually superior
(through synthetic intuition) or the emotionally superior
(through total feeling), --- it holds that it has an
individual character. That something affords consolation is a purely
subjective phenomenon, a psychological fact, that says
nothing about the objective connection of existence.
Because the individual has come to rest, or sees a way
% --- p. 151 ---
\setcounter{footnote}{0}%
\opage{151}open before him, it is not said that existence has come
to rest, or that a new way has really been opened.

But it is nevertheless a fact that in
this particular place in existence, in the inner life of this individual human being,
harmony has been reached, or the possibility of
a new unfolding. And if full and honest
use has been made of the force at one's disposal, then there has been achieved
in this place what can be achieved, and what
has been achieved has then also its objective significance.
A flower has burst into bloom, whatever its fate
may afterwards be. The world has then been so constituted
that this flower could unfold in this
place and at this time. ---

What, the world being what it is, can console
human beings can be learned only from experience. It is not abstract
theory but the life-experiences of working and suffering human beings
that is in question. And here is a large and
little cultivated field for comparative investigation.
A poet can contribute in this direction more readily than
a philosopher, and in any case the philosopher needs a large
and comprehensive material from his own and others' experience
in order to be able to see clearly here. The last task of science
is not that of finding an objective theory that fits
all the subjects at hand, but that of showing how
different human beings, --- different wanderers
on life's labyrinthine ways, --- must come
to different conceptions of life\footnote[1]{\textit{Den menneskelige Tanke}. p. 304.}. And here belong also
the differences in what has been able to console
human beings. ---

Illusions too can console. And when we see others
console themselves with what would not console us, we are inclined
to think that their consolation rests on illusion.
But here it holds that if the illusion is the fruit of
% --- p. 152 ---
\setcounter{footnote}{0}%
\opage{152}honest work, which under the given conditions could not
lead to any other result, others have no right to intervene,
unless they command abilities and means
to lead the illusorily consoled person to an objectively more important
result\footnote[1]{\textit{Etik}\textsuperscript{4}. p. 149 and 197.}.

On the question whether illusions can console,
Benedetto Croce rightly distinguishes between two cases.
If the illusion is not recognized as illusion and yet
affords full satisfaction, there must be grounds it
rests on, and it constitutes then the part of the real
truth that could be known under the given conditions.
Illusion it can be called only from the standpoint of a more complete
truth. It is otherwise when the illusion is recognized
as illusion by the person himself. It must
then awaken an anxiety and an eagerness that drives him to seek
after the truth, and if it is nevertheless held fast as
consoling, it must call forth unrest and a bad conscience,
until a serious striving after truth arises. Moreover
Croce denies that truth in the end could be
sad. The full truth must surely comprise both the
joyful and the sad and thus stand above both
parts\footnote[2]{B. \emph{Croce}: \textit{Logica}\textsuperscript{2}. Bari 1900. p. 346---347. --- The whole chapter on ``Consolatio philosophiea'' is very interesting.}. --- To this last remark it must however be said
that it will depend on whether a conception of the whole can be formed
that assigns both joy and sorrow
their definite place in existence. We come then in
any case back to Spinoza's and Goethe's conception.
And as before we must add that it also depends
on whether, in the face of life's conflicting experiences, a total feeling
can form itself that lets them all come
into their own. Reality is never present
fully concluded, and Croce's conception is therefore justified only
against a dogmatic pessimism.

% --- p. 153 ---
\setcounter{footnote}{0}%
\opage{153}H. C. Ørsted in his day (in ``Aanden i Naturen'')\footnote[1]{I have previously discussed Ørsted's remark in my \textit{Religionsfilosofi} § 115 and in \textit{Danske Filosofer} p. 56.}
objected to the proposition that our wishes are to be able to
determine what is truth, and demanded that one
should give the honor to truth: ``the full truth always brings
its consolation with it''. In this remark the great investigator's
love of truth and his optimism appear at
one and the same time. But he did not see that the remark can be understood
in two ways. If the weight is laid on its being the
full truth that consoles, then the question may arise
whether we ever reach it. In any case we must,
until we have it, console ourselves with hope. Instead of
having something to console ourselves with, we must then console ourselves with the thought
that continued striving can bring consolation. That is the old
meaning of the word Trøst, --- confidence, faith, reliability,
--- that here appears again in contrast to the
newer meaning of the word, which probably predominantly
leads to the idea of something passive, of a receiving,
a solace. --- If, on the other hand, the weight is laid, not
on the word ``full'', but on the word ``its'', the meaning is
that truth consoles merely by being truth, and this
can again be understood either to mean that the consolation lies in the
inevitability that is revealed to us, or that it is won
through the spiritual work required to find
the truth, that is, in the form of joy of knowledge. Inevitability
leads us back to the consolation of melancholy or despair;
joy of knowledge leads us back to
Spinoza and Goethe. Ørsted has hardly thought of any
definite one of these different interpretations. He has most likely
thought of truth as the disclosure of a reality
in which everything shows itself in harmony, however often disconsolateness
may seem to have the last word.

% --- p. 154 ---
\setcounter{footnote}{0}%
\opage{154}\romanhead{III}
Toward the end of classical antiquity (3rd cent.
B.C.) a peculiar literature of consolation developed, which
continued through the Roman imperial age and found its
conclusion in a work by a man who has been called
``the last Roman'', Boethius (d. 575 A.D.), a high
official under the Ostrogoth Theodoric the Great.
Through Boethius' work on the consolation that philosophy can
bring (De consolatione philosophiæ), a work that was
much read and commented on in the Middle Ages, the influence of
this literature extends into modern times.

The first renowned of these writers of consolation is
the Platonist Crantor (at the end of the third century
B.C.). His work ``on Sorrow'' has unfortunately
been lost, but its content can be traced in later
writers on the same subject, Cicero, Seneca and Plutarch.
Cicero calls it ``the golden book''. Even in
the consolatory addresses of modern preachers one thinks one can
find again the same thoughts, though naturally without the source
being known. What consolation is sought for is old age, exile,
loss of those nearest, poverty. The philosophers,
who in this field were soon outdone by the rhetoricians,
set reflection in motion in order to show that what
had happened was not sheer misfortune, perhaps not misfortune at all.
Many of these reasons strike us as prosaically sensible;
but at a given time and place they may have had
their effect, and often, no doubt, the will to console has made a
greater impression than the reasons that were brought into the field\footnote[1]{Cf. Constant Martha: Les Consolations dans l'Antiquité. (Etudes Morales sur l'Antiquité 1863. p. 135—189.)}.
In his Tusculans (III, 31) Cicero says that in a
great sorrow of his own (the death of his daughter) he had collected all the different
grounds of consolation he knew, since he sought at any price
a remedy for the turbulence of the mind (tumor animi). He
% --- p. 155 ---
\setcounter{footnote}{0}%
\opage{155}enumerates: reassurance or diversion from the sorrow, demonstration
of its imaginary character, reference to the fact that one
has done one's duty. And these different means must each
be applied at the right time and place.

Already the fact that one came to think about other
sides of what had happened, instead of circling around
the misfortune itself, could serve for diversion and prevent the
mind from staring at the irreparable. In his work on
old age Cicero says that the very composition of
this work of consolation (by which he wished to console his friend
Atticus and himself) has been such a pleasure to him
that it has swept aside for him the burdens of old age
and made his old days gentle and agreeable.
A characteristic expression of the fact that thinking about sorrow
can divert it! It is a consolation that is naturally only
possible so long as one still has strength for reflection.
But, says Cicero then, infirmity and weakness can also
befall young people, so that this side of the
matter does not belong especially in a work on old age.
For the rest he points out that old age is
the last act of a play which nature, who is truly
no botcher, has planned and carried through according to
a definite law. ``It could not, after all, have been
otherwise than that something had to be the last, and
just as it goes with the berries of trees, when they have reached
their ripeness in due time, so the last
section of life too had to become wizened, if I may say so, and
ready to fall; this the sensible man must then put up
with in meekness. Otherwise he makes war
on the gods after the manner of the giants; for what else does
this mean than precisely this: to fight against nature
itself\footnote[1]{\emph{Cicero}: \textit{Cato den Ældre. Om Alderdommen}. Gertz' translation p. 19.}.'' This utterance may be said to be typical of
% --- p. 156 ---
\setcounter{footnote}{0}%
\opage{156}ancient consolation. It rests on a naturalistic optimism,
into which, however, for more emotional natures, now
bitterness, now mild and melancholy resignation enters
as incorporated elements\footnote[1]{See my characterization of Marcus Aurelius in \textit{Mindre Arbejder}. III. p. 115.}.

With the Stoics of the imperial age the grounds of consolation take on a
more vigorous character. Appeal is made at one and the same time to
the will as the one essential thing and to the great cosmic
connection within which man with all his
joys and sorrows is only a vanishing point. Seneca
first represents to Marcia, whom he wishes to console for her
young son's death, that it accords best with her noble
and fine character herself to overcome her sorrow and not to
wait for the day when sorrow against her will leaves
her. Sorrow, he adds, becomes so violent only
because in our carelessness we do not expect misfortune, however often
it passes us by to strike others. Precisely that
which has the greatest value will as a rule be most exposed
to destruction. --- These considerations lead
Seneca to develop how vanishing human
fate is in relation to the great universe, within which
instability and change prevail. How can
anyone --- he says in a similar context in another
work (Ad Polybium) --- dare to believe that he and
those nearest him should be able to be exceptions within
the great order of nature, to which everything is subject?\footnote[2]{\emph{Seneca}: \textit{Ad Marciam}. (Dialogi ed. Gertz p. 173—201); \textit{Ad Polybium} p. 324.}

Epictetus stresses more strongly than Seneca that man's
essence is contained in the inner life, in the will. The god
would gladly have given man more than the will, but has
not been able to, any more than he can rob him of this!
But then the glory that man has in his inner life is
all the greater. You regard it as a misfortune
% --- p. 157 ---
\setcounter{footnote}{0}%
\opage{157}that you have not seen Phidias' Zeus at Olympia --- but in your
own will there is something worth beholding, and which
requires no long journey. And in there dwells the god
and your genius, so that you will never be alone, even if
you close the doors of your soul! One is then to be concerned only
with the will, with having mastery over one's ideas.
Everything else is not in your power. You are to love human beings,
but remember that they are mortal. The misfortunes of other human beings,
even of those nearest you, do not depend on your
will. Sigh, when such misfortune comes, but take care
that you do not sigh in your innermost being! You are a part of nature
and must bow to its course. In all sorrow and lamentation
there is expressed really a dogmatic assertion about what
is good and evil. It must be combated, for at the world's
great spectacle there is no use for a discontented
spectator. If you are not satisfied, then go your way. Say that you
are not playing any more!\footnote[1]{\textit{Epiktetus}. ed. Schmidt passim.}

In Boethius Stoic and Neoplatonic thoughts meet in
his remarkable work on ``The Consolation of Philosophy'', of which
an excellent historian of literature\footnote[2]{\textit{Fridrich Leo} in Kultur der Gegenwart. I,8. p. 395.} has said that it reflects
the radiance of antiquity's setting sun as a symbol
of its rising in the future. The book was written in the
prison into which King Theodoric had cast him, and which
he left only to suffer death. There has been dispute
whether he was a Christian or a pagan. From the content of the book
nothing can be concluded as to his Christianity; it gives a
purely human view in the track of Stoicism and Neoplatonism.
Since he is said to be the author of theological
writings, whose spuriousness there does not seem to be ground to assume,
it is now generally held that he was
a Christian, ``brought up in an enlightened and culture-eager Christianity,
which at that time among the aristocracy formed the foundation of the
% --- p. 158 ---
\setcounter{footnote}{0}%
\opage{158}world-view''\footnote[1]{\textit{Eduard Norden}: Kultur der Gegenwart. I,8. p. 412.}. In any case in distress
and mortal danger he was able to console himself without taking refuge
in Christianity. Philosophy visits him in his prison
and drives away the Muses, who are only able to nourish sorrow,
not to heal it. The true consolation, she teaches him, consists
in the knowledge of the universe, which --- as Plato
has shown in ``Timaios'' --- is permeated by the
spirit of the Divinity. Because misfortune comes upon us, the order of the world
has not changed; rising and setting alternate
now as before for all finite beings. But mortals
seek happiness outside, not in their own inner life, where the
highest good, which is one with God himself, has its dwelling. No
dispensation of fate can rob this inner possession, and
he who possesses it needs no external tribunal. We
live in a world of division and imperfection;
our earth is only as a point in relation to the whole
universe, our lifetime only as a moment in relation to
infinite time. But even in this imperfect existence
man can preserve his unity with the Divine\footnote[2]{\emph{Boëthii}: \textit{De Consolatione Philosophiæ}. ed. Obbarius. Jena 1843. (Især II, 3—7; III, 8—10).}.

\romanhead{IV}

It was reflection on life and the world that the ancient
thinkers took refuge in when consolation was to be sought.
They demanded a subordination of the life of feeling to the
life of thought, whose harmony and strength were to be preserved
through the course of fate. When in Boethius Philosophy
drives the Muses away from the old man's bed in
prison, a certain contrast between
ancient thinking and ancient poetry is thereby indicated. In Greek tragedies
mood and passion are often allowed to win full
expression. Think of Cassandra's passionate
% --- p. 159 ---
\setcounter{footnote}{0}%
\opage{159}lament, Antigone's wail at having to leave life
and Medea's despairing outbursts. Here is expressed only the
consolation that can lie in full self-surrender to sorrow.
But there are also examples of mild solace
in sensing the nearness of the Divinity. In the midst of his mortal agony
Hippolytus hears the voice of Artemis, his goddess, and
already that eases his pain:

\begin{quote}
O fragrance of divinity that reaches me!\\
I sense you, even in my suffering,\\
and I feel relief!\footnote[1]{\emph{Euripides}' \textit{Hippolytos}. v. 1391 f. (Gertz' translation).}
\end{quote}

What poetry has put forward in these examples has
surely often stirred in life. But the spokesmen of the age
harbored distrust of the feelings, because they set the imagination
in motion and tear man out of the inner freedom,
man's firm stronghold, and make him dependent on things
by tying the mind to them. Therefore warning is given against
``sighing in one's innermost being'', and generally against letting feeling
have its full course, whether this led only to
getting air, or brought positive relief. Such
scruples Christianity did not have. It has
saved the poetry of life by taking up life's most opposed
moods, the highest jubilation and the deepest sorrow, into
its worship and giving them expression in great images that
represent the great drama of the race. The individual is led here
under a great and comprehensive view of life
and feels his lot as determined by the course of this drama.
In the great expectation of fulfillment and judgment there was
possibility of a rich and deep life of feeling, which could
raise the mind high above all outer and inner vicissitudes\footnote[2]{See my comparison between Stoicism and Christianity in the essay \textit{Hedenske Sandhedssøgere} (Mindre Arbejder I) and \textit{Epiktet og Kristendommen} (Mindre Arbejder III).}.

% --- p. 160 ---
\setcounter{footnote}{0}%
\opage{160}Deepening of the life of feeling through exalted images
and thereby the foundation of a great poetry of life is from
a general human standpoint Christianity's contribution
to the history of consolation. Yet it was not the proclamation
with which Christianity itself came forward
--- as indeed no great spiritual movement
has ever brought precisely the fruit it itself meant to
bring. Christianity wanted to bring the highest truth,
to put the demand for truth in a new and strict sense, and
the fulfillment of this demand was set as an unalterable
condition for consolation to be able to come.

The truth around which primitive Christianity gathers everything
is that the breakthrough of the Kingdom of God is near, and that what matters
is to be prepared to receive it. Consolation is granted only
to those who with a full heart enter into the great expectation,
let their life be inspired by it and esteem all else
as small. Therefore the founder of Christianity could say
that he had not come to bring peace on earth,
but to kindle fire and awaken strife.

It could not be otherwise than that Christianity,
in its interaction with the development of culture which
primitive Christianity had neither expected nor foreseen,
had to undergo changes and displacements. And the element of consolation
has then decidedly pushed itself forward before the element of truth.
And often consolation, the consolation Christianity
proclaimed, came to stand far out on the horizon as a
last refuge, while the immediate joy and security
in life prevailed in the foreground. And people perhaps even found
the truth that was revealed in Christianity
to consist precisely in the consolation it brought. Had it
been possible to show that the truth Christianity proclaimed
as the highest could embrace, or at any rate agree with,
what human inquiry can recognize as truth,
such displacement would not have been necessary.
% --- p. 161 ---
\setcounter{footnote}{0}%
\opage{161}But it may be said to be basically admitted on
all sides that the significance of religion, especially of Christianity,
can no longer consist in satisfying the intellectual
demands of the human spirit. Involuntarily one
then, when one wishes to come to terms with Christianity,
holds to other sides of it and puts them in the front rank.
On this I have spoken at length in my Religionsfilosofi and
therefore do not go further into it here. I have said in the
book mentioned (§111) that it is a sign of religion's
changed position in the life of the spirit that it more and
more takes on the significance of consolation: ``Religion was once
the pillar of fire that went before the human race on its great
march through history and showed it the way. Now it is becoming
more and more the ambulance that follows after the march and
picks up the weary and wounded.'' --- I have been
told that these words have been quoted from a Copenhagen
pulpit and regarded as a vilification
of Christianity. Whoever will read the words in the context
in which they occur will see that this is
a misunderstanding. It is a historical remark which
one may try to contest, but which one will probably be
obliged to let stand. --- If truth is to be the first
and decisive thing in Christianity, it can be maintained only
from a standpoint like that which Pascal and Kierkegaard
took\footnote[1]{For \textit{Kierkegaard's} part reference may be made especially to his \textit{Papirer} (ed. by Kuhr and Heiberg) VIII,1. p. 209. ``Christ is not love, and least of all in the human sense; he was \textit{the truth}, \textit{the truth}, absolutely, therefore he could not merely defend, but he had to let men become guilty of his death, i.e. make \textit{the truth} manifest to the utmost.'' The thought here expressed he has developed more fully in \textit{Tvende ethiske-religiøse Smaa-Afhandlinger}. (Samlede Værker VI)}.

What for a time halted Kierkegaard and kept him
from asserting his sharp and strict conception
% --- p. 162 ---
\setcounter{footnote}{0}%
\opage{162}was precisely regard for the many who, undisturbed by
intellectual and other scruples, had found
their consolation in Christian ideas and images, ---
``the weak, simple, sick and sorrowful'', ``the many
happy, lovable existences'', whom he did not wish to terrify.
But the halt could, for the prophet of intellectual and ethical
honesty, only be provisional.

Meanwhile there lies a real problem in the consideration
that caused the provisional halt. The right consolation
requires, as we have seen, a spiritual activity, puts a
demand to him who is to be consoled. In face of such a
demand people will --- as in face of ethical demands\footnote[1]{On this see my \textit{Etik} IV,2; VIII,6 where it is shown that one cannot, fundamentally, put the same ethical demand to all either in kind or in degree.}
--- meet with highly different presuppositions,
and what is to be able to console them must therefore be
of highly different kind and degree. What matters is to release
the minimum of self-activity that must be
possible as long as life lasts. Therefore it is an art
to be able to console; it requires knowledge of human nature and freedom from prejudice.
The true consoler is he who is able to
loose the living force and discharge it in a way and
in a direction that accords with the individuality of the individual.
Here no one is excluded from consolation, provided only that the
right consoler (whether it be the individual himself or
another) comes. Whether he always comes, or can
come, psychology can say nothing. When it
is asserted that consolation is always possible, this is perhaps well-meant
rhetoric, but it is not thinking.

% --- p. 163 ---
\setcounter{footnote}{0}%

\opage{163}\essayhead{FEATURES OF THE HISTORY OF PRAYER}

It is of great psychological interest when, with regard to
a significant manifestation of mental life, a certain course of development
can be demonstrated, at any rate within particular
periods of culture, --- a series of metamorphoses
or displacements in a certain direction, so that in interaction
with the different historical presuppositions
there emerge types of more or less pronounced
peculiarities. Not as if the final form of such a course of development
were to signify a result attained once for all. There may always be a possibility of
new beginnings, also of a repetition of types that
seemed to belong to an earlier, supposedly surmounted
stage. Different processes of development may also go on at the same time,
different series of types may form,
whose significance and whose mutual relation only a much
later age will be able to make clear. A psychological and historical
consideration must always keep an eye open for such
possibilities. For the individual human being in the
practice of life, on the other hand, the task is first and foremost
to find the type to which he himself belongs according to his peculiarity
and his experience of life. It is the greatness of the principle of personality founded by
Socrates that it makes
the same demand at all stages and under all forms
of human development. Everywhere it is a matter of knowing
oneself, one's ability and one's need in relation to the conditions of one's own
% --- p. 164 ---
\setcounter{footnote}{0}%
\opage{164}life and one's own time. Such self-knowledge
is the most important presupposition for honest and healthy development,
hence for the emergence of new valuable types, and
it also marks the point at which people of widely
differing spiritual orientation understand one another best.

As an example, some features of the history of prayer
within different periods of culture shall be given here. I
arrange the examples I have at my disposal under three
main forms. The arrangement follows a standard
that presents itself naturally the more reflection is directed
to this mental phenomenon's connection with
personal life and its relation to it in general. But this does not mean,
as already indicated, that such a mode of development
is run through once for all in the history of the human
life of the spirit.

\subhead{A. Prayer and Incantation. Social and Individual Prayer.}

a. There is a stage within otherwise widely different
cultures at which prayer, like sacrifice, is ascribed a
double effect, an effect on the god to whom prayer is made,
and an effect on the praying human being himself. The god
is strengthened and nourished by prayer, and his self-feeling is addressed
by one's expressing one's dependence, one's fear,
one's hope or one's admiration before him. Prayer can
itself be regarded as a sacrifice, in that man gives up
his own power and pride. But man himself receives
redress for what he sacrifices, in that prayer gives him
power over the god\footnote[1]{Examples from Indian, Persian, Egyptian, Greek and Roman religion in \emph{Westermarck}: \textit{Origin and Development of Moral Ideas}. II. p. 653---659. \emph{Westermarck} would even (relying especially on certain Moroccan customs) interpret the word religion as derived from religare, to bind: the god is bound by the prayers, wishes or curses of human beings (I. p. 584).}. Faith in the
% --- p. 165 ---
\setcounter{footnote}{0}%
\opage{165}efficacy of a solemn prayer rests on the presupposition that by the power
of the word one can enchant the god, compel his will --- if
one only names him by his right name\footnote[1]{\emph{Usener}: \textit{Gøtternamen}. p. 335.}.

In two famous poetic works the incantatory
power of prayer forms a decisive motif. In Aeschylus' ``Eumenides''
Apollo tells the goddesses of vengeance that he will
deliver Orestes from their pursuit, because he has called upon
the god in trust of his protection: ``For dreadful
is, to god and man, the wrath of him who prays.'' (Brøndsted's
translation, p. 129). In the Indian drama ``Sakuntala''
the heroine's misfortune is brought about by her, captivated
by the happiness of her love, not hearing a pious sage's
message, and he then pronounces his curse upon her,
a curse that he himself cannot lift without incantation.
(Hammerich's translation, p. 44).

On more primitive levels, too, prayer and sorcery pass
into one another without any clear
distinction being possible. The belief in the power of the word is so great that the
mere telling of an event can be regarded as a sufficient
cause to bring about the event itself.
Frazer supposes that many ancient tales which we
are accustomed to regard as mere myths have been
used in magical rites, it being thought that when they were recited
in a definite, traditional manner, they brought about
what they were about. A distinction can then be drawn between narrative
spells (narrative spells) and imperative
spells (imperative spells)\footnote[2]{\emph{Frazer}: \textit{The Spirits of the Corn and of the Wild}. I. p. 105—106.}.

The power of prayer can be dangerous. Among the Ainu of Japan
the women must not pray --- lest they gain power
over the men! And for the same reason the Santals avoid
revealing the names of their household gods to the women.

The power that in prayer and incantation is placed in
% --- p. 166 ---
\setcounter{footnote}{0}%
\opage{166}the hand of man can give rise to two different consequences.
Through the omnipotence of prayer it can make man
himself a god. Thus the Indian word
Brahma denotes both the force by which the gods live and
work, and the force by which man through
prayer wrests from the gods the fulfilment of his wish. An
identification of the two lay near at hand, and can be traced in the
peculiar development that in ancient India led
to monotheism\footnote[1]{Cf. Religionsfilosofi\textsuperscript{2}. p. 151.}. But another
consequence can also be drawn, by maintaining that only lower
beings are subject to the magical power of the word, not
the true gods. In various religions this leads
to the application of the distinction between gods and demons,
which originally\footnote[2]{See e.g. \emph{Bezold}: \textit{Die babylonisch-assyrische Religion}. (Kultur der Gegenwart. II,3,1). p. 44.} is due to the victory of a new religion over an
older one.

b. Prayer and incantation, and the magical power that
can be exercised through them, must of course always be
acts performed by definite human beings. But
are they performed by these human beings on their own account,
as the expression of an individual ability, or is it a social
work that is performed through them? It is the relation
between magic and religion that we are facing here. Frazer
has set up the theory that magic precedes religion.
Before man recognizes his dependence,
realizes that he is not capable of everything, does not have complete
power over nature, there passes, according to Frazer, a period
in which a single individual believes himself able by his
formulas and actions to intervene in the course of nature.
But gradually man learns that the hope of
directing nature by magical means must be given up, and
he regards the gods as the sole possessors and
% --- p. 167 ---
\setcounter{footnote}{0}%
\opage{167}holders of the power that he himself believed he possessed. Only now
does the difference between gods and human beings emerge decisively,
and man no longer feels himself the
equal of the gods. Prayer and sacrifice now become parts of a
common ritual, and it is regarded as outrage and sacrilege
when the individual on his own account and by his own
way sets these forces in motion. Magic is therefore,
according to Frazer\footnote[1]{See especially \textit{The Magic Art}. (1911). I. p. 233 f; 374. --- Earlier in \textit{Adonis, Attis, Osiris}. (1906).}, older than religion, if by religion one
understands a worship, sprung from dependence,
of beings more powerful than human beings. As
religion develops, magic becomes
``the black art''.

This theory gives a simple psychological explanation
of the origin of religion through the experience of the limitation of
human power and a feeling of dependence conditioned by it.
I have therefore earlier laid weight
on it. But historical development does not
always proceed according to simple psychological formulas. And one
must grant the French sociological school that, however
far back we go, we always find religion
socially organized. The individual frees himself only gradually
from full domination by social custom and tradition.
According to Durkheim and his disciples, religion therefore precedes
magic, just as the social always precedes
the individual. Special weight is laid here
on the fact that in primitive societies the magician is trained and practiced
through instruction and initiation, and the beings
with whom he comes into contact through his often ecstatic period of initiation
and from whom he draws his power are those whom the tribe
worships. The young magician must in himself and for himself through
visions experience the truth of the accounts that his tribe's
tradition contains. He is thus himself, in his greatest
% --- p. 168 ---
\setcounter{footnote}{0}%
\opage{168}fullness of power, a social product and a social representative.
Only later can magic enter into a sharp relation of opposition
to religion\footnote[1]{See, besides \emph{Durkheim}: \textit{Les Formes Elémentaires de la Vie Religieuse}. (1912). Livre I. Chap. 1. \emph{Hubert et Mauss}: \textit{L'Origine des Pouvoirs Magiques} (Mélanges d'Histoire des Religions). (1909). --- Cf., on my position toward the French sociological school, the essay \textit{Sociology and the Philosophy of Religion} in the present volume.}.

By virtue of its original close connection with
incantation, prayer acquires a magical character and at the same time
a social character. But just as it gradually
sheds its magical character, so it also acquires
more and more an individual character. The social forms
become individualized. It is no longer enough that the individual in
his personal experience [Erfaring] seeks to experience what is reported
of the beings whom the tribe worships, but he has
new experiences of his own, and his prayer thereby acquires a more
personal character. In primitive religions (i.e. those
lowest in our eyes) decisive examples in this direction can hardly
be found, although there is no lack of examples of individuals who take up a doubting
attitude toward tradition. In Egyptian
religion, on the other hand, both the differentiation of prayer from incantation
and the individualization of prayer come to light in a remarkable
way.

The history of this religion shows us all the forms of prayer,
from the magical stage, where prayer and sorcery
are one, to the inward prayer that springs from personal disposition. At earlier levels hardly any distinction can be made
between prayer and sorcery; it cannot be seen from the content of the prayer
whether it is more than a spell; it depends
on the use, not on the content. In the old
rigid Egyptian state religion all could pray the same
prayer, and the weight was laid on the words that were uttered. But
% --- p. 169 ---
\setcounter{footnote}{0}%
\opage{169}later, in the remarkable movement in the direction of personal
inwardness that marks the religious development in
Egypt in the thirteenth and twelfth centuries B.C.,
prayer became ``a revelation of inner personal experience
or an expression of an individual relation to God.''
There is an express warning against using too many words, since
it is not at all the words that matter. ``Be not
of many words, for in silence you will gain what is good.''
``Pray then with a yearning mind, whose words are all hidden.''
And God is called ``the God of the Silent.'' ``You sweet fountain for
him who thirsts in the desert, are closed to him who speaks,
but open to the Silent.''

The development in the direction of personality and inwardness
that marks a high point within Egyptian
religion was succeeded by a new stiffening under the dominion of a narrow-minded
priesthood\footnote[1]{\emph{James Henry Breasted}: \textit{Development and Thought in Ancient Egypt}. (1912). p. 95; 272 ff; 349; 356.}.

There is an inner connection between prayer's liberation
from the character of incantation and its transition to
being an inward expression of the need of the individual personality. In incantations it will be a matter of
using quite definite words and actions, just as
one must everywhere use certain definite means
in order to reach a definite goal. Through the process of liberation
of which we have seen examples, prayer is led onto
a path of inwardization that can lead ever further. ---

When the weight is laid on the words and on their use
for the attainment of external ends, it is natural that
weight is laid on many and long prayers. It is characteristic
that 1200 years before the Sermon on the Mount we hear a
warning against using many words; it is a sign of
a process of inwardization which, unfortunately, was interrupted
% --- p. 170 ---
\setcounter{footnote}{0}%
\opage{170}by the power that has caused so many inward and personal
mental processes of development to cease.

Lotze cites somewhere a German popular saying, which
has evidently been called forth by the experience of an inverse
relation between the quantity of words and the inwardness of mood:
``Those who pray too much pray themselves straight through
heaven and are set to guard geese on the other
side.'' ---

\subhead{B. The Inwardization of Prayer.}

When prayer does not stand as a magical force, and
when it is not merely an involuntary or voluntary repetition
of a formula given in tradition, the
way is open for inwardization. Then the external ends
of prayer can fall away, and prayer can more and
more have its value in itself as a closed, concentrated
state, a state of unity, in which in the end
the duality of him who prays and of him to whom prayer is made
can also fall away, not merely the duality of the prayer
itself and the good that is prayed for.

A transitional form is prayer as an expression of
praise, admiration or gratitude. The prayer of thanksgiving
is nobler than the prayer of petition. The most beautiful and
most sublime of the Old Testament psalms belong
here. The eighth psalm, with its praise and
thanks for the glorious place assigned to man in
the world of creation, has been called the most beautiful that the religious literature of
all times has produced in this direction\footnote[1]{\emph{Wundt}: \textit{Völkerpsychologie}. II,3. p. 665 f.}.
The Stoic Epictetus knows only the prayer of thanksgiving, not
the prayer of petition, and his thanks concern, like those of the ancient
Psalmist, the glory of man, --- only that he as
a Stoic finds this glory in the fact that man
% --- p. 171 ---
\setcounter{footnote}{0}%
\opage{171}finds his true happiness in himself, not in anything that
is outside it\footnote[1]{Bonhöffer: Epiktet und das neue Testament. (1911). p. 65.}.

Already when the prayer of petition concerns spiritual goods,
which in the end depend on man's own mind and own
will, it has lost its externality. Such a prayer
is offered by Socrates in Plato's ``Phaedrus'', when after a
profound conversation with his friends in the sacred grove he
proposes to them the following prayer before they depart: ``O dear
Pan and you other gods who are at home here, grant me to
gain a beautiful inner self, and that the outer goods may be in harmony
with the inner!'' (Phaedrus p. 279 B). --- In the dialogue ``Alcibiades
II'' Socrates cites with approval a verse by an
older poet, who prays Zeus to give us what is good,
whether we pray for it or not, but to avert what is
evil, even if we pray for it. (Alcibiades II p. 143 A).
That this dialogue is hardly by Plato is, in this connection,
immaterial. It betrays the influence of Stoic
thoughts.

Platonism and Stoicism have had a very great influence
on the later history of prayer. Neoplatonism
carries through what Plato and the Stoics had begun.
Plotinus taught that what is experienced in our highest states,
which in his case took on the character of ecstasy, cannot be
thought or said as long as they last; afterward one can
try to interpret what was experienced, but it always remains
patchwork. We can have the Highest in us, even if we
cannot say what it is, just as the inspired and the
seized only know that they have something greater in them, but do not
know what it is. (Ennead. V, 3, 10; 5, 9). Prayer proper
he laid no weight on, any more than on worship
in general; the gods were to come to him, not the reverse\footnote[2]{See the essay Experience and Interpretation.}.
But there were those among his successors who swung back
% --- p. 172 ---
\setcounter{footnote}{0}%
\opage{172}to worship and popular religion, in part in a very
crass form. A testimony to this is the work ``The Mysteries
of the Egyptians'', which has been attributed to Iamblichus and in any
case dates from his time (c. 300 A.D.). The opposition
between philosophy and mythology falls away. The endeavor
is to renew the old popular religion
under the influence of the new oriental currents and
to defend it against what was at that time the modern
doubt. Prayer now has value only when it is offered
in holy words, that is to say here, words approved by tradition
--- and preferably in words that derive from the
holy peoples of the Orient (Egyptians and Assyrians), to whom the gods have
revealed the mysterious language. Even if these words
are unintelligible to us, the gods nonetheless understand them!
Prayer thus becomes magic again. But it is also maintained,
--- what will prove to mark a significant
turning point, --- that it is not prayer that turns the gods'
mind toward man, but that it is the gods' own
gracious mind that shines down upon those who pray, in that
it has already before prayer and address worked upon them.
Three steps are distinguished in the effect of prayer, a
threefoldness that often recurs in the later history of mysticism
(thus among the Victorines and Santa Teresa). The
first and lowest step consists in collection, concentration,
address toward a single goal. The second consists in
an inner connection with the gods. At the third step
the unity with the Godhead is completed (the unity that
Plotinus attained through the self-immersion of thought)\footnote[1]{\emph{Harless}: \textit{Das Buch von den ægyptischen Mysterien}. (1857). p. 73---79.}.

Not all Neoplatonists went along with such a mixing
of thought and magic. The last more significant
Neoplatonist, Olympiodorus of Alexandria, lays
the weight on the purely inward character of prayer. He
% --- p. 173 ---
\setcounter{footnote}{0}%
\opage{173}develops in his commentary on the Phaedo, following a line of thought
in Plato, that the highest human striving
ultimately has its end in itself, and to him who will not
concede this he points out that we do after all give
ourselves up to prayer for the sake of prayer itself. This is thus
for him a self-evident proposition. And in another context
he returns to it. He shows that thinking
is not something we can settle once for all; it
matters not to have thought but to remain in
thinking activity, and he illustrates this by prayer:
since prayer is in itself happiness, it is something we will
remain in\footnote[1]{\textit{Olympiodori Philosophi in Platonis Phaedonem commentaria}. ed. W. Norvin. (1913). p. 19; 95.}.

The Neoplatonic thoughts recur in the mysticism of the Middle Ages,
which everywhere seeks to translate the external relations
that religious ideas denote into inner ones. That ``the Kingdom of God
is at hand'' means for Master Eckchardt that it is in us,
for God is nearer to us than we are to ourselves. As here the temporal relation
falls away, so too the difference between
prayer and the hearing of prayer falls away. God is moved as little
by all prayers and good works as if there were no such thing at all.
``In a first eternal glance --- if
we can assume a first glance --- God saw all things
as they were to happen\dots{} He saw also the least prayer
and good work that anyone would perform, and he saw
which prayer and which devotion he would hear.''
On the other hand, the right prayer is that which springs
from a free mind that is not bound to anything here in the world,
but is wholly sunk in the will of God. In inward prayer
man at last becomes one with him to whom he prays, with
God\footnote[2]{\emph{Mester Eckehart}: \textit{Prædikener og Traktater}. Translated and introduced by Aage Markus. (1917). p. 69; 78; 95.}.

% --- p. 174 ---
\setcounter{footnote}{0}%
\opage{174}Even as late as Kierkegaard's religious discourses we meet the
same thoughts, rediscovered and personally experienced, but belonging
to a great tradition that is connected with the most
inward forms of religiosity.

According to the naive conception --- says Kierkegaard ---
the main thing is that God gets to know what one has
on one's heart. But the reverse is the case. What matters
is that the one who prays continues his prayer
until he gets to know what God wills. This
thought is developed at length in one of the religious discourses, whose title
is: ``The one who prays rightly struggles in prayer and prevails ---
in that God prevails.'' The inwardness of prayer is the main thing,
and prayer is not merely a means to attain
something. Already the collection that prayer presupposes, ---
that man gathers his whole soul about a single
wish, --- already that ennobles, and it is the preparation
for giving up everything; for only he who has but a single wish
can give up everything. It is God's unchangeable will
before which man bows, and thereby, --- by not himself
wishing to be anything, --- he can be wholly shone through by God
and come to resemble God. --- In a later note
Kierkegaard maintains that one must learn always to give thanks.
``Otherwise all this childish nonsense arises, one
moment one is to give thanks, because now this or that is good
and this and that is so pleasant, the next moment one is to sigh
and pray for patience, because now this and that was so
bad. It is impossible to have a real relation to
God and remain in one's merely human and earthly idea
of good and evil, the pleasant and the unpleasant.''\footnote[1]{\emph{Kierkegaard}: \textit{Samlede Værker}. V. p. 149—167. Cf. \textit{Papirer} (ed. by Heiberg and Kuhr) VII,1. 24; VIII,1 p. 123.}.

In Neoplatonism and mysticism it is the thought of unity that
rules. The endeavor is to get beyond
the opposition between prayer and that which is prayed for,
% --- p. 175 ---
\setcounter{footnote}{0}%
\opage{175}and beyond the opposition between him who prays and
him to whom prayer is made. In Kierkegaard's religious discourses it is
the thought of the divine immutability, hence
also a Platonic thought, that leads beyond the usual
theory of prayer. --- It is of course not Kierkegaard's
definitive standpoint that we have before us here. In his most profound
work (``Sygdommen til Døden'') he says:
``To pray is to breathe, and possibility is for the self what
oxygen is for breathing\dots{}. God is this, that everything is
possible, or that everything is possible is God''\footnote[1]{\textit{Samlede Værker}. XI. p. 153.}. Immutability, on the other hand,
does not permit more or other possibilities
than there have always been. ---

It might seem that it was merely a terminological
question whether one would use the word prayer of an act
in which the mind is wholly one with its object, with
the Immutable, and in which there is no difference
between the act itself and that to which it is directed. Brøchner
maintains accordingly that when one sees the effect of prayer in
the indirect, subjective influence on the one who prays himself,
the essence of prayer is in reality undermined
by criticism. For prayer aims at the miracle, and
only at the miracle. Brøchner assumes, on the other hand, that there
is a rational equivalent for prayer, namely the mind's
address to the infinite Being with which it is in
its innermost ground one, and in the force that through
such an address to the ultimate source of our being can
be released in us\footnote[2]{\emph{Brøchner}: \textit{Problemet om Tro og Viden}. (1868). p. 95ff. --- \textit{Det Religiøse i dets Enhed med det Humane}. (1869). p. 17.}. It is Master Eckehardt's standpoint
that Brøchner takes up, except that he will not call the
inward address that is possible at such a stage
prayer. In a French dissertation whose subject was a psychological
study of prayer, the author (Segond) had
% --- p. 176 ---
\setcounter{footnote}{0}%
\opage{176}precisely found the essence of prayer expressed, not among primitive
people, whose inner life we cannot reconstruct, but
among the mystics, in whom prayer is conceived with clear
consciousness of its nature. His opponents (Durkheim,
Delacroix and Delbos) pointed out to him the
arbitrariness of this limitation and maintained that the relation
between the prayer of silence and the prayer of the word ought to
have been more closely examined\footnote[1]{See the interesting résumé of the discussion in Revue de Métaphysique et de Morale Janvier 1911. Supplément p. 30—35.}. The opponents were evidently
right that it is arbitrary to call a single type
the essential one. A whole series of problems is thereby
passed over.

The question is, however, not merely terminological. History
seems to show that a successive transition takes
place between the different types, even though, when they are
each fully developed, they may come into sharp opposition
to one another. This continuity is of great
importance for avoiding waste of force; for through it
a new type can draw nourishment from the stored-up
force of the old one, and moreover it makes understanding easier, even though
this understanding must essentially be won by way of analogy.
Brøchner is right, however, that for the sake of clarity it is
expedient to mark the difference between the
different types by different names, even though it
may become difficult to decide where the decisive
line is to be drawn. The church has often drawn it where
dogmatics is threatened by bold mysticism. From the standpoint of the psychology of religion
one must draw it on the other side of
mysticism, where the great problem of the psychological
possibility and nature of the equivalent Brøchner
mentions arises. --- A question different from this
is which type is the highest. For the time being everyone, no doubt, regards the stage on which he himself stands as
% --- p. 177 ---
\setcounter{footnote}{0}%
\opage{177}the highest\footnote[1]{Cf. the dispute between Baggesen and Grundtvig over whose stage was the highest. It could, according to Grundtvig, be decided only on Judgment Day. \textit{Danske Filosofer}. p. 70f.}. What will be decisive is whether a single
type can be shown to contain in itself everything that is
valuable and justified in other types. It lies outside the
present investigation to go further into this.

As already stated at the beginning of this essay
there is not a single series of development, but
several series at once, now ascending, now descending,
as regards the relation between duality and unity.
And in great personalities who have intervened deeply
in the development of positive religion one will often find
several different types asserting themselves. In Jesus of
Nazareth, as the Gospels portray him, most of the types mentioned in the foregoing occur.
Toward the demons he uses incantation, but
toward God he uses the prayer of petition and the prayer of thanksgiving,
and here too the prayer of thanksgiving stands as the more sublime,
since it is the expression of the most glorious experiences. In the Gospel of John
the prayer of unity rules, and it is this
that has borne fruit in the development of church dogma. The high point
is, however, the prayer in Gethsemane, --- the most religious
word that has been spoken on earth. Here Jesus struggles
his way forward from duality to unity, from self-will to full
absorption in the divine will, which expresses his
calling. --- In the greatest personality within the church,
Augustine, we encounter (see Confessiones) also, partly successively,
partly in the same period, several different types.
As a child he prayed that he might not be flogged
at school. In his youth, when oratory and philosophy occupied
him, his prayers also changed with his mood (affectus),
and now aimed at gaining wisdom. When the
need to raise himself above his sensual passions
% --- p. 178 ---
\setcounter{footnote}{0}%
\opage{178}began to make itself felt, he prayed for chastity
and self-mastery; but he confesses that he feared
being heard too quickly. On the other hand, when the new
faith was about to break through in him, he prayed for
its victory in his mind, and now he says to his God: ``Why
not at once?'' A childlike prayer of petition appears also
after the conversion, when through his own prayers and his friends'
he was delivered from toothache. But he prayed in vain
to be delivered from the deep sorrow over the death of his
pious mother, and he knows why he was not
heard; God wished to impress upon him the importance of
the persistence of an inward relationship, even after the other party's
departure. The high point in the development of prayer was it
for Augustine when he could pray: ``Give me what you
command, and command what you will!'' ---

In all its forms prayer belongs closely together with
need and impulse, wish and longing. But there is a
great religion according to which impulse and wish are precisely
what is not to be, since they increase or ever anew
bring forth the tormented existence which it is precisely a matter
of bringing to an end. For Buddhism\footnote[1]{\emph{Warren}: \textit{Buddhism in translation}. (1906). p. 177. 189. 245.} prayer is karma
(existence), and conversely karma is the source of the yearning
and striving that find expression in prayer. This is the most
radical standpoint toward prayer as a mental phenomenon.
Prayer is the expression of a craving for life, an impulse to
hold fast to existence, even if it is to be under
quite different forms from those hitherto experienced, and even if
it is not one's own individual existence that is sought
to be maintained. There is a possibility of metamorphoses into new mental
phenomena, which one perhaps will no longer call
prayer, into a free spiritual and material urge to development. But
Buddhism takes this whole basis and all the possibilities
it contains away from the root.

% --- p. 179 ---
\setcounter{footnote}{0}%
\opage{179}In the characterization I have attempted in my Religionsfilosofi
to give of Buddha and Jesus, I have precisely laid weight
on this point. Buddha rejects all yearning and striving,
Jesus awakens a great expectation and thereby lays
the weight on striving and work.

\subhead{C. Prayer as Gift --- and as Symbol.}

On several occasions in the foregoing the thought has emerged
that it is a gift to be able to pray, thus
among the Neoplatonists and in Master Eckchardt. If
this thought is carried through, we stand at a turning point in the history of
prayer. It then becomes not, or not merely, nourishment,
but effect, and it is this that becomes the essential thing.
And then in the end the form of prayer, that something is prayed for
and to someone, becomes symbolic, a poetic expression of
yearning and longing for something that fundamentally already
lies in our own being. Prayer then works not merely
as releasing force, but it is itself the consequence of a release of force
which at the given time it was not in the individual's
conscious power to bring about. --- This viewpoint
reappears with different grounding. I
believe one can distinguish between a supernaturalist, an
idealist and a naturalist grounding.

a. It is told,\footnote[1]{Tiele: Elements of the Science of Religion. II. p. 133.} that to the Persian prophet Jelaled-din
Rumi there came a man and complained that Allah
never answered his prayer, --- that he never sensed
Allah's presence. Then the Prophet said: ``God commands me
to say to you: You sorely tried man, was it not
I who drove you to worship me? Your very invocation
of Allah was an expression of my presence, and your sorrow,
your longing, your zeal were my messengers!'' --- Prayer itself,
that prayer is made, is thus a revelation of the Godhead.

% --- p. 180 ---
\setcounter{footnote}{0}%
\opage{180}A similar thought emerges when Paul writes:
``The Spirit comes to the aid of our weakness; for we do not
know what we should pray for as we ought, but the Spirit
himself intercedes for us with unutterable sighs. And
he who searches the hearts knows what the Spirit's yearning
is.'' (Romans VIII, 26). And thereupon there is talk of
divine predestination and calling. --- In prayer
man is thus not the giver or the one who sacrifices,
but essentially the recipient. In the very unformed
prayer, ``the unutterable sighs'', a divine spirit stirs
in man's interior, without man being conscious
of it\footnote[1]{Cf. on this remarkable passage \emph{Bernhard Weis}: \textit{Lehrbuch der Theologie des neuen Testaments}.\textsuperscript{6} p. 343.}.

Prayer seems to be directed forward and upward; but
for a deeper conception there is thus behind prayer as
forward-yearning a force that (like a vis a tergo) pushes
and drives forward. This conception becomes the diametrical
opposite of the stage where prayer and incantation
merge into one, and where prayer as a magical force
gives man power both over God and over the world.
We cannot expect to find this deeper conception
except in individual outstanding religious personalities
at different times. But it must be a strange
standpoint from which this should diminish the value
of the conception.

It is shown in the essay on Experience and Interpretation
how this conception emerges in Santa
Teresa in the decisive opposition she sets between
the prayer of reflection (oraison de méditation) and the prayer of quiet
(oraison de quiétude). To the prayer of quiet one cannot
raise oneself up, but must wait until one is drawn
up, just as a garden must wait for its best
% --- p. 181 ---
\setcounter{footnote}{0}%
\opage{181}watering, the lovely rain that falls quietly down from
above.

Pascal, for whom it was certain that he could not
have sought Jesus if he had not already found
him, maintains outright that prayer is a gift that God
gives to whom it pleases him, and he rejects definitely
the assumption that one has prayer from oneself. It is
a Jansenist proposition that Pascal here maintains. It goes
back to Augustine in his polemic against Pelagius.
Even Bossuet therefore adheres to it, in saying
that prayer, by which everything is given to us, is itself given to us\footnote[1]{\emph{Pascal}: \textit{Pensées}. ed. Brunschwicg. II. p. 408; 438, with the editor's note. --- Cf. \emph{Sibbern}: \textit{Gabrielis' Breve}\textsuperscript{1}. I. p. 100: ``To be able to pray is true, full grace.''}.

b. When, in contrast to the supernatural explanation
of the origin of prayer, which starts from a duality
between divine and human nature, I set
the idealistic explanation, it is because here the
divine nature is thought of as working in and through the
human nature, just as it works through nature
in general.

In the Indian poem Bhagavadgita both faith
in gods and prayer to gods are explained as the effect of
a power that does not itself come to man's consciousness
and cannot be observed, --- of the absolute, unchangeable,
supreme Being itself, which stirs in everything. Krishna,
who is an incarnation of Vishnu, says in the said
poem: ``Those whose knowledge has been taken from them by
one desire or another approach other gods than
me, and bound to their own nature they take upon themselves all manner of
religious duties. When now one of them wishes in
faith to serve one of these gods, then I make this faith
unshakable in him. Filled with this faith he then
worships this god with care and attains the good that he
% --- p. 182 ---
\setcounter{footnote}{0}%
\opage{182}wishes; --- but it is in reality I who grant it to
him''\footnote[1]{\textit{Bhagavadgita} VII, 20—24. (Garbe's German translation).}.

In the European world such a thought has, through
Neoplatonism and medieval mysticism, passed over to
modern theologians who do not feel bound by confessional
ideas as final. That prayer is made,
and that something is attained through prayer, stand
as two experiences, which both spring from one source
of which human beings themselves have no inkling when they do not follow the path of
reflection to its end. According to Schleiermacher a theory of
prayer that presupposed a reciprocal relation between the Creator
and the created would mean a regression to the standpoint of magic.
Prayer and the hearing of prayer are fundamentally one, and in right
prayer passivity and activity are united. Prayer
and its fulfillment or nonfulfillment are only parts of
one and the same divine order\footnote[2]{\emph{Schleiermacher}: \textit{Der christliche Glaube}. § 47.}. In the same direction
goes later Auguste Sabatier's conception\footnote[11]{See my characterization of Sabatier's standpoint in Mindre Arbejder. Anden Række. p. 75—78.}.

c. As examples of what I mean by a naturalistic
conception in its application to prayer, I will
adduce statements by two spirits widely different from one another,
one of whom is the founder of the modern conception of
nature, the other one of the greatest figures in the history
of philanthropy.

Thomas Hobbes seeks to show that prayer may very well
find its place, even if one assumes that there are definite
natural causes for everything that happens. ``Prayer,'' he says\footnote[12]{\emph{Hobbes}: \textit{Tripos}\textsuperscript{3}. (1684). p. 293.}, ``is God's gift just as blessing is
\dots{} And prayer to Almighty God is only thanksgiving
for God's blessings. Even if it precedes the
% --- p. 183 ---
\setcounter{footnote}{0}%
\opage{183}special things we pray for, it is yet not the cause or
means thereto, but a declaration that we expect
nothing other than what God wills, not what we will.''

Florence Nightingale was not only one of the greatest
and wisest friends of humankind who have ever lived,
but she was an independent thinker. Stuart Mill's Logic,
Jowett's Plato scholarship and Quételet's statistics constituted
the basis from which she built her
own thoughts. The wonderful thing about God is, according to her
conception, precisely that he is inexorable: if he
were to change his mind according to people's prayers, we would
be delivered over, at their mercy or disfavor, to those who prayed
to him! The true idea of God, which alone is compatible
both with the deepest morality and with the most comprehensive
view of nature and history, is the
idea that shows him, not as an individuality, but
as a universal being --- as a law. What the particular
laws are, we must find by the aid of science.
And on the basis of the lawfulness of nature, a basis,
in which God's being is given, there then arises
for every human being the task of making the best
of his own personality. More beautiful than any sentence in
the church rituals was in her eyes the prayer of Socrates
in Plato's Phaedrus, already mentioned above:
grant me to gain a beautiful inner self, and that the outer goods
may be in harmony with the inner! Characteristic of Florence
Nightingale is it also that figures such as Andromache
and Antigone stood in her eyes above the women the Bible
mentions\footnote[1]{\emph{Edward Cooke}: \textit{Florence Nigthingale}. Danish translation. (1917). II. p. 432—442.}.

d. When now prayer --- with one or another motivation
--- is regarded as gift, hence as effect in man's
% --- p. 184 ---
\setcounter{footnote}{0}%
\opage{184}mind, not merely as sprung from man's mind as
cause, it becomes a symbolic act, poetry, and
if it is earnest and honest, poetry of life. Already the prayer of unity
brings this consequence with it. For if one is certain
of having God in oneself, of being one with him, so that his
will is our own will, then there is nothing to pray for,
just as little as God himself can have anything to pray for.
The duality of prayer and that which is prayed for falls away
with the duality of him who prays and him to whom prayer is made.
But when prayer and the hearing of prayer become one, it
shows itself quite clearly that the formulation of prayer in words is not
necessary, or has only poetic significance. The need
for formulation, with address to a being different from the one who prays,
will rather be connected with a need for pictorial ideas than with the very
need that seeks satisfaction. In some people
there prevails a dread of the indefinite, a shyness of the unsurveyable;
they need definitely formed ideas,
preferably such as have, through the succession of generations, acquired
the character of spiritual homeliness. In other people
there prevails precisely a shyness of the limited, a dread
of narrow ideas; they discover at once where every
simile limps and where every image proves inadequate.
A genuine poetry of life points beyond this
opposition; for in it every image that is used is the expression
of a personal experience of the conditions and fate of the inner or outer
life and of an endeavor to hold fast to
its model. It is an ideal of life that the mind here turns
toward, --- exemplar vitæ, as Spinoza called it, an
ideal socius, as James has called it. Even the founder of
positivism, Auguste Comte, directed the ``solennelle
effusion des sentiments généreux'', as he called his
prayer, to Humanity, which was for him represented
by a woman who had been for him what
% --- p. 185 ---
\setcounter{footnote}{0}%
\opage{185}Beatrice was for Dante\footnote[1]{Cf. \emph{Auguste Comte}: \textit{Catéchisme positif}. p. 105 f. La vraie nature de culte intime consiste dans l'adoration journalière des meilleures personifications que nous puissions assigner à l'Humanité --- In \textit{Politique positivs}. (IV p. 114) priere is defined as commémoration suivie d'effusion.}. For many it will be natural
to give such a deep mood, which is called forth by the
endeavor to hold fast to the ideal of life, whatever conditions and
fates we may struggle with, expression in words. When a thought
breaking forth with inwardness and force gets air
in words, when others are present, why should this not
also be able to happen in solitude? A Master has indeed said that one
should shut one's door when one prays. In common prayer, which is
formulated beforehand and is to be held at a fixed time and
place, there is always something artificial, which repels many
natures, not least such whose inner life unfolds
in a peculiar and deep way. In their best
thoughts, in their most inward feelings there stirs for them
the Highest they know, and they can very well hold fast to
the faith that something corresponding also stirs in the interior of other human beings,
without needing common formulas
to establish it.

That prayer is effect does not of course exclude that
it is cause. It is not merely itself a sign of a release
of energy, but it itself leads to a new release.
It is seen from this side that prayer, which is often called
the central religious phenomenon, cannot disappear
from the spiritual economy without an equivalent. And we have
seen the possibility that the naive prayer of petition, through the prayer of unity
and through prayer regarded as gift, may pass
over into becoming a free expression of courage, trust and enthusiasm
or of inward holding fast to an ideal of life.
There are certainly not a few who in their lives have run through
this series of types. Whether in the concluding
% --- p. 186 ---
\setcounter{footnote}{0}%
\opage{186}poetry of life they can find full compensation for what at the
earlier stages was granted them through prayer, may
be a difficult matter for them to decide, and even more difficult
for others. And whether such a series of types as
we on the basis of comparative consideration can
set up will also be run through by the race as a whole, or
whether the different types we have described will not be present
simultaneously at all times, that too is a difficult
thing to decide. ---

When the talk is of equivalence for earlier religious
forms of life, there is one thing one must not forget. The prayer of petition,
in its naive forms, is no doubt that of the types mentioned
above for which one will find it hardest to imagine
a replacement. It also belongs especially in times
when religion stood as the only spiritual power. There had
not yet occurred any division of labor in the spiritual
domain. It is different now. In divided labor
religion (whether one takes the word in its widest
or its narrowest sense) cannot enter in as a
single member, but it must either be everything, or stand above
everything, or be the deepest presupposition in all special
work without itself having a field of work of its own.
Outside these three possibilities no way will be
found. --- To be sure, an attempt is sometimes made, indeed perhaps
quite often, at a peculiar division of labor between
religion and culture, so that they go side by side without
having anything to do with each other. Thus the
well-known astronomer Father Secchi is said to have said: ``There is a
staircase leading up from my oratory (the prayer chamber)
to my observatory. When I have come up above a certain
step, I think only of the observatory; before I have
come to this step, I think only of the oratory.''
So external the relation cannot in the long run continue
to be. One cannot of course think of all things
% --- p. 187 ---
\setcounter{footnote}{0}%
\opage{187}equally energetically at the same moment. But one must nevertheless set
oneself the task of uniting the experiences from the oratory in
connection with the experiences from the observatory. And even
if one cannot and will not set oneself this task, experience
will show that there is a subterranean connection between
the two groups of experiences, which one fine day
may make itself felt in and for conscious thought. ---

The most inward collection upon and self-surrender to what
is for us the Ideal must show its force at our very
place of work, just as it is itself fruit of an inner work
of which we perhaps do not become conscious.

Postscript. --- After this sketch was written down,
there came into my hands Friedrich Heiler's book: Das Gebet.
Eine religionsgeschichtliche und religionspsychologische Untersuchung
(1918). In this book a great material is gathered
for the illumination of prayer and its
history. It does not occur to me to compare my
little sketch with Heiler's detailed presentation. I have
in it only wanted to put together some examples I have in the
course of the years collected, in order to demonstrate what changes
prayer can undergo, until a stage sets in
at which one can no longer speak of prayer in the proper
sense. The author characterizes excellently the two
great types, which he calls the mystical and the prophetic
prayer, and he considers the prophetic the highest. He
demands that prayer be measured by its own (religious)
norm and rejects the ``philosophical norm'' that applies
ethical and rational points of view. But he declares that
it is a ``coarsening'' of prayer when it is connected
with faith that God's will can be changed, a faith
that lies at the ground not only of primitive prayer,
% --- p. 188 ---
\setcounter{footnote}{0}%
\opage{188}but also of prophetic prayer, thus still of
Luther's prayer. Should it not be a ``philosophical norm''
that makes itself felt in him here? From the nature of the
prophetic prayer itself the corrective cannot in any case
have been taken, for for it the hearing of prayer is an essential
moment, and it signifies a great crisis when the hearing of prayer
does not occur. Would prayer without the influence of ethical
and rational motives have been able to develop from primitive
to mystic and prophetic prayer, not to mention the highly idealistic
form the author (moreover without closer grounding)
places highest? --- It is continually the influence of the
mentioned motives that finally leads to a stage like that
I have indicated at the end of my sketch.

% --- p. 189 ---
\setcounter{footnote}{0}%

\opage{189}\essayhead{TWO PASTORS}

\essaysource{Otto Møller and Skat Rørdam. A Correspondence. Copenhagen 1915---16.}

The correspondence presented here between two significant
and earnest ecclesiastical personalities of our own
day is of interest both for the psychology of religion and for the history of religion. At every time a religion has a greater content,
a greater sum of doctrines and inherited
ideas, than individuals make use of, or can
make use of. It is therefore of great interest to see which
elements of a given religion, at a given time
and in particular personalities, prove to exert
the greatest influence, --- which of them are in reality lived
by. What first aroused my attention in this correspondence
and made me want to write a little study was
the two men's relation to a problem which in my young
days gripped me so strongly that it became decisive for
the whole course of my life, and with which I have since
tried again and again to come to terms. It is the question of
the relation between ``the Christianity of the New Testament'' and
``modern Christianity.'' Yet there are also other questions
which this correspondence can set one thinking
about, such as the relation between self-experience and tradition,
and in connection with this the relation between the
individual and the ecclesiastical society to which he belongs. But
in a purely human respect too the correspondence offers \opage{190}interest through the beautiful and warm impression one gets of the
two men's friendship, which endured with inwardness
and fidelity from their young days until death parted them.
The correspondence runs from 1859 to 1909, and the friendship was
founded before that. There were reefs enough on which it might
have run aground, and it is attractive to see how a
deep harmony in that which for them both was the essential
helped them over the difficulties. It cannot
be denied that the readers here have most reason to admire
Otto Møller's magnanimity and peaceableness. The
conservative and orthodox Rørdam often expresses his
displeasure at his friend's heresies in a very
sharp manner. One is led to think how difficult
it must be for people who stand at different
stages to come to an understanding, when one sees
how near a split may lie even for people who
in the main wish to build on the same ground.

\romanhead{I}

Otto Møller characterizes himself as a melancholic,
who lives in something that is far removed from all that
stirs at the moment, and who easily feels confused
when the age and its movements press in upon him.
All his life, apart from his school years and his student days,
he lived in Gylling parsonage, where he succeeded his father.
He was a brooder and a scholar. ``I have far too much
time to go about chewing on my own thoughts,'' he says
himself. And his studies made him one of the most learned
pastors in Denmark, --- which, by his and
Rørdam's common opinion, admittedly meant not very much. He has
been a spiritual center to which many resorted
to win peace and clarity on ecclesiastical and religious
questions. An enthusiastic and faithful adherent of
Grundtvig he was; but it was the Grundtvig of the \opage{191}earlier period that was for him the real Grundtvig. In
Møller's eyes the Master, when from the forties on things went
downhill for him, was inclined to see everything too brightly and
to make things easier than they were and ought to be.
He did not think that Grundtvig's last period corresponded to
what his strong years had promised. In any case most
Grundtvigians now came to it too easily. There were too
many meetings, too much singing and clamor. ``There is a disgusting
fooling around with Christianity\dots{} God save us from the cheerful
Christianity; it can become worse than spiritual death,''
he writes in 1860. And much later (1892) he writes:
``The older I grow, the more I get a sense of
the one-sidedness and softness of the old `genuine' Grundtvigian\dots{}
Grundtvig himself was a peculiar personality,
all of a piece in his kind; but in his disciples
it became something mannered.''

Compared with the springtime of promise that
he had experienced in his youth, he felt poor
in his old days when he looked about him at the age. The pastors
of his time he had, like his conservative
friend, very little respect for. The bishops are ``muddleheaded,''
--- and the pastors are ``a lot of riffraff everywhere.'' ``The clergy
is always the same and grants the congregation
no real freedom.'' And as his political conservatism
unites with his radical church policy, he says
(1906): ``I have little taste for being governed by
parish clerks, but still far less by pastors.'' --- It is
not only the Grundtvigians that this strikes. The
High Church (the third tendency) and the Inner Mission
he also speaks out strongly against.

His despondency and feeling of sharp opposition to
the tendencies of the age, even those that like himself wished to
stand on the ground of the Church, are explained in part by
his temperament, by his lonely life and by his
% --- p. 192 ---
\setcounter{footnote}{0}%
\opage{192}ill health; but only in part. In a brief characterization
of the two friends the Bishop's wife Rørdam (Tilbageblik
paa et langt Liv, p. 118) speaks of Otto Møller as ``the man
with the pronounced critical sense and the heavy gaze.''
He is more of a thinker, more consistent and sharp in his
thinking than his friend. And his critical attitude is conditioned
not only by his inner and outer circumstances, but first
and foremost by his ever wakeful reflection. Of course
the reflection moves only within his own ecclesiastical
circle of ideas, but here it leads him to a
position on various theological and ecclesiastical questions
which draws constant disapproval from his more
limited friend, who was bound by authority.

Both friends are pronouncedly ecclesiastical natures. They both
have their last resort in ``the congregation,'' a concept
to which it proves very difficult to give a definite content
and extent. They are both deeply convinced
that in the tradition into which they had lived themselves in that part of their
youth which they called ``the springtime,'' there
is given the one thing needful for human beings, but also in itself
rich in content and profound, expression of the highest truth.
The difference between them emerges partly in their
conception of the constitution of the Church, partly in their position on
the right which personal experience has in relation to tradition,
partly, finally, in the extent to which they permit
freer reflection to unfold in the religious domain.

Skat Rørdam was, by his wife's testimony, a
more sanguine nature than his friend, with a brighter gaze
and greater cheerfulness. The brighter outlook expressed itself not
least in his having greater confidence in the existing
ecclesiastical institutions and, especially after he had become
Bishop of Zealand, greater faith that they could work
for good. He is also, partly because of his optimism,
% --- p. 193 ---
\setcounter{footnote}{0}%
\opage{193}partly because of less critical sense and need, more
orthodox than his friend, as is often told him in the letters.
But one gets the impression of a man who lives wholly
and fully in the tradition into which he has once for all been led.
The conclusions to which his reflection leads him agree
involuntarily with the conclusions which his Church had long since
drawn.\footnote[1]{See my characterization of ``ecclesiastical natures'' in my Religionsfilosofi § 96.} Otto Møller, on the other hand, has a greater latitude
between the ecclesiastical foundation and the lines of thought
that can be struck up from it, and he both saw and followed
more possibilities than his orthodox friend entered upon.

Very sharply Rørdam often expressed himself, indeed as a rule,
about Møller's authorship. Of Møller's book on
the Gospels Rørdam says, after having read it in
manuscript, that it can only be a joy to the Lord's
enemies, and it is on the whole a very harsh letter that he
sends his friend; one understands that the latter, when he read
it, had the feeling of getting a bucket of cold water
over his head. When Møller, despite Rørdam's urgent
admonition, published it nevertheless, Rørdam wished that it
might be read as little as possible, so as to do as
little harm as possible. The matter was that Møller had expressed
himself very strongly about Mark and Luke, whom he
set far below Matthew. So decisive an opposition
Rørdam in his opinion thought one ought not to assume between
authors who were all inspired by the Holy
Spirit. --- In a later book Møller attacked the orthodox
doctrine of atonement, according to which God could show grace toward human beings
only through a bloody sacrifice that had to be brought by his own Son.
In place of the doctrine of atonement he put the doctrine of redemption, which
presupposed no change in God from wrath
to love, but only a constant divine will to
% --- p. 194 ---
\setcounter{footnote}{0}%
\opage{194}help and deliverance. When Møller sent Rørdam the book,
he was sure that his orthodox friend would ``shake his
old head over it.'' And so it proved. Rørdam
was greatly offended that Møller will not concede
that ``something was rendered to God by the death of Christ,'' and that
he at bottom regards the orthodox doctrine of atonement
as heathen. Møller shows himself here too as the most
consistent. For a God with whom there is no change
or shadow of turning cannot of course
by any event or act pass over from wrath
to love. Naturally Møller too cannot here
draw all the consequences. --- Of a following book by
Møller on a related subject Rørdam is likewise ``so exceedingly
dissatisfied''; Møller, namely, will not concede that God
is angry with sinners; he relegates wrath to the Day of Judgment
and conceives Christ's death as testimony
that God has never been angry with human beings. --- The
opposition that appears on these points between
the two men is not new in the history of the Church, but reaches
far back. The strife between the views designated
by the words redemption and atonement came forth in antiquity
in the opposition between the Eastern and
the Western Church Fathers (Athanasius and Augustine),
in the Middle Ages in the opposition between Abelard and Anselm,
and in the time of the Reformation in the opposition between
Zwingli and Luther.

Møller takes the repeated rebukes with great
meekness. ``I see more and more that you
cannot understand me, and you are not alone in that; but it
cannot help me to see otherwise than I have
once come to see these things.''

In the conception of Jesus' personality and fate too
the two friends disagree. Møller has written an account
of the gospel history, in which he seeks
% --- p. 195 ---
\setcounter{footnote}{0}%
\opage{195}to show that Jesus in his public appearance suffered a whole series of
disappointments, since neither the learned in Jerusalem nor the common man
in Galilee recognized him, indeed not even the Apostles
understood him. Rørdam is scandalized. Is it supposed to have fared
with Jesus as with a human reformer, who must pass through
so many bitter torments because of resistance? How
can he then be the Lord and the Savior? It is modern
talk, by which Møller has let himself be captivated, and which must lead
to regarding Jesus as a mere man! --- Møller finds
it distressing that he must give up writing anything
that suits his older friend. But he declares that he
has taken the matter purely historically, and that he regards the
great disappointments as belonging to Jesus' suffering. Jesus
had wishes and hopes that were not fulfilled. Can one
be a true man without that? That Rørdam calls this
conception modern does not touch Møller: for ``how much looks
new that is nevertheless true?''

Everywhere it is confirmed, what Møller says somewhere,
that Rørdam has a peculiar love for what has been handed down,
while he himself has a drive to go behind what has been handed down.
This could easily have led them away from
each other, and perhaps also had done so, had not
oral companionship and conversation from time to time
replaced the discussion in letters. ---

\romanhead{II}
In close connection with the points of difference already mentioned
between the two pastors stands their different
conception of the relation between self-experience and
tradition.

It is a dangerous matter, says Otto Møller, to speak much
of what one has not yet experienced oneself. But most
people live by faith in authority, not by the
full certainty of self-experience. They believe what the Pope or Luther,
% --- p. 196 ---
\setcounter{footnote}{0}%
\opage{196}the Church or Scripture says, that is, at second hand. It
works only in calm weather. What one can know at first
hand is admittedly not much, --- really only the opposition
between sin and grace; but more may prove
to be contained therein. Of decisive importance it is
to begin from the right end, namely from consciousness of sin.
But whether self-experience can lead to faith
is conditioned by revelation's having been proclaimed to
the human being. --- Here Møller's ecclesiastical standpoint collides
with his approach to a purely psychological consideration.
Even consciousness of sin in the sense in which Møller
sets it forth is nevertheless a historically conditioned phenomenon
and presupposes proclamation. It is not a universally human
experience. Faith is really already a presupposition
for consciousness of sin; there must be faith in a supernatural
authority, before which the human being's striving
can show itself as sin. Møller is here himself too trustful of authority
toward the Lutheran doctrine of ``the order of salvation,''
which incidentally was already opposed in the time of the Reformation.
Can a self-experience not begin with
admiration and enthusiasm for an ideal, and will the feeling
of one's own distance from this ideal not become the deeper,
when self-knowledge is awakened by the thought of the ideal?
When one appeals to self-experience one must also
acknowledge that both its beginning and its course
may be highly different in different personalities.
Lutheranism has fixed a single type,
into which all were to be pressed, more or less violently\footnote[1]{See on various religious types in my Religionsfilosofi § 36—44 and 94—96, and the essay on some works in the philosophy of religion from the most recent time (Mindre Skrifter. Anden Række).}.

With his sense for problems Otto Møller
now also acknowledges that here for a pronounced ecclesiastical conception
% --- p. 197 ---
\setcounter{footnote}{0}%
\opage{197}there lies an insoluble riddle. Can there not be the most inward
need and striving present, without the proclamation
that alone can lead to faith being granted to the human being?
He admits that when faith in external authority has been
swept away entirely, there is no remedy any more. --- Indeed, Møller ought
to have been more of a pessimist than he is; for if one has
no faith, or rather some hope, that there can arise
new ideals by a purely human way, and that they can
awaken the self-knowledge without which life gets no
truth, then all channels are closed.

Skat Rørdam is now from the first apprehensive about
his friend's strong emphasis on self-experience, because it seems
to him to lead to building on bare ground. The human being,
he says, begins with faith in authority and remains with it
all his life. And even in the greatest human distress
the human being must still be able to glimpse and sense the value of faith
and the power of authority. The presupposition is that the human being
lives in a Christian country. How it goes when this
presupposition is not present, Rørdam does not speak of.
His more sharply thinking friend sees here an insoluble
riddle, because he so decidedly stresses both the personal
moment, self-experience, and the objective moment,
revelation and proclamation.

Rørdam trusts more in the congregation's influence on
individuals than Møller. In the congregation as it now exists
there lives, according to Rørdam, a testimony which goes back
to the apostolic congregation and to the writings of the Apostles, and
this apostolic congregation is the beginning set by God.
But Rørdam admits that within the
apostolic writings, which come to individuals through
the congregation, as ``the congregation's inheritance,'' it is not possible to distinguish
once for all between what is of God and what
is human; every Christian must draw the boundary
according to his
% --- p. 198 ---
\setcounter{footnote}{0}%
\opage{198}Christian consciousness nurtured through life in the congregation. And yet, as Rørdam maintains in opposition
to ``modern theology,'' the Church is ``the pillar and
ground of the truth.''

Otto Møller seems not to harbor the confidence in the congregation's
influence that Rørdam does. He thinks that ``when
the congregation is not directed to it, it does not take offense
at very much, but prefers to let everyone sweep before
his own door.'' And he is against a church constitution, because
under such a one it would at bottom be pastors and theologians
who came to rule. --- If one asks what then becomes of
``the congregation,'' I suppose that Møller would
refer to the sacraments. He is not afraid to
give thought a certain latitude, provided only the essential
in worship, the sacraments, is protected. If only this is observed,
``various capers'' may be allowed in other respects.
--- A solution naturally does not lie in this, for
the sacraments themselves presuppose, both as regards their tradition
and their interpretation, a dogmatic foundation
--- and where does that come from? ---

Only the Catholic Church can here bring the matter to a real conclusion.
For it there is here no riddle. Neither is there any insoluble
riddle, though certainly a problem, for those who stand
outside all confessional religion. For them the criterion of truth consists
in the most comprehensive and most coherent
experience [Erfaring] that it is possible to attain, and there
arise here, as regards the relation between individual and society,
ever new questions, which must find their answer
through historical and psychological investigation
of how the individual develops under the influence of society
to greater or lesser independence, and
how the individual on his side can react back upon
society and become co-determining for its further
development. Nowhere do we here stand before anything absolutely
first or last, though it may be difficult to
% --- p. 199 ---
\setcounter{footnote}{0}%
\opage{199}follow the ways by which the social influence asserts itself
in individuals, and to trace the first origin of individual
initiatives. An ``insoluble riddle'' there is
only when one proceeds from the necessity of an absolute authority
and yet will not be a Catholic.

\romanhead{III}
Hitherto I have dwelt predominantly on questions in
which the two friends stood against each other. There are of course
many points, quite apart from their common Grundtvigian
stage, in which they stand together. If their disagreements
have an interest for the psychology of religion by showing how
differently a view of life can be found in people
at the same stage, then their agreements have an interest
by showing how earnest pastors in our
day look upon the spiritual movements that may have
consequences for their Church, and upon the very historical foundation
of their Church.

They agree in their view of Neo-Grundtvigianism,
in which they see partly an effect of people's having
come too easily to their religious stage, partly an expression
of an immature need to bring the ecclesiastical conception into
connection with tendencies in the literature and politics of the age.
They both feel aversion toward the Inner Mission and display
an unequivocal contempt toward ``the third tendency''
and the ruling theology. They also look alike upon the
more recent investigations into the origin of the books of the Bible,
the literary history of the Bible, or, as it is
usually called, because in much and many things it rejects the traditional opinions about
these books, Biblical criticism. They are
both very apprehensive about this whole movement, but
it must be said to their credit that they see clearly that here
are questions which cannot be put aside. They both hold
that if one first concedes the grounds that lead to
% --- p. 200 ---
\setcounter{footnote}{0}%
\opage{200}placing ``the books of Moses'' in a time far later than
that of Moses, then it will be only a question of time when
the same method is applied to the books of the
New Testament. If Judaism is not of divine
origin, Christianity, which builds on it,
cannot be so either. As Møller puts it:
if Wellhausen gets in among the people, Christianity gets out!
It is a work of undermining that is going on. But
gradually they both perceive that here is something one cannot
get away from. Møller expresses their common standpoint in
his forceful manner: ``The painful thing is now that I am convinced
that there is something to this Old Testament
criticism; but it has fallen into some damned
hands.'' They both feel that they lack the knowledge
to take up the matter, and hope that the right scribe
will come who can help the congregation over the difficulties.
--- Whether it will help that the investigations
come into other hands, when the method is the same,
time will show.

From ``the Biblical critics'' to ``the freethinkers'' there is in the two
friends' eyes no long step, and yet their
judgments here fall far more sharply. As the ecclesiastical natures
they both are, each in his own way, living themselves in on all sides
into a way of thinking that never wholly frees itself from authority,
however much the one of them stresses self-experience
and reflection, they can see in free intellectual work
nothing but wickedness, when its results conflict
with their faith. Rørdam sets up the principle that ``what
we as God's children have to ask about is not the cause,
but the aim.'' He applies this principle to the
unhappy war of 1864 and twenty years later to the burning of the Palace.
Both misfortunes were to be understood as God's threatening speech
to the people. (What the cause of the Palace fire was he knows
moreover very well: that the Palace --- as he himself
% --- p. 201 ---
\setcounter{footnote}{0}%
\opage{201}saw with his own eyes --- burned so quickly, made it evident
enough that it was God's will that it should burn!). It
is clear that when one has access to knowing God's
purpose with what happens, or with the coming into being of certain writings,
the interest in finding the cause of these
things becomes very slight, or rather, the causes must be
given with the purposes, lie in the same will that
sets the purposes. From such a standpoint it is
hard to arrive at an understanding of the love of truth
that can lead to the abandonment of traditional assumptions
in the religious domain. For Rørdam freethinkers are
``deniers of God and conscience,'' who regard
the congregation as their real enemy, and all Christianity
that is not the purest individualism as hypocrisy.
When such people work for civil freedom,
it is only in order to use it to suppress Christianity,
which they hate because it condemns them.
Their ultimate aim is to make man a thrall of
nature, to enjoy instead of obey, to make man
like the dumb creatures. (One sees that Rørdam knows
not only about divine but also about human
purposes). Otto Møller makes some reservations
toward this forceful characterization; but he too
holds that what is really the matter with the freethinkers
is that they love darkness more than light, and that the fault lies
in the will, not in thought. At bottom there are no
excusable freethinkers.

It would surely be easier, however, to get along
with the pastor at Gylling than with the Bishop of Zealand. He
has a sense for the significance of self-experience and reflection,
and he has the need to go ``behind what has been handed down.''
At bottom he would be obliged to admit that
it is not the right to free inquiry and independent conduct of life
that is to be grounded, but rather the right of authorities
% --- p. 202 ---
\setcounter{footnote}{0}%
\opage{202}to set a limit to this right. This principle for
the distribution of the burden of proof in the spiritual domain is based
on this, that only in free self-surrender and on the basis
of personal experience can truth become the human being's
real possession.

\romanhead{IV}

I come now to a question that is of quite
special interest for the relation between modern Christianity
and primitive Christianity, namely how the representatives of the Church
take up the expectation of Christ's
imminent return, and with it of the close of the whole course of the world,
which is one of the fundamental thoughts of the New Testament,
without which in particular one does not understand at all the ethics of the New Testament. For an old Kierkegaardian there
is special reason to draw attention to
this point. For when Kierkegaard ended by declaring
that the Christianity of the New Testament does not exist, it is
the ethics of the New Testament he has in view. In dogmatic
respect he thought that nothing was the matter.
Yet Kierkegaard did not see that the ethics of the New Testament
is determined and conditioned by that great expectation which
is the real content of the faith of primitive Christianity. He
was even inclined to reverse the relation and to regard
the expectation as a psychologically necessary
illusion, without which the strict primitive Christian ethics could not
be maintained. This is of course incorrect. It is
the expectation of the imminent ``Kingdom of God'' that lies at the
foundation, and of which the ethics, which under these conditions
could only become an ``interim ethics,'' as a more recent investigator
has called it, is a necessary consequence. Let one try for once
to read the New Testament through with the constant
thought of the immediate, literal
% --- p. 203 ---
\setcounter{footnote}{0}%
\opage{203}nearness of ``the Kingdom of God,'' and everything will become quite otherwise intelligible than it
otherwise is. But the whole Church, the Catholic as well as
the Protestant, here goes the ways of rationalism, and
this from a natural instinct of self-preservation, since the Church's own
continued existence becomes inconceivable and unjustified
on that presupposition of an early close. And the whole
significance one attributes to Christianity for human
culture and for the development of the life of peoples comes to
lack a foundation. Rationalism begins early,
already in the Second Epistle of Peter, with the proposition that for
God a thousand years are as one day, and conversely. But if
God had precisely promised something for a quite definite
point of time? The God of rationalism has eternity
on his side, and differences of time dwindle to nothing; but the God of the Bible
acts in time and brings his Kingdom to completion
at definite times. One cannot appeal to Jesus
and the Apostles on all other points and reject their
statements on a point of such decisive significance as
this. For two wholly different views of life
here come to stand over against each other. For the
one, all human cultural work must be put in abeyance
before the one thing needful, which is to prepare for
the imminent close of all. The other lives and breathes
in the work for the human unfolding of life in social conditions,
in science and in art; here work is done
for the long term, and a close or interruption stands
only as a remote possibility. Between these two views of life
one must, or would have to, choose. The difference
between them became clear when the living expectation
paled or in any case lost its immediate
practical significance. Then one had time for many things,
and could have one's heart in many places, which within
primitive Christianity there could be no thought or talk of.
% --- p. 204 ---
\setcounter{footnote}{0}%
\opage{204}It is this difference that makes modern Christianity into a
different religion from primitive Christianity.\footnote[1]{I have previously expressed myself on this question, first in the concluding section \textit{Kierkegaard som Filosof} (1892), later in my \textit{Religionsfilosofi} (1901). (§ 14. 59. 99. 123).}

If one comes with this problem in view to the correspondence
between the two pastors, one finds indeed some
hints that they have glimpsed it, but they have not
put it in a decisive way and therefore not
answered it either.

Rørdam is in embarrassment over the interval between
death and judgment. This embarrassment does not exist
for the New Testament, since the interval of time is presupposed to
be quite short, and for many of the dying generation
not to exist at all. There is, says Rørdam quite rightly
(and Møller agrees with him in this), nothing said in
the Bible about what happens in that interval, in particular
nothing about a possible development of what was not
completed in and with the individuals before their death. --- But,
it must be added, there was indeed also no reason to say
anything about it, when no longer interval of time
was presupposed to exist.

Sometimes there stirs in the two pastors, under
the impression of the many ungodlinesses of the age, a presumption
that the end is near. They agree that the signs
of the times point strongly to the end, and that the world has grown
old. But as good Grundtvigians they are then disturbed
about what is to become of the Master's prophecies
of the mighty development of the life of the people. There seems indeed to
be growth; but there is still a long way ahead, and at bottom
it looks black everywhere now (1873). Hence it must
seem that it is a long way to the end. Otto Møller declares
also that when he agrees that the signs of the times
point toward the end, it is only the head that concludes
% --- p. 205 ---
\setcounter{footnote}{0}%
\opage{205}thus; the heart has no faith in the Lord's speedy coming,
and it is also for his part the season of growth
that Grundtvig inaugurated that he hopes must lead
to something despite all obstacles.

To the life of the people, for whose future the two pastors hope
in their hearts, so that it is only with their heads
that they can hold fast to the great, all-conditioning hope
of primitive Christianity, they undoubtedly reckon the life of the Church as well. They do not
speak, however, of the ecclesiastical development too being
broken off, if the Return occurred now. Not only culture
and the life of the people, the Church too has arranged itself for the long
term. There would be consternation among bishops, pastors
and theologians, if the end came now. Thor Lange has in
a little story (``Kommet igen.'' ``Ude og Hjemme''
1882) shown how Jesus' return would cause disturbance
in the existing Church. The Church has once for all
stepped into Jesus' place and is not at all arranged for his
taking the helm himself directly. Lange has set the scene
in the Middle Ages, but the depiction fits just as well
the Protestant Church, which has even arranged
itself quite as cozily as the Catholic.

Some 30 years later the two friends again come upon
the problem without noticing its sting. They wonder
that no special ethics is set forth in the New Testament.
Everything turns on awakening and nourishing faith in Jesus;
then it is presupposed that with this the manner and
direction of the conduct of life is also settled. Rørdam in particular occupies himself much
with this. ``There is here a good deal to think about, which I
am far from having come to the end of, but which I toss about
with on my morning walks.''

The question finds its simple answer when one
holds fast that faith in the New Testament is faith in Jesus
as the one who will soon bring the consummation. That ``the Kingdom of
Heaven'' is at hand is to be taken in the literal sense, despite
% --- p. 206 ---
\setcounter{footnote}{0}%
\opage{206}the ecclesiastical reinterpretations. If so, life is borne
by a great expectation, in comparison with which everything else
pales. Inward solidarity among those who live in
this expectation, and striving to spread it
are natural consequences. But concerning social order, art,
science one cannot expect to find anything in the
New Testament, and one also finds nothing. The ethics of the New Testament,
which Kierkegaard could not find in
modern Christianity, and which Bishop Rørdam could not
find in the New Testament itself, is the ethics of the expectation,
which naturally could subsist only as long as
the expectation was actual. The problem of the episcopal morning walks
finds here its simple, historical solution.

In his own way the pastor at Gylling comes upon the
same side of the problem. He has brooded over the parable
of the wise and the foolish virgins, which deals
with those who wait for Jesus' return, and he has been halted
by the feature that all the virgins, both those who had
oil in their lamps and those who had none, fell asleep.
This he expounds in the following way: ``It is quite
certain that no one can endure to live in uninterrupted
expectation of the Lord's coming; it leads only to unhealthy
and unnatural overstrain.'' Otto Møller starts here
from the supposition that it is to be a long time that is waited in. Even
a short waiting-time, when it is something great and absolutely
decisive that is awaited, can seem long, since ``hard
times pass slowly.'' But the time until the return is presupposed
precisely in itself to be short. In the short interval,
of which there is talk in the New Testament, the expectation might
very well hold out, and it did hold out in
the first generation. Otherwise it must naturally go,
when thousands of years pass, unsurveyable stretches of time,
beyond which perhaps the head, but not the heart can
see. But besides --- one thing is to sleep, another is what
% --- p. 207 ---
\setcounter{footnote}{0}%
\opage{207}the Church and human beings have done in that interval,
which indeed continues, namely to put force and interest
into wakeful and living cultural work and human development
and to have both one's heart and one's head occupied
by it. Sleep is a rest that the first Christians needed
like all human beings; they could very well from it,
or from other rest, wake up to the expectation
in which their whole spiritual life had gathered.
More difficult is it to wake up from the cultural work
in which Christians and non-Christians often take part with equal
zeal. The duality, the opposition, and finally the contradiction
between original Christianity and human culture
could not be illuminated more clearly than by Otto Møller's exposition
of the virgins' sleep. This duality is modern
Christianity's real problem, however little consciousness
of it there may seem to be among the Church's representatives.

The Catholic Church holds fast in its way to the ethics of expectation
through the conduct of life of its monks and nuns.
Therefore it sounds both from Saint Teresa in
the 16th century and from Cardinal Newman in the 19th
century: ``What would the world become, if there were
no monks!'' And according to Pascal Jesus' agony lasts
right to the end of the world, and one dare not sleep all that time!\footnote[1]{Jesus sera en agonie jusqu'à la fin du monde; il ne faut pas dormir pendant ce temps-là! \textit{Pensées}, éd Brunschwicg. Paris 1904. II. p. 435.}.
Only the living expectation's involuntarily transforming influence is here replaced by conscious asceticism. ---

All ideal life is a life in expectation, and it has
certainly (as I have already remarked in my Religionsfilosofi)
been of great importance that the religion which got
such great influence on Europe's spiritual life and the life of its peoples
was from the first a religion of expectation. But in order
to become the foundation for free human conduct of life
% --- p. 208 ---
\setcounter{footnote}{0}%
\opage{208}the expectation must be tied to tasks on which one
can work within the world that is accessible to
human experience. It is a great, and often difficult
shift of motive that must take place in the transition from
Christian to human ethics.

\begin{center}* * *\end{center}

The reading of the two pastors' correspondence is of no
small psychological and historical significance in the
light it casts on modern churchmen's attitude, partly
to their own doctrine, partly to connected
circumstances. --- I do not doubt that the most modern churchmen
present other character traits. Time will show whether
they will come to stand otherwise toward the problems
than their predecessors.

It should still be mentioned that the two friends in their old
age experienced something which, by the explanation of one of them,
was connected with their life's work. --- In Rørdam's last
letter it says: ``It has not gone with me as
I had expected in my younger days, that as I
grew older my capacity to rejoice in God and
to look upward and forward with a living hope would grow\dots{}
But of this I experience so little, not to say nothing.'' And Otto Møller answers: ``I am no richer in it
than you are\dots{} I feel so poor, tired and powerless,
as surely no one can, and moreover it even seems to me
that I am dull and do not feel it as humbling as I
ought.''

It is perhaps an experience which many old people
must have, and which finds its natural explanation
in declining vitality. But Otto Møller hints at an explanation
that gets special significance by coming from a
man who laid great weight on self-experience and serious
% --- p. 209 ---
\setcounter{footnote}{0}%
\opage{209}appropriation: ``Have we not perhaps talked and written ourselves tired?''
--- When it belongs to one's profession to speak again and again
of things that have value only as objects of
living feeling, the time will easily come when one notices
that the inner life cannot keep pace with the forms
that have been laid out for it, and disappointments of the mentioned
kind will lie near. A Greek author has rightly
said that from human beings we learn to speak, from the gods to
be silent. A little holy silence might perhaps also
do our most modern churchmen good.

% --- p. 210 ---
\setcounter{footnote}{0}%
\opage{210}\essayhead{MADVIG AS A PHILOSOPHER}

\essaysource{J. N. Madvig: Livserindringer (1887). \\ Supplementer til Livserindringer af J. N. Madvig. Udgivne af J. N. A. Madvig (1917).}

It was of great interest to me to read the great philologist's
autobiography, and no less the later published ``Supplementer,''
which contain, among other things, his philosophical
ideas and characterize his position on the religious
problem. I never belonged to his closer circle,
but in my younger days I repeatedly came
into contact with him and had conversations with him about
philosophical and religious questions. The hints he
gave on such occasions have been confirmed and
made clearer by the notes now available.

My first acquaintance with Madvig came when, like
so many schoolboys, I studied his Latin grammar.
It was naturally not always pleasant reading, but
I remember from the later school years how, with
great profit and with what I would now call logical and
systematic interest, I gained an overview of his
presentation of the syntax of the Latin language. It was, along
with mathematics, the most instructive in formal respect
of anything I came into contact with at school. And when I
sometimes, tired of going from paragraph to paragraph, turned
to the Preface and read his remark that ``this
grammar is written with a mind that allows the
% --- p. 211 ---
\setcounter{footnote}{0}%
\opage{211}present life its full due, while it turns
toward the past as that whose rich inheritance we are to preserve,''
--- then I caught a glimpse of a larger context, into
which my school studies could fit, and to which
they were to give me access. It is touching to me now in my
old age to read that it was only reluctantly that
he had undertaken to write a school grammar,
``to attach his name to the school bench''; he no doubt knew
how most schoolboys regard the author
of the grammars they must be tormented with. He overcame
this reluctance in recognition of the significance
the work could have. But it was connected with something
that even once led to a crisis in his spiritual life.
The philological study, which he cultivated so fruitfully
and so brilliantly, nevertheless stood before him as ``the greatest and most deeply
felt one-sidedness.'' Even in his older days he could
not ``bear to die leaving behind an unrectified idea
of himself as an essentially exclusive philologist,'' --- he wrote in
the year 1869 (65 years old). But already as a young professor
he felt (most acutely in the years 1830—32) in a
painful way this one-sidedness. His poor circumstances in
childhood and youth had not left him time for a free
and general orientation in the domains of the life of the mind,
or, as he expresses it, ``for working upon the personality
while surveying the various interests of life
and aspirations of life in relation to the particular calling and
the particular task.'' There arose in him a strong feeling
of the ``inadequacy of his general-scientific and philosophical
education,'' and ``only slowly and gradually
did he succeed in uniting the eager special study
with a quiet occupation with more general questions.''

There is hardly any doubt that this crisis was closely connected
partly with his conception of classical 
% --- p. 212 ---
\setcounter{footnote}{0}%
\opage{212}antiquity, partly with the direction in which, according to his gifts, he could above all cultivate
classical philology. He could not regard classical antiquity as the great
model, ``as a rule-giving, unattainable world of models,''
as the preceding age had conceived it, and as many of
his fellow practitioners have later regarded it. In contrast
to this he set a more sober conception. The culture of
Greco-Roman antiquity stood for him as primitive in relation to ours; but
precisely thereby it stood as simple and clear and as a
significant foundation for later European development,
a foundation that ought to be known at first hand
by those here at home who were to have a more thorough
education. This foundation required, in his eyes, supplementation,
and he was one of the first here at home who demanded
room in the school for a more extensive natural-scientific
and mathematical instruction than earlier times
have considered necessary. --- But nevertheless it was not, in
his work in the field of classical philology, particularly
the cultural-historical side that he laid weight on.
His epoch-making work went in the direction of critical
examination of manuscripts and of text-critical studies.
This work of procuring genuine and faultless texts,
a work that is so necessary for a fruitful use
of the old authors, was the chief work of his life,
and his fame rests on it. Only to a
relatively small extent has he contributed cultural-historical
investigations, although it was essentially on this side
of the classical studies that he grounded the necessity
of a considerable place for them in the school. Men like Boekh
and Wilamowitz, who quite especially occupied themselves with
the literary and cultural history of antiquity, have hardly to so
high a degree as Madvig felt the one-sidedness that so long
tormented him. On account of the direction of their studies they have
% --- p. 213 ---
\setcounter{footnote}{0}%
\opage{213}been better able to find points of attachment in the culture of more recent times, and thereby more directly to see the general significance
of their special activity.

But apart from this, something stirred in him
which, owing to the division of labor in the intellectual
domain, must assert itself in every researcher, ---
a feeling that the subject one cultivates brings one into
relation only with a single side of existence, --- a need
for a more many-sided content of thought than one's own subject can
yield. Naturally such a need stirs more strongly the
more specialized the subject is. But it can stir in all,
in the philosopher as well as in the grammarian or the entomologist.
In Madvig's case this need led
him first and foremost to philosophical studies. The Hegelian
philosophy's speculative procedure did not satisfy
him, although he himself at certain points had
occasion to take an interest in Hegel's points of view,
especially those of the philosophy of law and the philosophy of language. But he undertook
a coherent study of Spinoza and Kant, and in
his later years, when moreover his weakened sight made
philological work more difficult for him, he occupied
himself thoroughly with philosophical reading, or rather with
hearing philosophical writings read to him. And he
regarded it as a happiness that he could thereby ``conclude
and complete the spiritual content of his life,'' and
that he had gained ``peace to give himself up to contemplation
of the deeper questions of life and existence.'' Also
in other scientific domains he sought to orient
himself. He heard lectures on chemistry with Zeise
and let his friend the mathematician Jürgensen introduce him
to the foundations of higher mathematics. These studies
stood him in good stead in his activity as inspector of education
at the higher schools and gave him a many-sided
view of school questions. I once heard one of
% --- p. 214 ---
\setcounter{footnote}{0}%
\opage{214}his mathematical colleagues say that ``only a colossus
like Madvig'' could have a judgment about all sides of the
work of the higher school.

It was, however, not merely a need for a more comprehensive
content of thought that stirred in him. What is
perhaps most interesting in this universal researcher
is the need for a formation of the mind that could permeate
his whole personality --- the need for what he
himself has called ``working upon the personality while
surveying the various interests of life.'' In very
many men of science, not only does the special subject stand
isolated in relation to other subjects, but the whole content of thought stands
remote from one's own personal life, as this is
led outside the study or the laboratory. This
duality Madvig wanted to get beyond. He wanted his
thought-life also to develop his personality. He
strove ``to unite the need of thought for clarity with
the inwardness of feeling and the demand of love.'' And this
striving did indeed set its stamp on him, in the mild
and truth-loving mind that revealed itself in his
conduct. ---

Let us try to gain an impression of the direction in
which his reflection has moved. His notes
are brief and fragmentary, composed in his heavy,
but as a rule clear language. They give his provisional drafts.
But a definite direction of thought can be traced clearly\footnote[1]{In the following summary I have as far as possible retained Madvig's own expressions.}.

Madvig's philosophizing led him early to dissociate himself
from the speculative philosophy's attempt to derive
everything from pure thought. The task of philosophy was for him
% --- p. 215 ---
\setcounter{footnote}{0}%
\opage{215}to shape what is found in the world of observation into coherent
concepts, so that our own consciousness
forms the starting point. Philosophy is clarification of our consciousness's
undeniable content and pursuit thereof in all
its consequences. And we never find consciousness
without a definite content. The difference between our I
and its content is not original. We find from the very
beginning consciousness as double-sided, at once subjective
and objective, and the task is then to go back
to the most general form by the way of analysis.

We can sense or apprehend nothing without its
appearing in difference or opposition to something else
within the given, and it is precisely owing to such
differences in the given content that we can become
conscious of ourselves as comprehending these differences in unity.
The ruling subject becomes conscious of itself as a unity
over against a mobile and changing limitation. But
within the changing content Nothing can be apprehended as
really having happened except in connection with something Other that
has previously happened, and in continuous connection with this.
Such a continuity between every single fact
and other facts is expressed in the concept of cause, which has
its origin in the subject's basic consciousness, whereas each
single cause and effect must be found by the way of experience.
The first germ of the law of causality shows itself in involuntary
movements, e.g. such as ward off an
intervention; thereby it is presupposed both that a change
can occur and that it can be prevented. All that
thus asserts itself as a link in a continuous
connection we attribute existence, the ``to be.'' Thus
everything we can know stands for us as part of
a whole and becomes reality for us only thereby. And
all becoming in the particular gets, through the connection in
everything, influence on the whole.

% --- p. 216 ---
\setcounter{footnote}{0}%
\opage{216}What we call possible or contingent is only that of which we
have incomplete knowledge of all the conditions.
Error proves itself practically by a contradiction between
the action that certain ideas lead us to undertake,
and its result, in that the expectation we cherished is not
fulfilled. In error, but especially in doubt, we have
the strongest proof of the independence of the life of the mind. The acknowledgment
of something can be in conflict with reality,
or we can even wholly deny or suspend our acknowledgment.

The law of causality holds also in the domain of the will. In
every state in which one acts, the personality
itself is a definite factor, which never becomes zero over against the
influences coming from outside. The outcome of the decision is
determined by the relation of strength between the personality and
the influences. But perhaps not my whole nature
comes to exert its influence on the outcome; in that case
I can afterwards possibly react against it, which is expressed
by the conscience making itself felt. True
freedom consists in the ability to fight for the inner conviction
about one's goal against obstacles from other quarters. ---

The mental and the bodily differ in that
the former is observing, the latter only observed. There is no
transition from the bodily to the mental, although the latter
is throughout bound up with the former. But the mental comes
not from outside into this connection. The mental and the bodily are
two sides of the same. By way of analogy
we find ourselves in continuous connection, not only with
other organic beings, but also with what goes on
in the chemical processes under the opposition of the substances
and their interaction. In all chemical combinations there must,
together with the combination of the elements, occur a combination
of the elements of mind resting in these elements,
% --- p. 217 ---
\setcounter{footnote}{0}%
\opage{217}and something of the kind occurs especially at the coming-to-be of the soul
in the formation of the organism.

Throughout the various forms of nature there can be traced, especially
when we take Lamarck's hypothesis as a basis, an ascending
series of stages from inorganic matter through
the multiplicity of organic beings to the life of mind and love
that can develop in the human world. And through all matter, bodily form
and life of consciousness there unfolds a great will to existence.
It is no violation of the high value of the mental
as the innermost in man that the whole of nature
is regarded as a striving toward this culmination. ---

To these epistemological and natural-philosophical ideas
is joined a series of thoughts on the foundation of the
ethical.

Nor in this domain is pure thought
sole ruler. A morality of reason can have significance only
as clarification and rounding off of the life-drives that lie
in human nature; in itself it would have no content.
There manifests itself in man a willing, a drive
to existence, not only to one's own existence, but also
to the existence of fellow-living individuals. Both in sharing joy
and in compassion it shows itself that the individual feels connected
with other beings, especially those standing nearest.
And when the understanding of the nature and circumstances of others develops,
mere fellow-feeling becomes a total feeling,
which manifests itself as justice. The law for the individual,
which arises from this, is the expression of a divine
will, of God's nature flowing into us; for this
law indicates the definite direction of development for the
whole within which man has come to be. Only
in a developed social state can the sense of duty become firm
and clear. It develops in interaction with the various
forms of life that humanity through
% --- p. 218 ---
\setcounter{footnote}{0}%
\opage{218}immemorial times has groping and tentatively formed and constantly
forms for a harmonious life together. Every moral judgment
is grounded in the individual's relation to race, people and
humanity, thus to the whole to which it belongs. What is
at issue is a harmony between the drive to self-preservation and
the sympathy, both of which lie in man's nature. ---

These epistemological and ethical reflections bear witness
in all their brevity and provisionality that Madvig's
thoughts on the whole moved in the direction of what, at the
end of the 19th century, asserted itself
in the philosophical field. The interest and goodwill
he on so many occasions showed younger philosophizing
researchers becomes thereby understandable. The younger man who in
conversations with him came upon philosophical questions,
got the impression of standing before a man who thought in an
independent way. Publicly he did not express himself on
questions of this kind, and especially not on his position
on religious questions, the one which we will now seek to characterize.

In the autobiography he indicates clearly enough that his religious
opinions more and more diverged, not only from the church's
doctrine, but also from Christianity itself. As a schoolboy
at Frederiksborg he discussed with comrades
the difficulty for faith in revelation that lay in the fact that
the revealed content had come down to us with the
contingency and uncertainty that is bound up with all
historical tradition. Especially emphasized are the critical and
philological doubts concerning the text of the Bible. As a mature man
and researcher he came, after he (in November 1848)
had taken over the Ministry of Worship, to stand before
preparing the fulfillment of the new Constitution's promise of a
% --- p. 219 ---
\setcounter{footnote}{0}%
\opage{219}constitution for the Danish church. But he felt that he
could not with confidence, firmness and seriousness lead negotiations
in this direction, because he, although in his own way
inwardly religiously disposed, did not stand on the standpoint of
the whole faith in revelation. It was his duty to consolidate and build up;
tear down and weaken he did not wish to; --- and he
feared to waver or to play the hypocrite. These scruples
increased with the years. He was elected a member of the
Church Commission of 1868, but thinks he would have
done most rightly not to accept this election, because in
the intervening years he had removed himself farther and more definitely
from the firm ecclesiastical foundation. The consequence was
that within the Commission he acted more conservatively
than earlier, as he thought he ought to uphold the right
of what exists.

What held him back he has clearly stated in
the posthumous notes. He raises thus the
question how a spiritual companionship between man
and wife is possible in relation to religion, where the man
is a philosopher and man of science. And when he felt
the unhappy position between the recognized truth and
the fear that to express it to others could
work destructively without possibility of replacement and rebuilding,
he thought here especially of ``near and dear
beings, in whose feeling one ought not to intervene, because
one does not know whether one can lead them consolingly and satisfyingly
along new paths.'' What we here meet in Madvig
is connected with the mild and soft in his character.
And he has no doubt also felt that for a thinker the grounds are
more important than the results, and that the results by
those who do not understand the grounds are readily conceived wholly
externally and dogmatically and are associated with the ideas
that through centuries have become attached to
heresies. He behaved searchingly and could only be
% --- p. 220 ---
\setcounter{footnote}{0}%
\opage{220}understood by seekers --- and how many are they? It cannot be denied, however, that
it would have been of great significance
if a man with his scientific and social
authority had freely and openly stated his views.
Toward those nearest him it must nevertheless have been
possible for him to procure respect and
recognition for his personal conviction. There manifested itself here in him an anxiousness,
which also marked his attitude during the great
political struggle in the '80s.

In conversation he concealed neither his religious standpoint nor
what held him back from expressing himself. When I in
the year 1874 had given a lecture in the Students' Association
on Darwinism and Philosophy, which was printed in ``Nær
og Fjern,'' he expressed in a cordial way his sympathy
for my line of thought, adding that the thought
of a Fall seemed to him absurd. Then he said
with a sly smile: ``But I dare not say that to my
daughters!'' Some years later I turned to him to
propose that he sign an appeal to contribute
to a statue of Spinoza, which with funds from
all countries was to be erected at The Hague, and I reminded him
that he had often spoken of Spinoza with great sympathy
(a sympathy that now, through the posthumous notes,
is fully understandable). He answered then that he had
a scruple on account of the prejudice that was still harbored by
many against Spinoza, and that there were considerations
he was obliged to take. In order not to be misunderstood
he explained that it was consideration for his home he
was thinking of. Later he nevertheless decided to sign
and got considerable contributions collected.

His great considerateness toward those nearest him has
also been one of the reasons that he constantly went to church,
which had the consequence that he was counted one of the Church's
pillars. Yet the need not to separate himself from
% --- p. 221 ---
\setcounter{footnote}{0}%
\opage{221}those nearest him was not the only reason for his churchgoing.
He had a need for fellowship\footnote[1]{When he (according to the funeral oration) also went to communion with his family, this too was out of a need for spiritual fellowship. He could have said what Rousseau said when he went to communion in Neufchatel, though for him it was a mere ceremony: Je suis avec mes fréres!}, and he was
strongly captivated and felt strengthened by the organ-playing
and the church singing, --- just not by the sermon, as
one sees from the following note.

``In church. Organ-playing and singing, well led and well executed,
lift me into a mood in which there is at once
a feeling of the life of humanity, longing, love and
courage and a participation in the human life together of all times and places
in spirit and love. But then comes the sermon
with its so often halting and self-contradictory dogmatics,
its obscurity and crooked conclusions. Only the purely
ethical exhortation appeals.''

Also in a statement, as it seems, written down later
he sharply emphasizes the opposition between the
mood-filled liturgy and the dogmatic sermon.

After reading these statements one gets a strange
impression of the speech that was held at his funeral,
and which the editor of the ``Supplementer'' has printed
at the back of the book. He is praised here as the Church's faithful son.
``Sunday after Sunday he sat here in the church in his accustomed
place, while the Gospel was proclaimed.'' The preacher
has lived in happy ignorance of what his most eminent
hearer thought of his ``proclamation.''

Now it is indeed so that the great, often mighty poetry,
which is formed in the ecclesiastical worship, is not the Church's
property, but belongs to us all. It is the work of humanity,
as all poetry is. Everyone has here, as in other
spiritual domains, the right to take what he
can use in his spiritual household. And it is likewise
% --- p. 222 ---
\setcounter{footnote}{0}%
\opage{222}a natural need that one wants to share with those nearest what gives life loftiness and
value for them. But misinterpretation lies near when one tries
to take the kernel without making it evident that one is cracking
the shell; the consequence is that one confirms others in thinking that
shell and kernel cannot be separated. In any case one must
hope that respect for the individual personality and
its right to take its own standpoint in the questions of life
will be acknowledged in earnest, not only in larger circles,
but also in the most inward relations between people.
Then an accommodation that conflicts with what personal
truth requires will become impossible in the eyes of all parties. The
more woman develops to be able to take an independent
stand toward the questions of life, the more will also
Madvig's anxious question, how a philosopher
and man of science can have spiritual companionship with
his wife, find its answer --- of itself. ---

Let us now see what thoughts were fostered about the religious
problem in the great philologist's secret chamber.

The idea of God is due to a need for conclusion, for
a last ground or cause, a dissatisfaction with the multiplicity of existence
and its constant changes. But thought realizes that the concepts ground and cause only
find justified application within the existence given in experience.
God can then not be sought outside the world,
but must be conceived as the very unity and totality of existence,
--- as that which stirs in everything, in all thinking and willing,
as in all movement. Only through its forms of manifestation and
its activities do we know the Deity. To call it
personal is unjustified, since the concept of personality presupposes
a relation of interaction with something external and therefore
% --- p. 223 ---
\setcounter{footnote}{0}%
\opage{223}fits only finite beings. The right concept of God
is found only by seeing the inner connection among all
laws of nature and forces of nature and by heeding the direction of development
that makes itself known in existence. He who
is a link in the progressing development feels himself a participant
in a great, universal striving, and this feeling is the religious
feeling. Through it the sap of the whole tree of life streams
into every bud and every blossom. But when thought
tries in a graphic manner to express what
lies at the ground of this feeling, it has only images
from the finite world at its disposal. Thereby are understood
the opposite impressions that the ecclesiastical liturgy and the dogmatic
sermon make.

It is possible that most people cannot do without
mystical and dogmatic forms of religion. But
it is not certain that the Christian faith in its dogmatic
forms has to a very high degree penetrated into the real consciousness
of congregations and individuals; people lived
and live rather in a general, more or less pure
feeling of self-surrender to God and confidence, and only at certain times
was fanaticism created for uncomprehended
formulas and names.

This was a quiet man's quiet thoughts, sprung
from deep longing for truth, but held back by tenderness
toward those who thought differently, who stood near him. It must
not be forgotten that the distinguished man's real strength
found application in faithful work in other domains,
and it is no wonder that he did not wish to let these
fragments of thought come forth, he who otherwise with critical
seriousness considered every line he wrote. He had
the view that very few think things through
on this point. Also in the thinkers
% --- p. 224 ---
\setcounter{footnote}{0}%
\opage{224}and poets of the 19th century --- Kant and Hegel, Schiller and Goethe ---
he found unclarity in their relation to the Church. What
has been --- he asked --- the real conviction
and feeling among the thinking people who at all
in the last century have dared to fix their gaze on
this point? We go about disguised among one another.

By this last sentence (a quotation) I assume he
means that when the great common symbolism that the Church
has developed is given up as insufficient, there is no
common language in which serious people, who
deeply feel the riddles of life, can communicate with one another.
In the negations one can meet; but they are, after all,
not enough for a language. ---

For my part I think that we owe Madvig's
son thanks, because he has resolved to publish
these ``Supplementer,'' which are a new testimony that
more brooding and thinking goes on in the secret chambers than is often
believed. Behind the outer forms of the spiritual life there stirs
in the best an inner life, which perhaps some day will be
able to procure for itself forms of expression and forms of life. Yet
there is hardly any possibility that the opposition between
the outer and the inner life will ever disappear,
and it was not to be wished either.

% --- p. 225 ---
\setcounter{footnote}{0}%

\opage{225}\essayhead{ARISTOTLE ON FRIENDSHIP}

\romanhead{I}

The most discussed and at the same time most frequently misunderstood
thought in Aristotle's Ethics is his determination of
virtue, i.e. the sum of valuable traits of character,
as the right mean between opposites. One
has found a superficial middle-way theory therein and has not
noticed that Aristotle by mean does not understand
a mechanical average, so that, if e.g. 6 is
too much and 2 too little, 4 should be the right. Since
human individuals come with highly different need and
capacity, the mean, that is, properly speaking, the harmony of
the mental life's different elements, will get a different character
for each individual. What is rashness in
the one, can be cowardice in the other. It depends
on what inner conditions the individual meets
with. Just as the director of bodily exercises does not prescribe
the same diet for all participants, but takes account
of their bodily constitution, so one must
in fixing valuable traits of character take
account of the mental constitution. Ultimately it holds
for each one to find his own right mean.

This thought has seldom found recognition in
philosophical ethics, perhaps because it makes ethical questions
% --- p. 226 ---
\setcounter{footnote}{0}%
\opage{226}much more difficult than most care to think
them\footnote[1]{Cf. \emph{Holberg}: \textit{Moralske Tanker} (1744) p. 539: ``One who has no gall cannot be praised for composure. To ride a feeble ass is no merit, but to govern a wild and spirited horse, that deserves applause and clapping of hands. This, however, people commonly do not take into account.''}.

This regard for individual differences stands in
peculiar opposition to the pronounced social character
of Aristotle's Ethics. It is, after all, not said that because
the individual succeeds in harmonizing his drives and
capacities, he is therefore equipped to satisfy the demands
that society and the state for their part make on him. It is
perhaps for this reason that Aristotle in the closer carrying-out
of his doctrine of the right mean involuntarily leaves the
individualistic standpoint and applies a socially determined
average.

But also from another side the question can be raised
whether the individual can really be a closed,
harmonious little world for himself. There could be elements
in his very nature which pointed beyond this inner
world. In Plato, who also will assert a doctrine of harmony,
though not with the far-reaching recognition of individual
differences as in Aristotle, there stirs
decidedly such an element, namely Eros, the need
to attach oneself to other beings or endeavors and
thereby merge into a more comprehensive whole --- a love relation,
a life in social activity or in a world of thought.
Where Plato, as in the ``Symposium,'' lays such a
need at the ground, he presents a quite different type
than where the standpoint of harmony rules, as in the third and
fourth books of the ``Republic.'' In a more subdued, but equally principled
form an analogous element appears in Aristotle.
In the eighth and ninth books of his Ethics he speaks about
% --- p. 227 ---
\setcounter{footnote}{0}%
\opage{227}friendship (filia), treating, on the whole, of a sympathetic relation to
other people. It is a somewhat different ethical type that
appears here than in the books that deal with the ``mean.''
In the need for friendship, the sympathetic need, there is something
that points beyond the individual and can hinder harmony,
in that it ties man to other people, their
weal and woe, their fates. Is it now certain that the
inner harmony can be preserved when the individual is thus
drawn out into a larger circle of interests and fates?

In the first book of his Ethics, which deals with happiness,
Aristotle has noticed the problem and expressed it in a
naive way, in that he, after having established that happiness
presupposes a certain completeness and self-sufficiency
(autarky), adds that this self-sufficiency does not mean
isolation; it holds not for the individual in a
solitary life, but also for his parents and children and
wife, furthermore for friends and fellow citizens, since
man by nature is a social being. But --- now
he becomes uneasy. A limit must be set here, for otherwise
it goes on to infinity, since we should have to count ancestors
and descendants and friends' friends in! --- Aristotle
stands here right at the problem. The question is, by what
principle the limit is to be set --- starting from the individual himself
or from a larger circle and its conditions of life. Unrest
enters the closed concept of happiness. --- As
we shall see, Aristotle returns to the problem in another
connection --- and precisely starting from the nature of
sympathy itself.

I mention here in passing that the problem is sharpened
in Stoicism, which will make the individual independent
of others' wickedness and misfortune, and in Christianity, in that
faith (as condition for one's own salvation) and love
come to stand as opposites. Already among the Indians
the problem stirs. The sage Bharata was just
% --- p. 228 ---
\setcounter{footnote}{0}%
\opage{228}about to reach the highest perfection --- but he
missed it, because he pitied a little sick gazelle!
--- We need not go far to feel the sting of the problem. Who
can feel happy when the greater part of mankind
mutually destroy each other? ---

Even if it is correct that there appears in Aristotle
an opposition between sympathy and individual harmony,
it is yet natural to ask whether there is not
a standpoint from which Aristotle's conception of harmony
among the elements of mental life as the highest trait of character
can be seen in connection with the weight
he lays on the relation of friendship in its different degrees
and kinds. And such a standpoint Aristotle himself gives,
and one that places us before one
of the deepest thoughts in his philosophy and before one of
the most beautiful and finest traits in his personality.

\romanhead{II}

Precisely the happy man needs friends, says Aristotle.
And that not merely because having friends is
the greatest of all external goods, but because the able
and serious man needs someone to whom he can do good deeds.
Man is, as Aristotle again in
this connection declares, a social being, and happiness
consists in activity, not in enjoyment. What a living
Danish poet says: ``I adjure you, friends, to
love your happiness!'' --- agrees entirely with Aristotle's
line of thought. Happiness was for the ancients a serious matter,
a matter of will, one with the assertion of personality against
the external and internal forces that could split life or
make it one-sided. It must not be forgotten that the Greek
harmony was won through struggle with chaotic forces,
just as the Olympians had had to fight the Titans,
and just as the Greeks had had to fight with the
% --- p. 229 ---
\setcounter{footnote}{0}%
\opage{229}Orientals before Greek clarity and harmony could
come into being. Also in the individual great characters this must
have been so. In his characterization of Socrates Professor
Heiberg has shown how much indicates that the
Socratic self-mastery was won through struggle.

And as happiness consists in activity, in the unfolding of the soul's
powers, so friendship too reveals itself
in an activity. Friendship is a virtue,
or it presupposes virtue. It shows itself more in loving
than in being loved. What use is wealth, power
and honor, when one does not use these goods for an activity
fruitful for others? And to this friendship gives the
most beautiful occasion.

In the thought that friendship, like happiness and virtue,
consists in activity, Aristotle finds the explanation of the fact
that benefactors so often love those to whom they have shown a benefaction
more than they are loved by these. The reason for this
could, he thinks, lie in the fact that in him to whom they have done
good, and in his welfare, they see their own work, in which they
have laid down something of themselves. Our work is indeed a
part of ourselves, a realization of our possibilities. We
love our work because we love to act, and this again,
because to be is to act. Therefore, Aristotle adds,
artists love their works more than they would be
loved by these, if they became living. And poets especially!
They have an altogether excessive love for their productions.
--- Especially when the work has cost trouble, it will
be an object of love. Therefore mothers are apt
to love their children more than fathers do. Birth has
cost pains, and mothers have a more emphatic
knowledge that the children are theirs. So much, says Aristotle
in another connection, can mothers love their children
that they can hand them over to others, if it is best for
the children. ---

% --- p. 230 ---
\setcounter{footnote}{0}%
\opage{230}Already this whole line of thought shows that there is a
connection between love of others and love
of oneself. One does for one's friend the good
one can, and that for the friend's own sake. But quite
so stands the able and serious person toward
himself, especially toward the best in himself, the life of thought.
He wants to live in and with himself, in his memories and in
his expectations, and he feels at home in this
inner world. And quite in the same way he lives in his
friends' thoughts, their joys and sorrows. The friend is an
other I. One can conceive the relation to oneself as
an analogy to the relation to one's friend, or conversely friendship
as an extension of love of oneself.

When one blames love of oneself, it is because
one thinks only of the self-assertion that manifests itself at
the expense of others. But there is a love of oneself
that consists in care for the most beautiful goods, for
the development of our spirit and our character. Such a
love of oneself can even lie at the ground of the sacrifice
of life and property. One will then rather perform one great
deed than many small, and in that one for one's country or
for one's friends wins what is necessary for them,
one wins the beautiful for oneself. (Necessary Aristotle calls
what has value only as means, beautiful what
has value in itself). Yes, a good man will
even be able to renounce coming to perform
a beautiful exploit himself, in order that his friend may come
to perform it. He can find it more beautiful that his friend
gets the honor than that he himself gets it.

A closer ground for this line of thought
Aristotle finds in this, that we can better form a graphic
image of others than of ourselves. We wish to see ability
and energy, because life and happiness consist in vigorous unfolding,
--- and that we can see in our friends in common
% --- p. 231 ---
\setcounter{footnote}{0}%
\opage{231}life and cooperation with them for human tasks.
It is namely not enough to graze in the same field, as
cattle do, side by side. Cooperation in theoretical
pursuit is probably what Aristotle especially thinks of,
but he sees that human cooperation can have many
forms, and whatever common interest then operates, it can
ground a relation of friendship. One can hold a feast
together with one's friends, or play dice with them,
or exercise gymnastically with them, or go hunting with
them, just as well as one can philosophize with them.

He who is not a friend with himself can, according to Aristotle,
not be a friend with others in the true sense either.
There must be harmony in a person, in order that he
may be able to come into harmonious relation to other people.
And harmony here means in particular agreement
between what man is and does, and what he
wishes and wills. Without such an agreement
life becomes split and restless, and a firm relation to others
impossible. Only the ``able and serious'' person can
enter into a true relation of friendship. In contrast to such
a character Aristotle sets the ``bad,'' who cannot
be a friend with himself, because his life is split;
he wills one thing and does another; from cowardice and sluggishness
he never comes to carry out what he wills. Therefore
he finds no pleasure in himself and does not like to dwell in
his own inner being, his memories and his expectations.
If he seeks company with others, it is in order to forget himself.
Real friendship he cannot enter into.

It is perhaps characteristic of Aristotle as a Greek,
that in this connection he knows only two types,
the harmonious and the disharmonious, and has an inclination
to let the opposition between them coincide
with the opposition between goodness and badness. He
proceeds from the assumption that the bad man is always disharmonious, and turns
% --- p. 232 ---
\setcounter{footnote}{0}%
\opage{232}it then without more ado around, so that the disharmonious
person also must be bad. But there are disharmonious,
depressed, wavering and irresolute natures,
who are indeed unhappy, but not ``bad,'' and who may very
well enter into a relation of friendship, when the other
party has the right understanding. Although such a person
finds nothing lovable in himself, has no
sympathy with himself, cannot be a friend with himself,
he can nevertheless, precisely on account of the contrast, understand,
admire and love a composed and harmonious nature,
especially when this meets him with understanding.
Shakespeare has depicted such a relation of friendship.
Hamlet loves Horatio, this firm, in spite of all suffering
harmonious nature, and Horatio understood his friend and had
discovered the noble heart in his split inner being. Therefore
Hamlet also in his hour of death charges him to explain
him and his cause to those who do not understand. An example
from real life is Schleiermacher's relation to
Friedrich Schlegel (in the latter's young, pre-Catholic time). In a
letter Schleiermacher states that everything that seems erroneous,
self-contradictory and wrong in Schlegel is connected
with his circumstances, and that as the innermost in
his poetry and striving something great, rare
and beautiful constantly stirs. --- The presupposition for a friendship between
a harmonious and a disharmonious nature will precisely
be that in the midst of the cleavage and contradictions
in the latter there is discovered a striving, a working
will, even if this is not able to reach up to the
harmony that rules in the inner and outer
conduct of more fortunate natures.

Had Aristotle held fast to his doctrine of the dependence of harmony,
of the ``right mean,'' on individual presuppositions,
he could have extended his doctrine of friendship
with an interesting section. It is not always the
% --- p. 233 ---
\setcounter{footnote}{0}%
\opage{233}meanest and poorest natures that lack harmony. It
may be precisely the fullness, the many possibilities, the rich
powers, that hinder unity and closedness. Perhaps
this kind of nature is more frequent in modern times than in antiquity.

Aristotle has, moreover, himself in another section of his
Ethics depicted a disharmonious nature, whom he in any
case does not call wholly ``bad.'' It is people who lack
self-control. Either they follow every impulse,
without reflecting, or they do not hold fast the
result of reflections at the moment of action. And yet they
constantly acknowledge the right; the goal of their striving is good,
but cannot break through in the individual moments.
--- Here was in any case a starting point for an understanding
and a relation of friendship present, namely for him
who could catch sight of the again and again repressed
goal.

Moreover Aristotle himself comes to hint at a peculiar
reason why unhappy natures so often are not suited
for friendship. Manly natures, he says, do not want
their friends to share their misfortune and grieve for their
sake; it is enough for them that they themselves are unhappy. Thus
also, we may add, disharmonious natures will shut
themselves up within themselves, unable or unwilling to open themselves
to others, and that precisely out of sympathy with those with whom they otherwise
could have come into the most inward relation.
In Søren Kierkegaard's Papers one can study this
type, which presupposes self-control outwardly, though
this is lacking toward the inner chaos.

\romanhead{III}

The noblest kind of friendship is according to Aristotle friendship of character,
and in what precedes we have kept
to this type. It has the firmest ground, in contrast
to the friendship that is due to utility or pleasure.
% --- p. 234 ---
\setcounter{footnote}{0}%
\opage{234}Utility and pleasure can change, and the friendship will fall away
when the one no longer has any advantage from the other
(just as the friendship of states ceases with the mutual
advantage), or when charm, and with it infatuation, has
vanished. Yet Aristotle holds that the friendship that
is grounded on infatuation can more readily pass over into friendship of character
than that which is grounded on utility.

A condition for perfect friendship is longer
living together and cooperation, whereby the characters can stand
their test. One does not know one another before one has
eaten a certain amount of salt\footnote[1]{Aristotle alludes to a well-known saying, which the Greeks already used. Salt was an expression of guest-friendship. In a later Greek author the expression occurs ``to have eaten a bushel of salt together.''} with one another. All that
can bring about living together and cooperation can give occasion to a
relation of friendship. Thus kinship, traveling company, relation of fellow citizenship
etc. In the particular case friendship will get
its own character according to the different social and political
relations. ---

Relations of friendship between unequals require a special
investigation. A very great inequality in nature and
circumstances, with excessive superiority on the one side, will
make friendship impossible. Thus between gods and men.
Although one therefore wishes one's friends all good, one will
therefore not wish them to become gods. As regards men,
the powerful man will easily obtain relations of friendship
for the sake of utility and pleasure. But an able and serious
man will not readily become a friend with one standing higher in power,
unless this one possesses quite special
traits of character, so that the admiration and love
these awaken can counterbalance the inequality in the external
circumstances.

Other relations of inequality --- between older and younger,
parents and children, man and wife, master and slave ---
% --- p. 235 ---
\setcounter{footnote}{0}%
\opage{235}nuance the relation of friendship in manifold ways, in
analogy with the relation between the powerful and the lowly.
The relation between man and wife is thus compared
with an aristocracy, where power is exercised for the benefit of
the lesser party. The man rules by virtue of the dignity of his nature,
but he grants the wife the influence that belongs to
her; if he rules merely for his own advantage, marriage
becomes comparable to an oligarchy. As regards the relation
to slaves, Aristotle teaches that one can
be a friend with one's slave, insofar as he is a human being,
but not insofar as he is a slave. In his doctrine of the state he defends
expressly slavery as an institution grounded in nature. ---

Not least in this whole section one detects Aristotle's
personal experiences. In the Platonic school, to which he
belonged for twenty years, he was able to experience how
friendship can be built up through common interests and labors,
and how living together and cooperation can lead to
mutual understanding of the characters. Also the relation
between older and younger he was here able to see the significance
of, and when he himself became head of a school, he has
learned to know this relation from the other side toward
his disciples. The relation to the more powerful he lived through
in his position toward Alexander, whose teacher he had
been, and Hermias, the ruler of Atarneus, who took
him under his protection. From his will it is seen that
he was a good husband and a good master toward
his slaves.

Just as living together can call forth friendship, so duties
and rights within the relation of friendship follow
from the nature of the relation of living together. They are not the same within
a traveling company, within a family life, within relations of fellow citizenship,
within guest-friendship etc. Justice
and friendship both come into being within relations of living together,
% --- p. 236 ---
\setcounter{footnote}{0}%
\opage{236}and justice grows in refinement with the relation of friendship,
both as regards degree and form. A higher
right holds between friends than between those who stand
farther from one another. And at its culmination justice takes on
the stamp of love.

This last utterance is probably to be understood so,
that the perfect understanding of individual character
and individual circumstances, which the relation of friendship partly
presupposes, partly conditions, makes possible a consideration of
the individual personality's character and circumstances, which the justice
formulated in general rules is not able to
exercise.

Just as Aristotle here passes from sympathy to justice,
so he conversely comes in the fifth book
of his Ethics, which deals especially with justice, to
the thought that the general laws can never exhaust all
the individual circumstances and conditions of an action, but
formal justice is supplemented or corrected by
understanding of this individual side of the action, and
the quality by virtue of which such a supplementing
or correcting becomes possible, he there calls equity
(Epieiki).

\romanhead{IV}

When perfect friendship shows itself in cooperation and
in sympathy, it must be limited in extent. We have already
seen how Aristotle from a quite different side,
namely from the individual's happiness, comes to the necessity
of a limitation of human sympathies.
But the ground for this in the books on friendship is of
special interest. Aristotle has a clear eye for the fact that
in the domain of the life of feeling a law of inverse
relation between intensity and extent holds. Therefore the number of
friends must be limited. The individual cannot split
% --- p. 237 ---
\setcounter{footnote}{0}%
\opage{237}himself, and at its culmination friendship can concern only a
single person. To this it is added that our friends should preferably
be friends with one another, and this becomes more difficult the more of them
there are. This consideration Aristotle applies also to
the number of citizens in a state, although there it is only
a question of political friendship, which has a greater extent than
personal friendship. There too there is a limit,
and 10 myriads (100,000) Aristotle thinks are too many
for a state. He lays in his doctrine of the state great weight
on all citizens knowing one another. Furthermore
he applies this consideration to a critique of Plato's
communistic ideal of the state, according to which family life
and property are abolished, for the ruling class,
so that all who belong to this class can regard
one another as brothers. The sympathy that can
rule among men when all family bonds fall
away will, thinks Aristotle, become very thin, like when
one mixes a little bit of wine into a great quantity of water! ---

We see here the Greek thought of the necessity of limitation
from its most beautiful side. And we understand better
Plato's and Aristotle's zeal for asserting the opposition
between Hellenes and barbarians. It is through life in the
narrowest and most inward human relations that human beings
can develop to live with warmth and strength in the
most comprehensive relations. If the general ideal of humanity
is to lead to more than slack sentimentality,
it must work its way forward step by step. The ancient thinkers
were right that the way goes through, not around,
the nationalities.

Yet the decisive step beyond the national
limitation was already taken in classical antiquity. In guest-friendship
and in humaneness toward strangers, which
already in Homer stands under Zeus's special protection,
there lay from ancient time a germ that could later be
% --- p. 238 ---
\setcounter{footnote}{0}%
\opage{238}unfolded. Aristotle recognizes the possibility of friendship between
the free and slaves, and since he grounds slavery especially
on the fact that the barbarians are slaves by nature, this
recognition leads beyond national limitation on an essential
point. There developed toward the end of antiquity,
under the influence of the social conditions that the conquests of Alexander
and the Romans had called forth, a universally human
fellow-feeling. It seems to have been the
Stoic school through which this development came to
consciousness, whereby it became possible that the Greeks' Latin
apprentice Cicero could speak of caritas generis humani.

\end{document}
